% L’élève produira un rapport — Français Tle Tech — Cours signé OnBuch+
% Programme officiel MINESEC. Compiler avec : tectonic N05-eleve-produira-rapport.tex
\newcommand{\DOCMATIERE}{Français}
\newcommand{\DOCNIVEAU}{Tle Tech}
\newcommand{\DOCMODULE}{Français – classe de Terminale technique}
\newcommand{\DOCLECON}{N05}
\newcommand{\DOCTITRE}{L’élève produira un rapport}
% OnBuch+ — préambule « cours signés » ENRICHI (compilé avec tectonic / XeLaTeX)
% Système visuel « Pop » d'OnBuch+ : fond crème (jamais blanc), bordures encre
% épaisses, ombres « dures » décalées, titres Archivo Black en capitales.
% Version MATHÉMATIQUES : graphes (pgfplots/tikz), boîtes pédagogiques
% (définition, propriété, méthode, attention, le savais-tu, exercice résolu,
% exercice, TP, bilan), page de garde et signature OnBuch+ ancrée en bas.
%
% Macros attendues dans main.tex AVANT \input{preamble} :
%   \DOCMATIERE  \DOCNIVEAU  \DOCMODULE  \DOCLECON (numéro)  \DOCTITRE
\documentclass[11pt]{article}

\usepackage{fontspec}
\usepackage[french]{babel}
\usepackage[a4paper,margin=2cm,top=3.4cm,bottom=2.8cm,headheight=1.4cm,footskip=1.3cm]{geometry}
\usepackage{amsmath,amssymb,amsfonts}
\usepackage{mathtools} % \xmapsto, \coloneqq... souvent utilisés par les modèles
\usepackage{cancel} % simplifications barrées (bilans de forces)
% \abs{z} (module, valeur absolue) et \norm{u} (norme), utilisés par les modèles
\providecommand{\abs}{}\let\abs\relax\DeclarePairedDelimiter{\abs}{\lvert}{\rvert}
\providecommand{\norm}{}\let\norm\relax\DeclarePairedDelimiter{\norm}{\lVert}{\rVert}
\usepackage[table]{xcolor} % [table] : fournit aussi \rowcolors (alternance de lignes)
\usepackage{pagecolor}
\usepackage{tikz}
\usetikzlibrary{arrows.meta,positioning,calc,shapes.geometric,shapes.misc,shapes.arrows,decorations.pathmorphing,decorations.pathreplacing,decorations.markings,patterns,shadows,mindmap,trees,fit,backgrounds,matrix,intersections,angles,quotes}
\usepackage{pgfplots}
\usepackage[version=4]{mhchem} % équations nucléaires \ce{^{235}_{92}U}
\usepackage[european,siunitx]{circuitikz} % schémas électriques
\pgfplotsset{compat=1.18}
\usepgfplotslibrary{fillbetween,groupplots} % aires entre courbes (calcul intégral)
\usepackage{enumitem}
\usepackage{fancyhdr}
% [most] charge toutes les bibliothèques tcolorbox (listings, minted, raster,
% documentation...) et gonfle inutilement la mémoire TeX sur un document long
% (~300 boîtes) : on ne charge que ce qui est réellement utilisé.
\usepackage[skins,breakable]{tcolorbox}
\usepackage{graphicx}
\usepackage{colortbl}
\usepackage{booktabs}
\usepackage{tabularx}
\usepackage{multirow}
\usepackage{diagbox} % case d'en-tête diagonale des tableaux à double entrée
% Colonnes extensibles centrée / gauche / droite (C, L, R), souvent utilisées par les modèles
\newcolumntype{C}{>{\centering\arraybackslash}X}
\newcolumntype{L}{>{\raggedright\arraybackslash}X}
\newcolumntype{R}{>{\raggedleft\arraybackslash}X}
\usepackage{array}
\usepackage{multicol}
\usepackage{siunitx}
% parse-numbers=false : accepte aussi \SI{2,5 \times 10^{-3}}{...}
\sisetup{locale=FR,output-decimal-marker={,},group-minimum-digits=4,parse-numbers=false}
\DeclareSIUnit\ohm{\ensuremath{\symup{\Omega}}} % U+2126 absent de la police
\DeclareSIUnit\division{div} % réglages d'oscilloscope (ms/div, V/div)
\DeclareSIUnit\tr{tr} % tours (vitesse de rotation, ex. \SI{1500}{\tr\per\minute})
\DeclareSIUnit\molar{M} % concentration molaire (ex. \SI{3}{\molar}, \SI{1}{\micro\molar})
\DeclareSIUnit\an{an} % année (le modèle confond parfois avec \year, un compteur TeX) — ex. \SI{5}{\kilo\meter\per\an}
\DeclareSIUnit\FCFA{FCFA} % franc CFA (ex. \SI{500}{\FCFA\per\kilo\gram})
\DeclareSIUnit\degreeBrix{\textdegree Brix} % teneur en sucre d'un sirop/jus (réfractométrie)
\DeclareSIUnit\month{mois} % le modèle utilise parfois \month, un compteur TeX, comme unité de durée
\DeclareSIUnit\microliter{\micro\liter} % le modèle écrit parfois \microliter en un seul token
\DeclareSIUnit\million{millions} % dénombrement (ex. \SI{15}{\million\per\mL} spermatozoïdes)
\usepackage{qrcode}
% \degree : commande fréquemment utilisée par les modèles pour un angle ou une
% température en dehors de \SI{}{} (non fournie par les paquets chargés).
\providecommand{\degree}{^{\circ}}
% \qed (fin de démonstration), utilisé par les modèles sans amsthm
\providecommand{\qed}{\hfill$\square$}
% \barre : double barre des valeurs interdites dans un tableau de variations (tkz-tab)
\providecommand{\barre}{\ensuremath{\|}}
% environnement proof (amsthm non chargé) utilisé par les modèles
\ifdefined\proof\else\newenvironment{proof}[1][Démonstration]{\par\noindent\textit{#1.}\ }{\qed\par}\fi
\usepackage{lastpage}
\usepackage{hyperref}
\hypersetup{hidelinks}

% ============================================================
% Palette « Pop » OnBuch+ — mêmes hex que la classe P (plus_ui.dart)
% ============================================================
\definecolor{poporange}{HTML}{F59321}
\definecolor{popdark}{HTML}{E07A0C}
\definecolor{popcream}{HTML}{FFF6E8}
\definecolor{popink}{HTML}{1C1714}
\definecolor{poppurple}{HTML}{7A5AE0}
\definecolor{popgreen}{HTML}{2E9E6A}
\definecolor{poppink}{HTML}{E0457B}
\definecolor{popblue}{HTML}{2F7DE1}
\definecolor{popgold}{HTML}{C98A12}
\definecolor{popmuted}{HTML}{8A7F76}
\definecolor{popline}{HTML}{EADFD2}
\definecolor{popgray}{HTML}{8A7F76}
\colorlet{popgrey}{popgray}
% Alias tolérants (le modèle nomme parfois les couleurs différemment)
\colorlet{popwhite}{white}
\colorlet{popblack}{popink}
\colorlet{popred}{poppink}
\colorlet{popyellow}{popgold}
% Teintes claires (fonds de tableaux/encadrés) — le modèle nomme parfois
% les couleurs avec un suffixe L (ex. popblueL) sans les avoir définies.
\colorlet{poporangeL}{poporange!20}
\colorlet{popdarkL}{popdark!20}
\colorlet{poppurpleL}{poppurple!20}
\colorlet{popgreenL}{popgreen!20}
\colorlet{poppinkL}{poppink!20}
\colorlet{popblueL}{popblue!20}
\colorlet{popgoldL}{popgold!20}
\colorlet{popredL}{popred!20}
\colorlet{popgrayL}{popgray!20}
% Teintes claires pour fonds de tableaux / zones
\colorlet{popblueL}{popblue!10!white}
\colorlet{poporangeL}{poporange!14!white}
\colorlet{popgreenL}{popgreen!12!white}
\colorlet{poppurpleL}{poppurple!12!white}
\colorlet{poppinkL}{poppink!12!white}
\colorlet{popgoldL}{popgold!14!white}

\pagecolor{popcream}

% ============================================================
% Typographie
% ============================================================
\defaultfontfeatures{Ligatures=TeX}
\setmainfont{PlusJakartaSans-Medium.ttf}[Path=../../../fonts/,BoldFont=PlusJakartaSans-Bold.ttf]
% Maths en sans-serif (Fira Math) avec chiffres et lettres droites de Plus
% Jakarta, pour que les formules se fondent dans le texte.
\usepackage[math-style=ISO,bold-style=ISO]{unicode-math}
\setmathfont{FiraMath-Regular.otf}[Path=../../../fonts/]
\setmathfont{PlusJakartaSans-Medium.ttf}[Path=../../../fonts/,range={up/num,up/{Latin,latin},\mathup}]
\usepackage{icomma} % virgule décimale sans espace en mode math
% Caractères mathématiques Unicode absents des polices de texte (Plus Jakarta,
% Archivo Black) : rendus par leur équivalent mathématique, en texte comme en
% formule — sinon ils s'affichent en carré vide (ex. « Arithmétique dans ℤ »).
\usepackage{newunicodechar}
\newunicodechar{ℝ}{\ensuremath{\mathbb{R}}}
\newunicodechar{ℤ}{\ensuremath{\mathbb{Z}}}
\newunicodechar{ℂ}{\ensuremath{\mathbb{C}}}
\newunicodechar{ℕ}{\ensuremath{\mathbb{N}}}
\newunicodechar{ℚ}{\ensuremath{\mathbb{Q}}}
\newunicodechar{𝔻}{\ensuremath{\mathbb{D}}}
\newunicodechar{α}{\ensuremath{\alpha}}
\newunicodechar{β}{\ensuremath{\beta}}
\newunicodechar{γ}{\ensuremath{\gamma}}
\newunicodechar{δ}{\ensuremath{\delta}}
\newunicodechar{ε}{\ensuremath{\varepsilon}}
\newunicodechar{θ}{\ensuremath{\theta}}
\newunicodechar{λ}{\ensuremath{\lambda}}
\newunicodechar{μ}{\ensuremath{\mu}}
\newunicodechar{σ}{\ensuremath{\sigma}}
\newunicodechar{τ}{\ensuremath{\tau}}
\newunicodechar{φ}{\ensuremath{\varphi}}
\newunicodechar{ψ}{\ensuremath{\psi}}
\newunicodechar{ω}{\ensuremath{\omega}}
\newunicodechar{Σ}{\ensuremath{\Sigma}}
% majuscules produites par \MakeUppercase dans les titres (α-aminés -> Α)
\newunicodechar{Α}{\ensuremath{\alpha}}
\newunicodechar{Β}{\ensuremath{\beta}}
\newunicodechar{Γ}{\ensuremath{\Gamma}}
\newunicodechar{Δ}{\ensuremath{\Delta}}
\newunicodechar{Ε}{\ensuremath{\varepsilon}}
\newunicodechar{Θ}{\ensuremath{\Theta}}
\newunicodechar{Λ}{\ensuremath{\Lambda}}
\newunicodechar{Μ}{\ensuremath{\mu}}
\newunicodechar{Τ}{\ensuremath{\tau}}
\newunicodechar{Φ}{\ensuremath{\Phi}}
\newunicodechar{Ψ}{\ensuremath{\Psi}}
\newunicodechar{Ω}{\ensuremath{\Omega}}
\newunicodechar{↦}{\ensuremath{\mapsto}}
\newunicodechar{∈}{\ensuremath{\in}}
\newunicodechar{∉}{\ensuremath{\notin}}
\newunicodechar{∧}{\ensuremath{\wedge}}
\newunicodechar{⇔}{\ensuremath{\Leftrightarrow}}
\newunicodechar{⟺}{\ensuremath{\Longleftrightarrow}}
\newunicodechar{⇒}{\ensuremath{\Rightarrow}}
\newunicodechar{⟹}{\ensuremath{\Longrightarrow}}
\newunicodechar{∘}{\ensuremath{\circ}}
\newunicodechar{∪}{\ensuremath{\cup}}
\newunicodechar{∩}{\ensuremath{\cap}}
\newunicodechar{≡}{\ensuremath{\equiv}}
\newunicodechar{⊂}{\ensuremath{\subset}}
\newunicodechar{⊥}{\ensuremath{\perp}}
\newunicodechar{∃}{\ensuremath{\exists}}
\newunicodechar{∀}{\ensuremath{\forall}}
\newunicodechar{∛}{\ensuremath{\sqrt[3]{\,}}}
\newunicodechar{∜}{\ensuremath{\sqrt[4]{\,}}}
\newunicodechar{⃗}{\ensuremath{{}^{\to}}}
\newunicodechar{ⁿ}{\ifmmode^{n}\else\textsuperscript{\ensuremath{n}}\fi}
\newunicodechar{ˣ}{\ifmmode^{x}\else\textsuperscript{\ensuremath{x}}\fi}
\newunicodechar{ᵇ}{\ifmmode^{b}\else\textsuperscript{\ensuremath{b}}\fi}
\newunicodechar{ʳ}{\ifmmode^{r}\else\textsuperscript{\ensuremath{r}}\fi}
\newunicodechar{ᵘ}{\ifmmode^{u}\else\textsuperscript{\ensuremath{u}}\fi}
\newunicodechar{ʸ}{\ifmmode^{y}\else\textsuperscript{\ensuremath{y}}\fi}
\newunicodechar{ᵃ}{\ifmmode^{a}\else\textsuperscript{\ensuremath{a}}\fi}
\newunicodechar{ᵉ}{\ifmmode^{e}\else\textsuperscript{\ensuremath{e}}\fi}
\newunicodechar{ᵏ}{\ifmmode^{k}\else\textsuperscript{\ensuremath{k}}\fi}
\newunicodechar{ᵖ}{\ifmmode^{p}\else\textsuperscript{\ensuremath{p}}\fi}
\newunicodechar{ᴺ}{\ifmmode^{N}\else\textsuperscript{\ensuremath{N}}\fi}
\newunicodechar{ᶜ}{\ifmmode^{c}\else\textsuperscript{\ensuremath{c}}\fi}
\newunicodechar{ʰ}{\ifmmode^{h}\else\textsuperscript{\ensuremath{h}}\fi}
\newunicodechar{ᵠ}{\ifmmode^{q}\else\textsuperscript{\ensuremath{q}}\fi}
\newunicodechar{ₙ}{\ifmmode_{n}\else\textsubscript{\ensuremath{n}}\fi}
\newunicodechar{ₐ}{\ifmmode_{a}\else\textsubscript{\ensuremath{a}}\fi}
\newunicodechar{ᵢ}{\ifmmode_{i}\else\textsubscript{\ensuremath{i}}\fi}
\newunicodechar{ᵣ}{\ifmmode_{r}\else\textsubscript{\ensuremath{r}}\fi}
\newunicodechar{ₖ}{\ifmmode_{k}\else\textsubscript{\ensuremath{k}}\fi}
\newunicodechar{ᵦ}{\ifmmode_{\beta}\else\textsubscript{\ensuremath{\beta}}\fi}
\newunicodechar{ₓ}{\ifmmode_{x}\else\textsubscript{\ensuremath{x}}\fi}
\newfontfamily\popdisplay{ArchivoBlack-Regular.ttf}[Path=../../../fonts/]
\newfontfamily\poplogo{SpaceGrotesk-Bold.ttf}[Path=../../../fonts/]
\newfontfamily\popsemi{PlusJakartaSans-SemiBold.ttf}[Path=../../../fonts/]
\newfontfamily\popxbold{PlusJakartaSans-ExtraBold.ttf}[Path=../../../fonts/]
\newfontfamily\popxboldls{PlusJakartaSans-ExtraBold.ttf}[Path=../../../fonts/,LetterSpace=3.0]

\setlist[enumerate]{leftmargin=1.7em,itemsep=5pt,topsep=4pt}
\setlist[itemize]{leftmargin=1.7em,itemsep=4pt,topsep=4pt,label={\color{poporange}\textbullet}}
\setlength{\parindent}{0pt}
\setlength{\parskip}{5pt}
\renewcommand{\arraystretch}{1.25}

% pgfplots : style Pop par défaut
\pgfplotsset{
  every axis/.append style={axis line style={popink,line width=1pt},tick style={popink},
    grid style={popline,line width=0.5pt},label style={font=\small},tick label style={font=\footnotesize}},
  popcurve/.style={poporange,line width=1.8pt,no markers,smooth},
}
% Même style disponible dans un \draw TikZ simple (hors pgfplots)
\tikzset{popcurve/.style={poporange,line width=1.8pt}}
\tikzset{popgrid/.style={popline,very thin}}
\colorlet{popcurve}{poporange} % le nom du style est parfois utilisé comme couleur

% ============================================================
% Pill / squiggle
% ============================================================
\newcommand{\poppill}[3]{%
  \tikz[baseline=-0.55ex]\node[draw=popink,line width=1.2pt,rounded corners=2.6pt,%
    fill=#2,inner xsep=6.5pt,inner ysep=3pt]{%
    \popxboldls\fontsize{7}{7}\selectfont\color{#3}\MakeUppercase{#1}};%
}
\newcommand{\popsquiggle}[1]{%
  \tikz[baseline=-0.4ex]{
    \draw[#1,line width=2.6pt,line cap=round]
      plot [smooth] coordinates {(0,0.09) (0.34,0.01) (0.68,0.21) (1.02,0.01) (1.36,0.10)};
  }%
}
% Mot-clé surligné (marqueur orange sous le texte)
\newcommand{\cle}[1]{{\popxbold\color{popdark}#1}}

% ============================================================
% En-tête / pied
% ============================================================
\pagestyle{fancy}
\fancyhf{}
\renewcommand{\headrulewidth}{0pt}
\renewcommand{\footrulewidth}{0pt}
\lhead{\raisebox{-0.15ex}{\poplogo\fontsize{13}{13}\selectfont\color{popink}OnBuch\color{poporange}+}}
\rhead{\poppill{\DOCMATIERE\ \textbullet\ \DOCNIVEAU\ \textbullet\ Leçon \DOCLECON}{white}{popink}}
\lfoot{\popxbold\fontsize{7.6}{7.6}\selectfont\color{popmuted}\MakeUppercase{Cours signé OnBuch+}\ \textbullet\ {\popsemi\DOCTITRE}}
\rfoot{\tikz[baseline=-0.6ex]\node[rounded corners=8.5pt,fill=popink,minimum height=17pt,minimum width=17pt,inner xsep=5pt,inner sep=0]{\popxbold\fontsize{8}{8}\selectfont\color{popcream}\thepage\,/\,\pageref*{LastPage}};}

% ============================================================
% Titres
% ============================================================
\newcommand{\chaptitre}[2]{%
  \begin{tcolorbox}[enhanced,colback=poporange,colframe=popink,boxrule=2.6pt,arc=11pt,
      drop shadow=popink,left=15pt,right=15pt,top=12pt,bottom=11pt,before skip=3pt,after skip=18pt]
    \poppill{#2}{popink}{popcream}\par\vspace{9pt}
    {\raggedright\hyphenpenalty=10000\exhyphenpenalty=10000\popdisplay\fontsize{22}{24}\selectfont\color{popcream}\MakeUppercase{#1}\par}
  \end{tcolorbox}
}
% Section numérotée : pastille orange avec le numéro + titre Archivo
\newcounter{popsec}
\newcommand{\coursec}[1]{%
  \stepcounter{popsec}\setcounter{popsub}{0}%
  \par\vspace{16pt}\noindent
  \begin{minipage}{\linewidth}
  \tikz[baseline=-0.7ex]\node[draw=popink,line width=1.6pt,fill=poporange,rounded corners=4pt,minimum size=22pt,inner sep=0]{\popdisplay\fontsize{12}{12}\selectfont\color{popcream}\thepopsec};%
  \hspace{8pt}{\popdisplay\fontsize{15}{17}\selectfont\color{popink}\MakeUppercase{#1}}\par
  \vspace{4pt}\noindent{\color{poporange}\rule{36pt}{3.6pt}}
  \end{minipage}\par\nopagebreak\vspace{8pt}
}
\newcounter{popsub}
\newcommand{\courssub}[1]{%
  \stepcounter{popsub}%
  \par\vspace{9pt}\noindent{\popxbold\fontsize{11.5}{13}\selectfont\color{popink}{\color{poporange}\thepopsec.\thepopsub}\hspace{6pt}#1}\par\nopagebreak\vspace{4pt}
}

% ============================================================
% Boîtes pédagogiques — même recette : fond blanc, bordure épaisse colorée,
% ombre dure encre, badge plein. Titre optionnel [#1].
% ============================================================
\tcbset{popbase/.style={breakable,enhanced,colback=white,boxrule=2.3pt,arc=9pt,
  shadow={3pt}{-3pt}{0pt}{popink},left=14pt,right=14pt,top=11pt,bottom=11pt,before skip=11pt,after skip=15pt,
  fontupper=\popsemi\fontsize{10.6}{14.5}\selectfont\color{popink},
  fontlower=\popsemi\fontsize{10.6}{14.5}\selectfont\color{popink}}}
% Coche dessinée (le glyphe ✓ n'existe pas dans les polices du document)
\newcommand{\popcheck}[1]{\tikz[baseline=-0.5ex]\draw[#1,line width=1.6pt,line cap=round,line join=round](-0.13,0.0)--(-0.04,-0.09)--(0.13,0.1);}
\providecommand{\coche}{\popcheck{popgreen}}
\providecommand{\resultat}[1]{\par\smallskip\noindent\fcolorbox{popgreen}{popgreenL}{\parbox{\dimexpr\linewidth-2\fboxsep-2\fboxrule\relax}{\textbf{Résultat.} #1}}\par\smallskip}
\newcommand{\popboxhead}[3]{\poppill{#1}{#2}{white}\ifx\relax#3\relax\else\hspace{6pt}{\popxbold\fontsize{10.6}{12}\selectfont\color{popink}#3}\fi\par\vspace{7pt}}

\newtcolorbox{aretenir}[1][]{popbase,colframe=popgreen,before upper={\popboxhead{À retenir}{popgreen}{#1}}}
\newtcolorbox{exemplebox}[1][]{popbase,colframe=poporange,before upper={\popboxhead{Exemple}{poporange}{#1}}}
\newtcolorbox{definition}[1][]{popbase,colframe=popblue,before upper={\popboxhead{Définition}{popblue}{#1}}}
\newtcolorbox{propriete}[1][]{popbase,colframe=poppurple,before upper={\popboxhead{Propriété}{poppurple}{#1}}}
\newtcolorbox{methode}[1][]{popbase,colframe=popgold,before upper={\popboxhead{Méthode}{popgold}{#1}}}
\newtcolorbox{attention}[1][]{popbase,colframe=poppink,before upper={\popboxhead{Attention}{poppink}{#1}}}
\newtcolorbox{experience}[1][]{popbase,colframe=popblue,colback=popblueL,before upper={\popboxhead{Expérience}{popblue}{#1}}}
% « Le savais-tu ? » : carte violette pleine (contexte, culture, Cameroun)
\newtcolorbox{savaistu}[1][]{popbase,colback=poppurple,colframe=popink,
  fontupper=\popsemi\fontsize{10.4}{14.2}\selectfont\color{white},
  before upper={\poppill{Le savais-tu\,?}{white}{poppurple}\ifx\relax#1\relax\else\hspace{6pt}{\popxbold\color{white}#1}\fi\par\vspace{7pt}}}
% Exercice résolu : énoncé en haut, \tcblower puis la solution sur fond crème
\newcounter{popexo}\newcounter{popexr}
\newtcolorbox{exoresolu}[1][]{popbase,colframe=poporange,colbacklower=popcream,
  segmentation style={popink,line width=1.2pt,dashed},
  before upper={\stepcounter{popexr}\popboxhead{Exercice résolu \thepopexr}{poporange}{#1}},
  before lower={\poppill{Solution}{popink}{popcream}\par\vspace{6pt}}}
% Exercice (non résolu en place) — la correction est dans la partie Corrigés
\newtcolorbox{exercice}[1][]{popbase,colframe=popink,
  before upper={\stepcounter{popexo}\popboxhead{Exercice \thepopexo}{popink}{#1}}}
% Corrigé
\newtcolorbox{corrige}[1][]{popbase,colframe=popgreen,colback=popgreenL,
  before upper={\popboxhead{Corrigé}{popgreen}{#1}}}
% Objectifs / prérequis / bilan
\newtcolorbox{objectifs}{popbase,colframe=popink,colback=poporangeL,
  before upper={\popboxhead{Objectifs de la leçon}{popink}{}}}
\newtcolorbox{prerequis}{popbase,colframe=popink,colback=popblueL,
  before upper={\popboxhead{Prérequis}{popblue}{}}}
\newtcolorbox{bilan}{popbase,colframe=popink,colback=white,boxrule=2.8pt,
  before upper={\popboxhead{Fiche bilan}{poporange}{L'essentiel à connaître pour l'examen}}}
% Figure encadrée avec légende
\newtcolorbox{popfigure}[1][]{enhanced,breakable=false,colback=white,colframe=popline,boxrule=1.4pt,arc=8pt,
  left=10pt,right=10pt,top=10pt,bottom=8pt,before skip=10pt,after skip=12pt,halign=center,
  lower separated=false,halign lower=center,
  fontlower=\popsemi\fontsize{9}{11.5}\selectfont\color{popmuted},
  #1}
\newcommand{\legende}[1]{\par\vspace{6pt}{\popsemi\fontsize{9}{11.5}\selectfont\color{popmuted}#1}}

% ============================================================
% Signature OnBuch+
% ============================================================
\newcommand{\poplogoinline}[1]{{\poplogo\fontsize{#1}{#1}\selectfont\color{popink}OnBuch\color{poporange}+}}
\newcommand{\signatureonbuch}{%
  \begin{tcolorbox}[enhanced,colback=popink,colframe=popink,boxrule=2.2pt,arc=10pt,
      drop shadow=poporange,left=14pt,right=14pt,top=10pt,bottom=10pt,before skip=0pt,after skip=0pt]
    \tikz[baseline=-0.7ex]\node[circle,fill=popgreen,draw=popcream,line width=1.3pt,minimum size=22pt,inner sep=0]{\popcheck{white}};%
    \hspace{9pt}\begin{minipage}[c]{0.62\linewidth}
      {\popxbold\fontsize{12}{13}\selectfont\color{popcream}Signé\ \textbullet\ {\poplogo OnBuch\color{poporange}+}}\par\vspace{2pt}
      {\raggedright\popsemi\fontsize{8}{10.5}\selectfont\color{popline}\MakeUppercase{Équipe pédagogique OnBuch+}\\\MakeUppercase{Conforme au programme officiel MINESEC}\par}
    \end{minipage}\hfill
    {\poplogo\fontsize{22}{22}\selectfont\color{popcream}OnBuch\color{poporange}+}
  \end{tcolorbox}%
}

% ============================================================
% Page de garde
% #1 titre · #2 matière · #3 niveau · #4 module · #5 n° leçon · #6 durée ·
% #7 description / objectifs (texte ou itemize)
% ============================================================
\newcommand{\pagegarde}[7]{%
  \thispagestyle{empty}
  \newgeometry{a4paper,margin=1.7cm,top=1.6cm,bottom=1.4cm}%
  \begin{tikzpicture}[remember picture,overlay]
    \fill[poporange!18] ([xshift=-3.2cm,yshift=-3.6cm]current page.north east) circle (4.6cm);
    \fill[poppurple!14] ([xshift=2.4cm,yshift=7.6cm]current page.south west) circle (3.8cm);
  \end{tikzpicture}%
  {\poplogo\fontsize{30}{30}\selectfont\color{popink}OnBuch\color{poporange}+}\hfill
  \raisebox{0.6ex}{\poppill{Cours signé \textbullet\ Premium}{popink}{popcream}}\par
  \vspace{1pt}\popsquiggle{poppurple}\par
  \vspace{8pt}
  \begin{tcolorbox}[enhanced,colback=poporange,colframe=popink,boxrule=2.8pt,arc=13pt,
      drop shadow=popink,left=18pt,right=18pt,top=12pt,bottom=13pt,after skip=10pt]
    \poppill{#2 \textbullet\ #3}{popink}{popcream}\hspace{5pt}\poppill{#4}{white}{popink}\par\vspace{8pt}
    {\popdisplay\fontsize{14}{14}\selectfont\color{popink}LEÇON #5}\par\vspace{4pt}
    {\raggedright\hyphenpenalty=10000\exhyphenpenalty=10000\popdisplay\fontsize{25}{27}\selectfont\color{popcream}\MakeUppercase{#1}\par}
    \vspace{8pt}
    {\popxbold\fontsize{9.5}{9.5}\selectfont\color{popink}\MakeUppercase{Durée conseillée : #6}}
  \end{tcolorbox}
  \begin{tcolorbox}[enhanced,colback=white,colframe=popink,boxrule=2.2pt,arc=10pt,
      drop shadow=popink,left=16pt,right=16pt,top=10pt,bottom=10pt,after skip=8pt]
    \poppill{Dans cette leçon}{poppurple}{white}\par\vspace{5pt}
    \popsemi\fontsize{9.3}{12.6}\selectfont\color{popink}#7
  \end{tcolorbox}
  \noindent\begin{minipage}[t]{0.487\linewidth}
    \begin{tcolorbox}[enhanced,colback=popgreen,colframe=popink,boxrule=2.2pt,arc=10pt,
        drop shadow=popink,left=11pt,right=11pt,top=10pt,bottom=10pt,height=3.1cm,valign=center]
      \tcbox[colback=white,colframe=popink,boxrule=1.3pt,arc=5pt,boxsep=0pt,nobeforeafter,
        left=3pt,right=3pt,top=3pt,bottom=3pt,valign=center]{\qrcode[height=1.9cm]{https://play.google.com/store/apps/details?id=com.luvvix.onbuch&hl=fr}}%
      \hspace{8pt}\begin{minipage}[c]{0.5\linewidth}
        {\popdisplay\fontsize{11}{12.5}\selectfont\color{white}\MakeUppercase{Télécharge OnBuch}}\par\vspace{4pt}
        {\raggedright\popsemi\fontsize{8.4}{11}\selectfont\color{white}Gratuit \textbullet\ Google Play\\Tuteur IA, annales, concours, résultats\par}
      \end{minipage}
    \end{tcolorbox}
  \end{minipage}\hfill
  \begin{minipage}[t]{0.487\linewidth}
    \begin{tcolorbox}[enhanced,colback=poppurple,colframe=popink,boxrule=2.2pt,arc=10pt,
        drop shadow=popink,left=11pt,right=11pt,top=10pt,bottom=10pt,height=3.1cm,valign=center]
      \tikz[baseline=-0.7ex]\node[circle,fill=white,draw=popink,line width=1.3pt,minimum size=1.6cm,inner sep=0]{\popdisplay\fontsize{24}{24}\selectfont\color{poppurple}+};%
      \hspace{8pt}\begin{minipage}[c]{0.6\linewidth}
        {\popdisplay\fontsize{11}{12.5}\selectfont\color{white}\MakeUppercase{Passe à OnBuch+}}\par\vspace{4pt}
        {\raggedright\popsemi\fontsize{8.4}{11}\selectfont\color{white}Tous les cours signés, Tuteur IA illimité, fiches PDF — onglet OnBuch+ de l'app\par}
      \end{minipage}
    \end{tcolorbox}
  \end{minipage}
  \vspace{8pt}
  \popcontenu
  \vfill
  \signatureonbuch
  \restoregeometry
  \newpage
}

% Rangée « Ce cours contient » (page de garde)
\newcommand{\poptile}[3]{% #1 couleur #2 gros texte #3 légende
  \begin{tcolorbox}[enhanced,colback=#1,colframe=popink,boxrule=2pt,arc=9pt,shadow={2.5pt}{-2.5pt}{0pt}{popink},
      left=6pt,right=6pt,top=9pt,bottom=9pt,halign=center,valign=center,height=1.75cm,width=\linewidth,nobeforeafter]
    {\popdisplay\fontsize{12.5}{13}\selectfont\color{white}\MakeUppercase{#2}}\par\vspace{3pt}
    {\centering\popsemi\fontsize{7.8}{9.5}\selectfont\color{white}#3\par}
  \end{tcolorbox}}
\newcommand{\popcontenu}{%
  \par\noindent{\popxboldls\fontsize{7.5}{7.5}\selectfont\color{popmuted}CE COURS SIGNÉ CONTIENT}\par\vspace{7pt}
  \noindent\begin{minipage}[t]{0.235\linewidth}\poptile{popblue}{Cours}{complet, illustré, pas à pas}\end{minipage}\hfill
  \begin{minipage}[t]{0.235\linewidth}\poptile{poporange}{Méthodes}{et exercices résolus}\end{minipage}\hfill
  \begin{minipage}[t]{0.235\linewidth}\poptile{poppink}{Exercices}{type examen + corrigés}\end{minipage}\hfill
  \begin{minipage}[t]{0.235\linewidth}\poptile{popgreen}{Fiche bilan}{l'essentiel à retenir}\end{minipage}\par}

% Sommaire
\newtcolorbox{sommaire}{enhanced,colback=white,colframe=popink,boxrule=2.2pt,arc=10pt,
  drop shadow=popink,left=18pt,right=18pt,top=13pt,bottom=13pt,before skip=6pt,after skip=20pt,
  before upper={\poppill{Sommaire}{poppurple}{white}\par\vspace{9pt}\popsemi\fontsize{10.6}{14.5}\selectfont\color{popink}}}

% Signature de fin de cours — poussée tout en bas de la dernière page
\newcommand{\findecours}{%
  \par\vfill\nopagebreak
  \begin{center}\popsquiggle{poporange}\end{center}\vspace{4pt}
  \signatureonbuch
}
\AtBeginDocument{\def\llbracket{\mathopen{[\mkern-3.5mu[}}\def\rrbracket{\mathclose{]\mkern-3.5mu]}}} % intervalles d'entiers (glyphe ⟦ absent de la police)
% Glyphes absents de Fira Math : redéfinis à partir de glyphes disponibles
\AtBeginDocument{%
  \def\setminus{\mathbin{\backslash}}\let\smallsetminus\setminus
  \def\vdots{\mathord{\vbox{\baselineskip3.2pt\lineskiplimit0pt\kern4pt\hbox{.}\hbox{.}\hbox{.}}}}%
  \def\ddots{\mathinner{\mkern1mu\raise6.5pt\hbox{.}\mkern2mu\raise3.6pt\hbox{.}\mkern2mu\raise0.7pt\hbox{.}\mkern1mu}}%
  \def\triangle{\mathord{\scalebox{0.9}{$\symup{\Delta}$}}}%
  \def\nabla{\mathord{\raisebox{\depth}{\scalebox{1}[-1]{$\symup{\Delta}$}}}}%
  \def\checkmark{\popcheck{popdark}}%
  \def\star{\mathbin{\ast}}\def\bigstar{\mathord{\ast}}%
  \def\diamond{\mathbin{\scalebox{0.6}{$\Diamond$}}}%
  \def\bigcup{\mathop{\mathchoice{\raisebox{-0.3em}{\scalebox{1.7}{$\cup$}}}{\scalebox{1.25}{$\cup$}}{\cup}{\cup}}\displaylimits}%
  \def\bigcap{\mathop{\mathchoice{\raisebox{-0.3em}{\scalebox{1.7}{$\cap$}}}{\scalebox{1.25}{$\cap$}}{\cap}{\cap}}\displaylimits}%
  \def\lesseqqgtr{\mathrel{\lessgtr}}\def\bigoplus{\mathop{\oplus}\displaylimits}%
}
\providecommand{\ssi}{si et seulement si} % abréviation fréquente
\AtBeginDocument{\providecommand{\wideparen}{\overparen}} % arc de cercle

%% FIGFIX-BEGIN v5 — correctifs globaux des figures (docs/onbuch-generation/tools/figfix)
\usepackage{adjustbox}\usepackage{etoolbox}\tcbuselibrary{raster}
% toute figure TikZ/pgfplots plus large que la ligne est réduite à la largeur disponible
\BeforeBeginEnvironment{tikzpicture}{\begin{adjustbox}{max width=\linewidth}}
\AfterEndEnvironment{tikzpicture}{\end{adjustbox}}
% légendes pgfplots lisibles par-dessus les courbes
\pgfplotsset{every axis/.append style={legend style={fill=white,fill opacity=0.92,text opacity=1,draw=black!15,inner sep=3pt}}}
% étiquettes d'arêtes (syntaxe edge["..."]) avec fond blanc
\tikzset{every edge quotes/.append style={fill=white,inner sep=1.2pt}}
%% FIGFIX-END
\begin{document}

\pagegarde{L’élève produira un rapport}{Français}{Tle Tech}{Français – classe de Terminale technique}{N05}{15 séances (fiche de progression harmonisée)}{Cette leçon conduit l'élève de Terminale technique à maîtriser le rapport, texte fonctionnel de la sphère professionnelle : repérage de ses caractéristiques informatives, étude de sa structure externe et interne, mobilisation du lexique commun et spécialisé, puis rédaction complète et améliorée d'un rapport authentique à partir d'une situation concrète.
\vspace{4pt}
\begin{multicols}{2}\small
\begin{itemize}
\item À la fin de cette leçon, je sais identifier les caractéristiques d'un texte informatif (objectivité, précision, clarté, concision).
\item À la fin de cette leçon, je sais distinguer la structure externe (en-tête, titre, introduction, corps, conclusion, signature) et la structure interne (situation, constats, analyses, recommandations) d'un rapport.
\item À la fin de cette leçon, je sais différencier le lexique commun du lexique spécialisé et les employer à bon escient dans un rapport.
\item À la fin de cette leçon, je sais rédiger l'introduction et la conclusion d'un rapport (contexte, objet, méthode, annonce du plan ; bilan, recommandations).
\item À la fin de cette leçon, je sais rédiger le corps d'un rapport en paragraphes articulés et hiérarchisés.
\item À la fin de cette leçon, je sais préparer la rédaction d'un rapport : collecte de données, prise de notes, plan détaillé.
\end{itemize}\end{multicols}}

\begin{sommaire}
\begin{multicols}{2}
\begin{enumerate}
\item Découverte du rapport et du texte informatif
\item La structure externe et interne du rapport
\item Lexique commun et lexique spécialisé
\item Rédiger l'introduction et la conclusion
\item Rédiger et améliorer le corps du rapport
\item De la préparation à la rédaction complète du rapport
\item Activité d'intégration
\item Méthodes et exercices résolus
\item Exercices
\item Corrigés des exercices
\item Fiche bilan
\end{enumerate}
\end{multicols}
\end{sommaire}

\begin{objectifs}
\begin{itemize}
\item À la fin de cette leçon, je sais identifier les caractéristiques d'un texte informatif (objectivité, précision, clarté, concision).
\item À la fin de cette leçon, je sais distinguer la structure externe (en-tête, titre, introduction, corps, conclusion, signature) et la structure interne (situation, constats, analyses, recommandations) d'un rapport.
\item À la fin de cette leçon, je sais différencier le lexique commun du lexique spécialisé et les employer à bon escient dans un rapport.
\item À la fin de cette leçon, je sais rédiger l'introduction et la conclusion d'un rapport (contexte, objet, méthode, annonce du plan ; bilan, recommandations).
\item À la fin de cette leçon, je sais rédiger le corps d'un rapport en paragraphes articulés et hiérarchisés.
\item À la fin de cette leçon, je sais préparer la rédaction d'un rapport : collecte de données, prise de notes, plan détaillé.
\item À la fin de cette leçon, je sais rédiger un rapport complet à partir d'une situation concrète de la vie courante ou professionnelle.
\item À la fin de cette leçon, je sais améliorer un rapport (relecture, correction, reformulation) et justifier ma démarche rédactionnelle.
\end{itemize}
\end{objectifs}

\begin{prerequis}
\begin{itemize}
\item Les types de textes (narratif, descriptif, informatif, argumentatif) vus en classe de Première
\item Le compte rendu et la prise de notes (classes de Seconde et Première)
\item Les connecteurs logiques et l'articulation du texte
\item La grammaire de la phrase : expansions du nom, subordination
\item Les textes fonctionnels déjà étudiés (lettre administrative, CV, lettre de motivation)
\end{itemize}
\end{prerequis}

\newpage

\coursec{Découverte du rapport et de la communication}

\courssub{Situation d'accroche}

À la fin de son stage dans une entreprise de transformation de manioc à Bafoussam, Marlyse, élève de Terminale technique, doit rendre compte par écrit de ses observations à son chef de service et à son lycée. Elle a pris des notes, des photos et des chiffres : nombre de bassines traitées par jour, température de séchage, horaires de l'équipe, incidents survenus sur la machine à râper. Mais son premier texte est confus : pas d'introduction claire, des informations mélangées, un vocabulaire imprécis (« la machine marchait un peu mal »), des opinions personnelles (« je trouve que les ouvriers sont sympas »). Son chef de service le lui rend avec cette remarque : « Ce n'est pas un rapport, c'est un journal intime. »

Comment transformer des données brutes en un document professionnel clair, structuré et crédible ? C'est tout l'enjeu du \cle{rapport}, texte fonctionnel incontournable de la vie professionnelle et administrative camerounaise. Toute cette unité d'apprentissage va nous apprendre à le construire pas à pas. Mais avant d'écrire, il faut comprendre ce qu'est un rapport, à quoi il sert, et quelles qualités rendent un texte informatif fiable.

\textbf{Problématique :} qu'est-ce qu'un rapport, quelles sont les caractéristiques du texte informatif qui le compose, et comment le distinguer d'autres textes proches comme le compte rendu ou l'article de presse ?

\courssub{Lecture et réception de textes fonctionnels authentiques}

Avant de définir, observons. Voici trois extraits de textes fonctionnels, de natures voisines. Lis-les attentivement : notre travail de réception consiste à repérer ce qu'ils ont en commun et ce qui les différencie.

\begin{exemplebox}[Trois textes fonctionnels à comparer]
\textbf{Texte A — Rapport de stage (extrait).} « Du 5 janvier au 13 février 2026, j'ai effectué un stage d'observation et de pratique à l'entreprise SOCAMANIOC de Bafoussam, spécialisée dans la transformation du manioc en farine et en bâtons. La production journalière moyenne s'élève à 480 kg de farine. J'ai constaté que la machine à râper n°2 présente une usure des lames qui réduit le rendement d'environ 15 \%. Je recommande son remplacement avant la saison des pluies. »

\textbf{Texte B — Rapport d'inspection sanitaire (extrait).} « Le 18 mars 2026, l'équipe d'inspection de l'hygiène communale a visité le marché Mokolo (Yaoundé II). Sur 120 étals inspectés, 84 étaient conformes aux normes d'hygiène, soit un taux de conformité de 70 \%. 36 étals présentent des manquements : absence de point d'eau (21 cas), denrées exposées au sol (15 cas). Il est proposé une mise en conformité sous quinze jours. »

\textbf{Texte C — Rapport de réunion (extrait).} « La coopérative agricole COOPAGRO de l'Ouest a tenu son assemblée générale le 7 février 2026 à Bafang, en présence de 58 membres sur 72 inscrits. L'ordre du jour comportait trois points : bilan de la campagne 2025, fixation du prix du kilogramme de maïs, élection du bureau. Le bilan fait état d'une production de 312 tonnes, en hausse de 8 \%. Le prix est fixé à 190 FCFA/kg. »
\end{exemplebox}

\textbf{Observation guidée.} Que remarque-t-on ? Les trois textes donnent des \cle{faits} précis (dates, lieux, chiffres), s'adressent à un destinataire identifiable (chef de service, mairie, membres de la coopérative), et évitent les sentiments personnels. Aucun ne dit « j'ai trouvé le marché sale » ou « la réunion était ennuyeuse ». Ils informent pour permettre une \cle{décision} : remplacer une machine, mettre en conformité des étals, fixer un prix. C'est exactement la fonction du rapport.

\courssub{Définition du rapport}

\begin{definition}[Le rapport]
Le \cle{rapport} est un \cle{texte fonctionnel} qui rend compte de façon \textbf{objective} et \textbf{structurée} d'une situation, d'une mission, d'une visite, d'un stage ou d'un événement, à destination d'un \textbf{destinataire identifié} (hiérarchie, institution, administration), afin de l'informer et de l'aider à \textbf{décider}. Il comporte généralement un constat, une analyse et souvent des recommandations.
\end{definition}

Décortiquons cette définition, car chaque mot compte :

\begin{itemize}
\item \textbf{Texte fonctionnel} : il n'est pas écrit pour le plaisir ou pour émouvoir (comme un poème ou un roman), mais pour \emph{servir} à quelque chose : informer, justifier, obtenir une décision.
\item \textbf{Objectif} : il rapporte des faits vérifiables, pas des impressions. « La machine n°2 perd 15 \% de rendement » est un fait ; « la machine est nulle » est une opinion, interdite dans un rapport.
\item \textbf{Structuré} : les informations sont ordonnées (introduction, développement, conclusion), pas mélangées comme dans le premier essai de Marlyse.
\item \textbf{Destinataire identifié} : un rapport n'est jamais écrit « pour tout le monde ». Il s'adresse à quelqu'un de précis qui attend ce document pour agir.
\end{itemize}

\begin{savaistu}[Le rapport, pilier de l'administration camerounaise]
Dans l'administration camerounaise, tout déplacement en mission donne \textbf{obligatoirement} lieu à un rapport de mission transmis à la hiérarchie. Un inspecteur pédagogique qui visite un lycée de Garoua, un ingénieur qui contrôle un chantier à Douala, un médecin qui supervise un centre de santé à Bamenda : tous doivent rédiger un rapport. Savoir rédiger un rapport, c'est donc posséder une compétence directement professionnelle, recherchée dans la fonction publique comme dans les entreprises privées.
\end{savaistu}

\courssub{Le texte informatif et ses quatre qualités}

Le rapport appartient à la grande famille des \cle{textes informatifs} : des textes dont la fonction principale est de transmettre une information. Mais toute information ne se vaut pas : une rumeur informe mal, un bon rapport informe bien. Quatre qualités distinguent le bon texte informatif.

\begin{definition}[Le texte informatif et ses quatre qualités]
Le \cle{texte informatif} est un texte qui transmet des faits de manière neutre et vérifiable. Il obéit à quatre qualités :
\begin{enumerate}
\item l'\textbf{objectivité} : des faits vérifiables, sans opinion personnelle ni sentiment ;
\item l'\textbf{exactitude} : des chiffres, dates, lieux et noms précis ;
\item la \textbf{clarté} : des phrases simples et un vocabulaire adapté au destinataire ;
\item la \textbf{concision} : aller à l'essentiel, sans détails inutiles.
\end{enumerate}
\end{definition}

\textbf{Vérifions ces qualités sur nos textes.} Reprenons le texte B : « Sur 120 étals inspectés, 84 étaient conformes, soit un taux de conformité de 70 \%. »

\begin{itemize}
\item \emph{Objectivité} : le nombre d'étals se compte, le taux se calcule ($84/120 = 0,70$, soit 70 \%). Nul ne peut contester.
\item \emph{Exactitude} : « le 18 mars 2026 », « marché Mokolo », « 120 étals », « 21 cas » : tout est daté, localisé, chiffré.
\item \emph{Clarté} : phrases courtes, un fait par phrase.
\item \emph{Concision} : aucune description pittoresque du marché, aucune anecdote.
\end{itemize}

\begin{attention}[Le piège de l'imprécision]
L'erreur la plus fréquente des élèves est de remplacer les faits par des formules vagues : « beaucoup d'étals », « il y a quelque temps », « une grande quantité ». Dans un rapport, « beaucoup » ne veut rien dire : écrivez « 36 étals sur 120 ». « Récemment » ne veut rien dire : écrivez « le 18 mars 2026 ». Toute affirmation doit pouvoir être vérifiée par le destinataire.
\end{attention}

\courssub{Les marques d'objectivité dans la langue}

L'objectivité n'est pas seulement une intention : elle se voit dans la langue. Quatre marques linguistiques signalent qu'un texte est objectif.

\begin{aretenir}[Les marques d'objectivité]
\begin{enumerate}
\item \textbf{Le présent de vérité générale} : pour les faits permanents. Ex. : « La coopérative compte 72 membres inscrits. »
\item \textbf{Le passé composé et l'imparfait du récit factuel} : pour les actions accomplies et les constats. Ex. : « L'équipe a visité le marché ; 36 étals présentaient des manquements. »
\item \textbf{L'absence de la 1\textsuperscript{re} personne} ou l'usage contrôlé du « je » professionnel : on préfère « l'équipe d'inspection a constaté » à « j'ai vu que c'était sale ». Le « je » n'est admis que pour engager la responsabilité du rédacteur : « Je recommande le remplacement de la machine. »
\item \textbf{Le lexique dénotatif} : les mots sont pris dans leur sens propre, précis, sans images ni exagération. On écrit « usure des lames » et non « la machine est à l'agonie ».
\end{enumerate}
\end{aretenir}

\begin{exemplebox}[Marques d'exactitude dans un rapport de visite du marché Mokolo]
Soulignons dans cet extrait tout ce qui rend l'information exacte : « Le \textbf{18 mars 2026}, l'équipe d'inspection de l'hygiène communale a visité le \textbf{marché Mokolo (Yaoundé II)}. Sur \textbf{120} étals inspectés, \textbf{84} étaient conformes, soit un taux de conformité de \textbf{70 \%}. \textbf{36} étals présentent des manquements : absence de point d'eau (\textbf{21} cas), denrées exposées au sol (\textbf{15} cas). »

On compte ici : une date complète, un lieu précis (marché + arrondissement), cinq données chiffrées et un pourcentage calculé. Vérifions la cohérence interne : $84 + 36 = 120$ étals, et $21 + 15 = 36$ cas de manquements. C'est cette densité de faits vérifiables et cohérents qui rend le rapport crédible.
\end{exemplebox}

\courssub{Rapport, compte rendu, article de presse : ne pas confondre}

Ces trois textes informent tous sur des faits réels. Pourtant, ils ne se ressemblent pas. Le \cle{compte rendu} restitue fidèlement ce qui s'est passé ou ce qui a été dit (une réunion, un match, un cours) ; le \cle{rapport} va plus loin : il analyse et propose ; l'\cle{article de presse} informe le grand public avec un ton parfois plus libre.

\begin{center}
\begin{tabularx}{\linewidth}{|l|X|X|X|}
\hline
\rowcolor{popblueL} \textbf{Critère} & \textbf{Rapport} & \textbf{Compte rendu} & \textbf{Article de presse} \\
\hline
Destinataire & Une hiérarchie, une institution identifiée & Les absents, les membres d'un groupe & Le grand public \\
\hline
Ton & Objectif, analytique, administratif & Objectif, neutre, fidèle & Informatif, parfois accrocheur ou engagé \\
\hline
Structure & Introduction, constat, analyse, recommandations, conclusion & Restitution chronologique ou thématique des faits & Titre, chapeau, pyramide inversée (essentiel d'abord) \\
\hline
Finalité & Aider à décider & Informer de ce qui s'est dit ou passé & Informer, parfois susciter l'intérêt \\
\hline
\end{tabularx}
\end{center}

\begin{attention}[Erreur fréquente : confondre rapport et compte rendu]
Beaucoup d'élèves croient qu'un rapport se contente de raconter ce qui s'est passé, comme un compte rendu. C'est faux : le rapport comporte une \textbf{analyse} et souvent des \textbf{recommandations}, pas seulement une restitution. Le compte rendu dit : « Le prix du maïs a été fixé à 190 FCFA/kg. » Le rapport ajoute : « Ce prix, inférieur de 10 FCFA au prix du marché, risque de décourager les producteurs ; il est recommandé de le réviser à la prochaine assemblée. » Au Baccalauréat technique, un « rapport » sans analyse ni proposition est sanctionné comme hors sujet partiel.
\end{attention}

\courssub{La chaîne de communication du rapport}

Un rapport n'existe que dans une situation de communication précise : quelqu'un a observé quelque chose et en rend compte à quelqu'un qui doit décider. On peut schématiser cette chaîne ainsi :

\begin{popfigure}
\begin{center}
\begin{tikzpicture}[
box/.style={draw=popblue, fill=popblueL, rounded corners, minimum width=3.2cm, minimum height=1.3cm, align=center, font=\small\bfseries},
fleche/.style={-{Stealth[length=3mm]}, very thick, poporange},
note/.style={font=\footnotesize\itshape, text=popdark}]
\node[box, align=center] (redac) at (0,0){Rédacteur\\ (celui qui observe)};
\node[box, align=center] (dest) at (5.6,0){Destinataire\\ (hiérarchie, institution)};
\node[box, fill=popgreenL, draw=popgreen, align=center] (dec) at (11.2,0){Décision\\ (action à prendre)};
\draw[fleche] (redac) -- (dest) node[midway, above, note]{rapport : faits + analyse};
\draw[fleche] (dest) -- (dec) node[midway, above, note]{lecture, évaluation};
\draw[fleche, poppurple, dashed] (dec.south) to[bend left=25] node[midway, below, note]{suite donnée, retour d'information} (redac.south);
\end{tikzpicture}
\end{center}
\legende{La chaîne de communication du rapport : 1. le rédacteur observe une situation et la transmet ; 2. le destinataire lit et évalue ; 3. une décision est prise, qui peut revenir au rédacteur sous forme de consignes.}
\end{popfigure}

Ce schéma montre l'essentiel : le rapport est un \textbf{maillon} dans une chaîne de décision. S'il est confus ou inexact, toute la chaîne est bloquée — c'est pourquoi le chef de service de Marlyse a refusé son premier texte.

\courssub{Repérer la situation d'énonciation : un outil pour le Bac}

À l'épreuve de compréhension de texte du Probatoire et du Baccalauréat technique, on te demandera souvent d'analyser un texte fonctionnel. Le premier réflexe est d'identifier sa \cle{situation d'énonciation}, c'est-à-dire les circonstances dans lesquelles le texte a été produit.

\begin{methode}[Repérer la situation d'énonciation en trois questions]
\begin{enumerate}
\item \textbf{Qui écrit ?} Identifie l'émetteur : son nom, sa fonction, son statut (inspecteur, stagiaire, secrétaire de séance). Cherche les indices : signature, en-tête, « je » professionnel.
\item \textbf{À qui ?} Identifie le destinataire : supérieur hiérarchique, assemblée, administration. Indices : formule d'appel (« Monsieur le Directeur »), mention « à l'attention de », objet du document.
\item \textbf{Dans quel but ?} Détermine la fonction du texte : informer, justifier une dépense, proposer une mesure, rendre compte d'une mission. Indices : l'objet, les recommandations finales, les verbes d'action (« il est proposé de », « je recommande »).
\end{enumerate}
La réponse à ces trois questions donne la clé du texte : on comprend toujours mieux un document quand on sait qui parle, à qui, et pourquoi.
\end{methode}

\begin{exoresolu}[Analyse de la situation d'énonciation d'un texte fonctionnel]
\textbf{Énoncé.} Voici l'en-tête et la fin d'un document : « \emph{Rapport d'inspection sanitaire du marché Mokolo. À l'attention de Monsieur le Maire de la Commune d'arrondissement de Yaoundé II. [...] Il est proposé une mise en conformité des 36 étals défaillants sous quinze jours, sous peine de fermeture temporaire. Fait à Yaoundé, le 20 mars 2026. Le chef de l'équipe d'inspection, Ngo Bell Martine.} » Répondez aux trois questions de la situation d'énonciation.
\tcblower
\textbf{Solution détaillée.}
\begin{enumerate}
\item \emph{Qui écrit ?} Ngo Bell Martine, chef de l'équipe d'inspection de l'hygiène communale : un agent investi d'une mission officielle, ce qui donne autorité au document.
\item \emph{À qui ?} Monsieur le Maire de Yaoundé II : un supérieur hiérarchique, seul habilité à ordonner la fermeture d'étals. La mention « à l'attention de » est l'indice décisif.
\item \emph{Dans quel but ?} Rendre compte d'une inspection et obtenir une décision : la mise en conformité des 36 étals sous quinze jours. La formule « il est proposé » signale la recommandation, marque typique du rapport (et non du simple compte rendu).
\end{enumerate}
Conclusion : ce texte est bien un rapport administratif, écrit par un subordonné à son autorité de tutelle pour l'aider à décider.
\end{exoresolu}

\courssub{Bilan de la section et ouverture}

\begin{aretenir}[L'essentiel]
\begin{itemize}
\item Le \cle{rapport} est un texte fonctionnel, objectif et structuré, destiné à un destinataire identifié pour l'aider à décider.
\item Le \cle{texte informatif} repose sur quatre qualités : objectivité, exactitude, clarté, concision.
\item L'objectivité se lit dans la langue : présent de vérité générale, passé composé et imparfait factuels, « je » contrôlé, lexique dénotatif.
\item Le rapport se distingue du compte rendu (simple restitution) et de l'article de presse (grand public) par son analyse et ses recommandations.
\item Pour analyser tout texte fonctionnel, pose trois questions : qui écrit ? à qui ? dans quel but ?
\end{itemize}
\end{aretenir}

Marlyse sait désormais ce qu'on attend d'elle : un document objectif, exact, clair et concis, qui aidera son chef de service à décider du remplacement de la machine. Mais comment ce document est-il construit ? Quelles parties doit-il comporter, et dans quel ordre ? C'est l'objet de la prochaine section : la structure externe et interne du rapport.

\begin{exercice}[À vous de jouer]
Voici trois phrases. Indiquez pour chacune si elle peut figurer dans un rapport, et justifiez : a) « Le marché de Mokolo est vraiment dégoûtant, je n'y retournerai jamais. » b) « Sur 120 étals inspectés le 18 mars 2026, 36 présentent des manquements aux normes d'hygiène. » c) « Il y a pas mal de problèmes avec les étals depuis quelque temps. »
\end{exercice}

\begin{corrige}[Exercice : À vous de jouer]
a) Phrase refusée : opinion personnelle (« dégoûtant », « je n'y retournerai jamais »), lexique connotatif, aucune donnée vérifiable : elle viole l'objectivité. b) Phrase acceptée : date, chiffres exacts, fait vérifiable, ton neutre : elle respecte l'objectivité et l'exactitude. c) Phrase refusée : « pas mal » et « quelque temps » sont vagues ; l'exactitude exige des chiffres et des dates précis.
\end{corrige}

\coursec{La structure externe et interne du rapport}

Après avoir découvert ce qu'est un rapport et en quoi le texte informatif se distingue du texte littéraire ou argumentatif (section précédente), nous allons maintenant « ouvrir le capot » pour examiner l'architecture précise de ce genre professionnel. Comprendre la structure, c'est acquérir le squelette sur lequel viendra se greffer la chair de votre rédaction. Sans squelette solide, le rapport s'effondre ; sans chair, il reste une coquille vide.

\courssub{Réception d'un texte fonctionnel : un rapport de stage en BTP à Douala}

Commençons par l'observation d'un modèle authentique. Voici un extrait significatif d'un rapport de stage effectué chez \textbf{SOCIÉTÉ DE CONSTRUCTION ET DE TRAVAUX PUBLICS (SCTP)} à Douala, par un élève de Terminale Génie Civil.

\begin{exemplebox}[Extrait d'un rapport de stage — SCTP Douala]
\textbf{REPUBLIQUE DU CAMEROUN}\\
\textbf{Paix — Travail — Patrie}\\
\medskip
\begin{minipage}{\linewidth}
\textbf{LYCÉE TECHNIQUE DE BONABÉRI}\\
\textbf{Service des Stages et Insertion Professionnelle}
\end{minipage}
\medskip

\noindent\textbf{DOUALA, le 15 juillet 2025}\\
\medskip

\noindent\textbf{OBJET :} Rapport de stage de fin d'études — Période du 02 juin au 11 juillet 2025\\
\textbf{RÉF :} N° 045/LTB/STG/2025
\medskip

\noindent\textbf{Monsieur le Directeur Général de la SCTP}\\
\textbf{B.P. 1234 Douala}
\medskip

\begin{center}
\textbf{\Large RAPPORT DE STAGE DE FIN D'ÉTUDES}\\
\textbf{\large Exécuté à la Société de Construction et de Travaux Publics (SCTP)}\\
\textbf{\large Du 02 juin au 11 juillet 2025}
\end{center}
\medskip

\noindent\textbf{PRÉFACE / REMERCIEMENTS}\\
Je tiens à exprimer ma profonde gratitude à M. le Directeur Général de la SCTP...
\medskip

\noindent\textbf{INTRODUCTION}\\
Dans le cadre de la formation en Génie Civil option Bâtiment au Lycée Technique de Bonabéri, j'ai effectué un stage de six semaines...
\medskip

\noindent\textbf{I. PRÉSENTATION DU CADRE DU STAGE}\\
\textbf{I.1. Présentation de l'entreprise}\\
\textbf{I.2. Description du chantier : Construction du bâtiment administratif R+3 à Bonanjo}\\
\textbf{I.3. Organisation du travail et encadrement}
\medskip

\noindent\textbf{II. ACTIVITÉS RÉALISÉES ET CONSTATS TECHNIQUES}\\
\textbf{II.1. Suivi des travaux de terrassement et fondations}\\
\textbf{II.2. Participation au ferraillage et coulage des poteaux du rez-de-chaussée}\\
\textbf{II.3. Constats sur la gestion des matériaux et la sécurité sur le chantier}
\medskip

\noindent\textbf{III. ANALYSE ET INTERPRÉTATION DES DIFFICULTÉS RENCONTRÉES}\\
\textbf{III.1. Écarts entre le planning prévisionnel et l'exécution réelle}\\
\textbf{III.2. Problèmes d'approvisionnement en ciment CIMENCAM}\\
\textbf{III.3. Analyse des non-conformités sur le ferraillage (écartement des étriers)}
\medskip

\noindent\textbf{IV. RECOMMANDATIONS ET PROPOSITIONS D'AMÉLIORATION}\\
\textbf{IV.1. Mise en place d'un planning hebdomadaire révisé avec marges de sécurité}\\
\textbf{IV.2. Négociation d'un contrat-cadre avec deux fournisseurs de ciment}\\
\textbf{IV.3. Renforcement du contrôle qualité par un chef d'équipe dédié}
\medskip

\noindent\textbf{CONCLUSION}\\
Ce stage m'a permis de confronter la théorie apprise en classe à la réalité du terrain...
\medskip

\noindent\textbf{ANNEXES}\\
Annexe 1 : Plan de situation du chantier\\
Annexe 2 : Fiches de contrôle quotidien (semaines 1 à 6)\\
Annexe 3 : Attestation de stage signée par le chef de chantier
\medskip

\noindent\begin{minipage}{0.45\linewidth}
Fait à Douala, le 15 juillet 2025\\
\medskip
\textbf{Le Stagiaire}\\
\textit{Ngono Essomba Jean-Pierre}\\
Matricule : 2025-GC-014
\end{minipage}
\hfill
\begin{minipage}{0.45\linewidth}
Vu et validé par le Tuteur\\
\medskip
\textbf{M. Fouda André}\\
Chef de Chantier SCTP
\end{minipage}
\end{exemplebox}

\courssub{La structure externe : la disposition matérielle visible}

La structure externe est l'enveloppe physique du document. Elle répond aux normes de la communication administrative et technique. Elle permet au destinataire d'identifier immédiatement l'expéditeur, l'objet, la date et de naviguer dans le document.

\begin{definition}[Structure externe du rapport]
La \cle{structure externe} désigne l'ensemble des éléments matériels et visuels qui encadrent le texte : en-tête, titre, numérotation des parties, mise en page, signature et annexes. Elle assure l'identification administrative et la lisibilité du document.
\end{definition}

\begin{popfigure}
\begin{tikzpicture}[
    box/.style={draw=popblue, thick, rounded corners=4pt, fill=popblue!5, minimum height=1.8cm, text width=0.9\linewidth, align=center, font=\small},
    arrow/.style={-Stealth, thick, popdark},
    labelbox/.style={font=\tiny, text=popmuted, anchor=north}
]
\node[box, align=center] (header){
    \textbf{EN-TÊTE ADMINISTRATIF}\\
    République du Cameroun — Paix-Travail-Patrie\\
    Établissement / Service / Direction\\
    Lieu, Date (ex. Douala, le 15 juillet 2025)\\
    Objet : \textit{Précis, complet, avec références}\\
    Destinataire (qualité, nom, adresse)
};
\node[labelbox] at (header.north) {Bloc identitaire obligatoire};

\node[box, below=0.8cm of header, align=center] (title){
    \textbf{TITRE PRINCIPAL}\\
    \textit{RAPPORT DE STAGE DE FIN D'ÉTUDES}\\
    \textit{Exécuté à [Entreprise] — Période [Dates]}\\
    Centré, police grande, gras
};
\node[labelbox] at (title.north) {Identification du document};

\node[box, below=0.8cm of title, align=center] (intro){
    \textbf{INTRODUCTION}\\
    Contexte, mission, objectifs, démarche, plan annoncé\\
    \textit{Non numérotée, 1 à 2 pages max}
};
\node[labelbox] at (intro.north) {Porte d'entrée du contenu};

\node[box, below=0.8cm of intro, align=center] (body){
    \textbf{CORPS DU RAPPORT — PARTIES NUMÉROTÉES}\\
    \textbf{I. Contexte \& Cadre} (I.1, I.2, I.3)\\
    \textbf{II. Constats \& Observations} (II.1, II.2...)\\
    \textbf{III. Analyse \& Interprétation} (III.1, III.2...)\\
    \textbf{IV. Recommandations} (IV.1, IV.2...)\\
    \textit{Titres clairs, hiérarchie respectée}
};
\node[labelbox] at (body.north) {Cœur technique — Plan logique};

\node[box, below=0.8cm of body, align=center] (concl){
    \textbf{CONCLUSION}\\
    Bilan synthétique, portée des résultats, ouverture\\
    \textit{Non numérotée, ½ à 1 page}
};
\node[labelbox] at (concl.north) {Fermeture argumentée};

\node[box, below=0.8cm of concl, align=center] (sign){
    \textbf{SIGNATURES \& VALIDATION}\\
    Lieu, Date\\
    \textbf{Le Rédacteur} (Nom, Matricule, Signature)\\
    \textbf{Le Tuteur / Supérieur} (Nom, Qualité, Signature, Cachet)
};
\node[labelbox] at (sign.north) {Force probante du document};

\node[box, below=0.8cm of sign, fill=poporange!5, draw=poporange, align=center] (annex){
    \textbf{ANNEXES (facultatif mais fréquent)}\\
    Plans, photos, tableaux bruts, attestations, cahiers de charges\\
    \textit{Numérotées : Annexe 1, Annexe 2... Renvois dans le corps}
};
\node[labelbox] at (annex.north) {Preuves et compléments};

\draw[arrow] (header.south) -- (title.north);
\draw[arrow] (title.south) -- (intro.north);
\draw[arrow] (intro.south) -- (body.north);
\draw[arrow] (body.south) -- (concl.north);
\draw[arrow] (concl.south) -- (sign.north);
\draw[arrow] (sign.south) -- (annex.north);
\end{tikzpicture}
\legende{Maquette annotée de la structure externe d'un rapport technique (normes administratives camerounaises).}
\end{popfigure}

\textbf{Focus sur l'en-tête (la « tête de lettre »)}.\par
C'est la carte d'identité du document. Au Cameroun, l'en-tête officiel comporte \textbf{impérativement} :
\begin{itemize}
    \item La mention \textbf{« REPUBLIQUE DU CAMEROUN — Paix — Travail — Patrie »} (centrée, majuscules).
    \item L'identification de l'organisme émetteur (Lycée, Ministère, Entreprise, Service).
    \item Le \textbf{Lieu} et la \textbf{Date} de rédaction (ex. \textit{Yaoundé, le 20 octobre 2025}).
    \item L'\textbf{Objet} : une phrase nominale complète résumant le contenu (ex. \textit{Rapport de stage de fin d'études — Période du 02 juin au 11 juillet 2025}).
    \item La \textbf{Référence} (numéro d'ordre, service, année) pour l'archivage.
    \item Le \textbf{Destinataire} : qualité, nom, adresse (aligné à gauche).
\end{itemize}

\begin{attention}[Piège fréquent : l'objet incomplet ou absent]
Ne confondez pas \textbf{Objet} et \textbf{Titre}. L'objet est une mention administrative obligatoire dans l'en-tête (ex. \textit{Rapport d'incident sur le transformateur n°4 — Poste ENEO d'Ebolowa}). Le titre est la première ligne visible du corps du document, souvent plus descriptive. \textbf{Oublier l'objet} rend le document irrecevable administrativement. \textbf{Mélanger objet et titre} (ex. écrire seulement « Rapport » en objet) est une faute de professionnalisme sanctionnée au Probatoire.
\end{attention}

\textbf{La numérotation des parties : cohérence et hiérarchie}.\par
Deux systèmes coexistent. Choisissez-en \textbf{un seul} et tenez-le jusqu'au bout.

\begin{center}
\begin{tabularx}{\linewidth}{|l|X|X|}
\hline
\textbf{Système} & \textbf{Exemple} & \textbf{Usage recommandé} \\
\hline
Chiffres arabes (décimal) & 1. / 1.1 / 1.1.1 & Rapports techniques longs, normes ISO, facilité de renvoi (ex. « voir § 2.3.1 ») \\
\hline
Chiffres romains + arabes & I. / I.A / I.A.1 ou I. / I.1 / I.1.1 & Rapports administratifs, scolaires, juridiques (tradition française) \\
\hline
\end{tabularx}
\end{center}

\begin{propriete}[Règles d'or de la numérotation]
\begin{enumerate}
    \item \textbf{Un niveau n'existe que s'il a au moins deux éléments :} pas de I.1 sans I.2 ; pas de 1.1.1 sans 1.1.2.
    \item \textbf{Parallélisme des titres :} même structure grammaticale (tous infinitifs, ou tous nominaux, ou tous verbaux à la même personne).
    \item \textbf{Maximum 3 niveaux de profondeur} (ex. I, I.1, I.1.1). Au-delà, le lecteur se perd.
    \item \textbf{Numérotation continue} : pas de retour à 1 dans chaque partie (sauf annexes).
\end{enumerate}
\end{propriete}

\courssub{La structure interne : l'organisation logique des idées}

Si la structure externe est le squelette, la structure interne est la circulation sanguine : elle assure la cohérence du raisonnement. Un rapport technique suit une \cle{progression thématique} immuable : du contexte vers l'action, du fait vers l'interprétation, du diagnostic vers la prescription.

\begin{definition}[Structure interne du rapport]
La \cle{structure interne} est l'enchaînement logique des blocs de contenu qui guide le lecteur de la compréhension de la situation (contexte) vers la proposition d'action (recommandations), en passant par l'exposition des faits (constats) et leur exploitation intellectuelle (analyse).
\end{definition}

\begin{popfigure}
\begin{tikzpicture}[
    level 1/.style={sibling distance=6cm, level distance=3.5cm},
    level 2/.style={sibling distance=3cm, level distance=3cm},
    block/.style={draw=popblue, thick, rounded corners=6pt, fill=popblue!10, text width=3.2cm, align=center, minimum height=2.2cm, font=\small},
    subblock/.style={draw=popgreen, thick, rounded corners=4pt, fill=popgreen!10, text width=2.8cm, align=center, minimum height=1.8cm, font=\footnotesize},
    arrow/.style={-Stealth, thick, popdark},
    link/.style={dashed, popmuted}
]
\node[block, fill=popblue!20, align=center] (root){\textbf{CONTEXTE \& MISSION}\\(Introduction + Partie I)\\
\textit{« D'où je viens, quoi, pourquoi, comment »}}
    child {node[subblock] (c1) {Présentation de la structure d'accueil}}
    child {node[subblock] (c2) {Objectifs du stage / mission confiée}}
    child {node[subblock] (c3) {Moyens, planning, encadrement}};
    
\node[block, fill=poporange!15, below=4.5cm of root, xshift=-5cm, align=center] (obs){\textbf{CONSTATS \& OBSERVATIONS}\\(Partie II)\\
\textit{« Ce que j'ai vu, fait, mesuré » — FAITS BRUTS}} 
    child {node[subblock] (o1) {Activités réalisées (chronologique ou thématique)}}
    child {node[subblock] (o2) {Données techniques collectées (mesures, plans, essais)}}
    child {node[subblock] (o3) {Dysfonctionnements, écarts, incidents (factuels)}};
    
\node[block, fill=poppurple!15, below=4.5cm of root, xshift=5cm, align=center] (ana){\textbf{ANALYSE \& INTERPRÉTATION}\\(Partie III)\\
\textit{« Ce que cela signifie » — RAISONNEMENT}} 
    child {node[subblock] (a1) {Comparaison prévu / réel (écarts, causes racines)}}
    child {node[subblock] (a2) {Explication technique des problèmes (théorie vs terrain)}}
    child {node[subblock] (a3) {Impact sur coûts, délais, qualité, sécurité}};
    
\node[block, fill=popgreen!20, below=9cm of root, align=center] (rec){\textbf{RECOMMANDATIONS \& PROPOSITIONS}\\(Partie IV + Conclusion)\\
\textit{« Que faire maintenant ? » — ACTION}} 
    child {node[subblock] (r1) {Mesures correctives immédiates (curatif)}}
    child {node[subblock] (r2) {Mesures préventives / amélioration processus (structurel)}}
    child {node[subblock] (r3) {Perspectives pour l'entreprise / la formation}};

\draw[arrow, thick] (root.south) -- ++(0,-0.5) -| (obs.north);
\draw[arrow, thick] (root.south) -- ++(0,-0.5) -| (ana.north);
\draw[arrow, thick] (obs.south) -- ++(0,-0.5) -| (rec.north);
\draw[arrow, thick] (ana.south) -- ++(0,-0.5) -| (rec.north);

\node[draw=popgold, thick, rounded corners, fill=popgold!10, fit=(root)(obs)(ana)(rec), inner sep=10pt, label={[font=\small\bfseries, popdark]above:PROGRESSION THÉMATIQUE : CONTEXTE $\rightarrow$ CONSTATS $\rightarrow$ ANALYSE $\rightarrow$ RECOMMANDATIONS}] {};
\end{tikzpicture}
\legende{Arbre de la structure interne : les quatre blocs logiques et leur articulation. Les flèches indiquent le flux logique obligatoire.}
\end{popfigure}

\textbf{Les quatre blocs logiques en détail}.\par

\begin{enumerate}
    \item \textbf{Contexte \& Mission (Introduction + Partie I).} Répond à \textit{Pourquoi ce rapport ?} On y trouve : l'origine de la mission (programme scolaire, commande hiérarchique), l'identité de l'organisme d'accueil, la période, les objectifs assignés, les moyens mis à disposition. \emph{C'est le « contrat de lecture ».}
    \item \textbf{Constats \& Observations (Partie II).} Répond à \textit{Qu'est-ce qui s'est passé ?} C'est le \textbf{seul} endroit pour les faits bruts : description des activités, relevés de mesures, comptes-rendus de réunions, photos légendées, incidents \emph{décrits sans jugement}. \textbf{Interdiction formelle :} y glisser une analyse ou une critique.
    \item \textbf{Analyse \& Interprétation (Partie III).} Répond à \textit{Qu'est-ce que cela veut dire ?} Ici, on mobilise les connaissances techniques (cours, normes, DTU) pour expliquer les constats : recherche des causes racines (méthode 5M, arbre des causes), évaluation de la gravité, comparaison aux normes. C'est la \textbf{valeur ajoutée intellectuelle} du rapport.
    \item \textbf{Recommandations \& Propositions (Partie IV).} Répond à \textit{Que faut-il faire ?} Elles découlent \textbf{logiquement} de l'analyse. Chaque recommandation doit être : \textbf{Spécifique}, \textbf{Mesurable}, \textbf{Atteignable}, \textbf{Réaliste}, \textbf{Temporelle} (critères SMART). On distingue correctif (urgence) et préventif (durable).
\end{enumerate}

\begin{methode}[Construire le plan d'un rapport en 4 blocs logiques]
\begin{enumerate}
    \item \textbf{Listez vos matières premières :} notes de terrain, photos, mesures, entretiens, documents d'entreprise.
    \item \textbf{Triage :} Classez chaque élément dans l'un des 4 tiroirs : \textit{Contexte} (cadrage), \textit{Constat} (fait brut), \textit{Analyse} (lien fait/théorie), \textit{Recommandation} (action future).
    \item \textbf{Regroupement thématique :} Dans chaque tiroir, regroupez par thème (ex. : Sécurité, Planning, Qualité, Gestion RH). Chaque thème = une sous-partie (I.1, I.2...).
    \item \textbf{Rédaction des titres :} Formulez des titres \textbf{parallèles} et \textbf{informatifs} (pas « Problèmes » mais « Écarts planning et causes »).
    \item \textbf{Vérification de la chaîne logique :} Chaque recommandation (IV.x) doit trouver sa cause dans l'analyse (III.x) qui explique un constat (II.x) situé dans le contexte (I.x).
    \item \textbf{Numérotation finale :} Appliquez le système choisi (I, I.1 ou 1, 1.1) sur toute la profondeur.
\end{enumerate}
\end{methode}

\courssub{Exemple chiffré : plan détaillé commenté — Stage chez CAMWATER}

Voici un plan complet (3 parties principales, 8 sous-parties) pour un rapport de stage « Étude de la qualité de l'eau distribuée au réseau de Yaoundé-Nord » chez CAMWATER. Observez la progression et la numérotation.

\begin{exoresolu}[Plan détaillé — Rapport de stage CAMWATER Yaoundé]
\textbf{INTRODUCTION}\\
Contexte : Stage de fin d'études BTS Génie de l'Eau. Mission : évaluer la conformité physico-chimique et bactériologique de l'eau au robinet dans 5 quartiers. Durée : 8 semaines.
\medskip

\noindent\textbf{I. CADRE ET MÉTHODOLOGIE DE L'ÉTUDE}\\
\textbf{I.1. Présentation de CAMWATER et du laboratoire central d'Odza}\\
\textbf{I.2. Description du périmètre d'étude : 5 quartiers, 50 points de prélèvement}\\
\textbf{I.3. Protocole de prélèvement et normes de référence (OMS, norme NC 234)}
\medskip

\noindent\textbf{II. RÉSULTATS D'ANALYSES ET CONSTATS DE TERRAIN}\\
\textbf{II.1. Paramètres physico-chimiques : pH, turbidité, chlore résiduel, conductivité}\\
\textbf{II.2. Paramètres bactériologiques : Coliformes totaux, E. coli, Entérocoques}\\
\textbf{II.3. Observations de terrain : état des canalisations, pressions, plaintes usagers}
\medskip

\noindent\textbf{III. ANALYSE CRITIQUE ET INTERPRÉTATION DES NON-CONFORMITÉS}\\
\textbf{III.1. Corrélation turbidité / chlore résiduel faible dans les zones périphériques}\\
\textbf{III.2. Identification de 3 points de contamination fécale (rupture de canalisation, reflux)}\\
\textbf{III.3. Écart entre fréquence de contrôle réglementaire (hebdo) et réel (mensuel)}
\medskip

\noindent\textbf{IV. RECOMMANDATIONS TECHNIQUES ET ORGANISATIONNELLES}\\
\textbf{IV.1. Réhabilitation urgente des 3 tronçons défectueux identifiés (quartiers Etoudi, Nkolbisson)}\\
\textbf{IV.2. Installation de 2 postes de rechloration automatique en bout de réseau}\\
\textbf{IV.3. Révision du plan de surveillance : passage à 2 prélèvements/semaine par zone}\\
\textbf{IV.4. Sensibilisation des usagers sur le signalement des anomalies (numéro vert)}
\medskip

\noindent\textbf{CONCLUSION}\\
Bilan : 86\% de conformité globale, mais poches de risque sanitaire majeures. Propositions chiffrées (coût estimé 45 millions FCFA). Ouverture : extension de la méthode à Yaoundé-Sud.
\medskip

\noindent\textbf{ANNEXES}\\
Annexe 1 : Carte de localisation des 50 points. Annexe 2 : Tableaux bruts d'analyses (6 semaines). Annexe 3 : Photos des ruptures. Annexe 4 : Devis estimatifs travaux.
\tcblower
\textbf{Commentaire du professeur :}
\begin{itemize}
    \item \textbf{Parallélisme parfait :} I.1/I.2/I.3 (Cadre), II.1/II.2/II.3 (Données brutes), III.1/III.2/III.3 (Explication), IV.1/IV.2/IV.3/IV.4 (Actions).
    \item \textbf{Chaîne logique traçable :} Le constat « turbidité élevée + chlore faible » (II.1) $\rightarrow$ Analysé comme « risque de protection virale insuffisante » (III.1) $\rightarrow$ Recommandation « poste de rechloration » (IV.2).
    \item \textbf{Titres informatifs :} On comprend le contenu sans lire le paragraphe.
    \item \textbf{Équilibre :} 3 parties principales, 3 à 4 sous-parties chacune. Ni trop fin, ni trop gros.
\end{itemize}
\end{exoresolu}

\courssub{Activité guidée : reconstitution d'un rapport mélangé}

\paragraph{Consigne.}
Les rubriques ci-dessous proviennent d'un rapport de stage « Maintenance préventive des groupes électrogènes — Centrale thermique de Limbé ». Elles sont mélangées. \textbf{1.} Remettez-les dans l'ordre de la structure externe. \textbf{2.} Attribuez à chaque rubrique de corps (II, III, IV) son bloc logique interne (Contexte, Constats, Analyse, Recommandations). \textbf{3.} Justifiez vos choix.

\begin{center}
\begin{tabularx}{\linewidth}{|>{\bfseries}l|X|}
\hline
\textbf{Rubrique mélangée} & \textbf{Contenu indicatif} \\
\hline
A. Annexe 2 : Fiches de maintenance constructeur (CATERPILLAR) & Documents fournisseur \\
B. III.2. Analyse des pannes récurrentes : filtres à gasoil encrassés & Cause racine identifiée \\
C. Objet : Rapport de stage maintenance G.E — Juin/Juillet 2025 & En-tête \\
D. II.1. Planning des interventions préventives réalisées (Semaines 1-6) & Activités faites \\
E. IV.1. Proposition de stock de sécurité filtres (niveau min/max) & Action corrective \\
F. Conclusion : Bilan compétences acquises + ouverture sur GMAO & Fermeture \\
G. I.2. Organisation du service maintenance : 3 équipes tournantes & Cadre du stage \\
H. Signature stagiaire / Validé par Chef Maintenance & Validation \\
I. II.3. Relevés de paramètres : T° huile, pression huile, T° eau, tension & Données brutes mesurées \\
J. Introduction : Contexte formation, objectifs stage, plan annoncé & Porte d'entrée \\
K. III.1. Corrélation entre non-respect délais vidange et usure prématurée & Raisonnement technique \\
L. Titre principal : RAPPORT DE STAGE MAINTENANCE PRÉVENTIVE G.E — CENTRALE LIMBÉ & Titre \\
M. I.1. Présentation de la centrale thermique de Limbé (4 $\times$ 25 MW) & Contexte entreprise \\
N. IV.2. Mise en place d'un planning GMAO avec alertes automatiques & Action préventive \\
O. II.2. Constats visuels : fuites injecteurs groupe n°3, corrosion bac rétention & Faits observés \\
\hline
\end{tabularx}
\end{center}

\begin{corrige}[Exercice de reconstitution — Centrale Limbé]
\textbf{1. Ordre structure externe (squelette) :}
\begin{enumerate}
    \item \textbf{C} (Objet — en-tête) $\rightarrow$ \textbf{L} (Titre principal) $\rightarrow$ \textbf{J} (Introduction) $\rightarrow$ \textbf{M} (I.1) $\rightarrow$ \textbf{G} (I.2) $\rightarrow$ \textbf{D} (II.1) $\rightarrow$ \textbf{I} (II.2) $\rightarrow$ \textbf{O} (II.3) $\rightarrow$ \textbf{K} (III.1) $\rightarrow$ \textbf{B} (III.2) $\rightarrow$ \textbf{E} (IV.1) $\rightarrow$ \textbf{N} (IV.2) $\rightarrow$ \textbf{F} (Conclusion) $\rightarrow$ \textbf{H} (Signatures) $\rightarrow$ \textbf{A} (Annexes).
\end{enumerate}

\textbf{2. Classification structure interne (logique) :}
\begin{itemize}
    \item \textbf{Contexte (Partie I) :} M (Présentation centrale), G (Organisation service). \textit{Justification :} Cadrent le « où » et le « comment » de la mission.
    \item \textbf{Constats (Partie II) :} D (Planning réalisé — fait), I (Relevés paramètres — mesures brutes), O (Fuites, corrosion — observations visuelles). \textit{Justification :} Données factuelles, sans interprétation.
    \item \textbf{Analyse (Partie III) :} K (Corrélation vidange/usure — lien cause/effet technique), B (Analyse pannes filtres — diagnostic). \textit{Justification :} Mobilisation savoirs techniques pour expliquer les constats II.
    \item \textbf{Recommandations (Partie IV) :} E (Stock sécurité — réponse directe à B), N (GMAO — réponse systémique à K). \textit{Justification :} Actions concrètes issues de l'analyse.
\end{itemize}

\textbf{3. Points de vigilance :} Remarquez que II.2 (relevés) précède II.3 (constats visuels) : ordre logique « mesures instruments » puis « inspection visuelle ». L'analyse III.1 porte sur la maintenance (vidange), III.2 sur la filtration : deux thèmes distincts, deux sous-parties. Les recommandations IV.1 et IV.2 répondent chacune à une analyse précise.
\end{corrige}

\begin{savaistu}[Les mentions officielles sur les rapports camerounais]
Au Cameroun, tout rapport émanant d'une administration publique, d'un établissement public ou d'un lycée technique porte \textbf{obligatoirement} en en-tête la devise nationale : \textbf{« REPUBLIQUE DU CAMEROUN — Paix — Travail — Patrie »}. Cette mention confère au document sa \cle{valeur probante} et son authenticité administrative. Elle est souvent accompagnée des \textbf{armes du Cameroun} (étoile, fasces, devise). Les références (ex. \textit{N° 0045/MINESEC/SG/DRETS/2025}) suivent une codification stricte : numéro d'ordre / Ministère / Direction / Service / Année. Omettre cette mention ou la devise expose le document à une irrecevabilité formelle lors d'un examen (Probatoire, BTS) ou d'un dépôt hiérarchique.
\end{savaistu}

\courssub{Synthèse et transition}

Vous maîtrisez maintenant l'architecture complète du rapport : l'\textbf{enveloppe externe} (en-tête normalisé, titre, numérotation rigoureuse, signatures, annexes) qui garantit sa recevabilité administrative, et le \textbf{moteur interne} (Contexte $\rightarrow$ Constats $\rightarrow$ Analyse $\rightarrow$ Recommandations) qui garantit sa pertinence technique. La prochaine étape consistera à outiller votre plume : comment choisir les mots justes — \cle{lexique commun} vs \cle{lexique spécialisé} — pour que ce squelette respire la compétence professionnelle. C'est l'objet de la section 3.

\coursec{Lexique commun et lexique spécialisé}

Dans la section précédente, nous avons appris à construire l'ossature du rapport : sa structure externe (page de garde, sommaire, introduction, corps, conclusion, annexes) et sa structure interne (l'enchaînement logique des idées). Mais un rapport bien structuré peut être gâché par un mauvais choix de mots. Écrire « la machine ne marche pas bien » dans un rapport de maintenance, c'est comme dessiner un plan sans cotes : le lecteur ne peut rien en faire. Cette section vous apprendra à choisir le bon vocabulaire : celui que tout le monde comprend, et celui, plus précis, qui appartient à votre métier.

\courssub{Intuition : deux façons de dire la même chose}

Imaginez deux techniciens qui constatent la même panne sur un groupe électrogène.

Le premier dit : « Le moteur fait un drôle de bruit et il chauffe. »

Le second dit : « Le moteur présente un cognement au ralenti et une surchauffe du liquide de refroidissement au-delà de 95 degrés. »

Le premier message est compréhensible par tous, mais il est inutilisable : quel bruit ? quelle température ? Le second message utilise des mots précis du métier : il permet au réparateur de diagnostiquer, de commander la bonne pièce, d'agir. Toute la différence tient au \cle{lexique} employé. C'est ce que nous allons maintenant définir rigoureusement.

\courssub{Définitions et critères de distinction}

\begin{definition}[Lexique commun]
Le \cle{lexique commun} est l'ensemble des mots de la langue française que tout locuteur comprend et emploie dans la vie quotidienne, indépendamment de son métier : \emph{eau}, \emph{machine}, \emph{travail}, \emph{maison}, \emph{malade}. Ces mots sont fréquents, polyvalents et accessibles à tout destinataire.
\end{definition}

\begin{definition}[Lexique spécialisé]
Le \cle{lexique spécialisé} (ou vocabulaire technique) est l'ensemble des termes propres à un domaine d'activité, à une science ou à un métier : \emph{chiffre d'affaires} (comptabilité), \emph{béton armé} (génie civil), \emph{paludisme} (médecine), \emph{fibre optique} (télécommunications). Ces termes désignent des réalités précises et ne sont pleinement compris que par les spécialistes du domaine.
\end{definition}

\begin{aretenir}[Les trois critères de distinction]
Pour décider si un mot relève du lexique commun ou du lexique spécialisé, on applique trois critères :
\begin{enumerate}
\item \textbf{La fréquence d'usage} : le mot apparaît-il dans la conversation quotidienne (\emph{eau}) ou seulement dans un contexte professionnel (\emph{ferraillage}) ?
\item \textbf{Le domaine d'emploi} : le mot est-il utilisable partout, ou lié à un champ d'activité précis (médecine, bâtiment, informatique) ?
\item \textbf{Le destinataire} : un non-spécialiste comprend-il le mot sans explication ? Si non, c'est un terme spécialisé.
\end{enumerate}
\end{aretenir}

\courssub{Observation : repérage dans des rapports authentiques}

Appliquons ces critères à deux extraits de rapports réels, comme vous le ferez en séance de réception des textes.

\textbf{Extrait 1 — Rapport médical, hôpital de district de Bafang.} « Au cours du trimestre, le service a enregistré 214 cas de \emph{paludisme} confirmés par \emph{test de diagnostic rapide (TDR)}. Les patients ont reçu une \emph{antibiothérapie} adaptée ; le taux de \emph{létalité} reste inférieur à 2 \%. »

Termes spécialisés à souligner : \emph{paludisme}, \emph{test de diagnostic rapide (TDR)}, \emph{antibiothérapie}, \emph{létalité}. Mots du lexique commun : \emph{trimestre}, \emph{service}, \emph{cas}, \emph{patients}, \emph{taux}.

\textbf{Extrait 2 — Rapport technique, chantier routier Yaoundé–Nsimalen.} « Le \emph{décapage} de la plate-forme est achevé sur le PK 12. La \emph{couche de fondation} en \emph{grave non traitée} a été mise en œuvre et compactée à 95 \% de l'\emph{optimum Proctor}. »

Termes spécialisés : \emph{décapage}, \emph{plate-forme}, \emph{couche de fondation}, \emph{grave non traitée}, \emph{optimum Proctor}. Mots communs : \emph{achevé}, \emph{mis en œuvre}, \emph{compactée}.

\begin{attention}[Erreur fréquente]
Ne confondez pas \emph{mot difficile} et \emph{terme spécialisé}. Un mot savant comme \emph{fastidieux} appartient au lexique commun soutenu : tout le monde peut le comprendre. Un terme spécialisé comme \emph{optimum Proctor} n'a de sens que dans un métier. Dans un rapport, l'erreur grave est d'employer un terme spécialisé \textbf{sans le définir} quand le destinataire n'est pas du métier : le proviseur qui lit le rapport du stage n'a pas besoin de savoir ce qu'est l'optimum Proctor, sauf si vous le lui expliquez en une périphrase.
\end{attention}

\courssub{La zone de recouvrement : les mots à double usage}

Certains mots vivent dans les deux mondes : ils ont un sens courant et un sens technique différent. C'est l'intersection des deux lexiques.

\begin{popfigure}
\begin{center}
\begin{tikzpicture}
\draw[fill=popblueL,opacity=0.7] (-1.4,0) circle (2.2);
\draw[fill=poporangeL,opacity=0.7] (1.4,0) circle (2.2);
\node[popdark,font=\bfseries] at (-2.6,1.7) {Lexique commun};
\node[popdark,font=\bfseries] at (2.6,1.7) {Lexique spécialisé};
\node[align=center,font=\small] at (-2.6,0.3) {eau\\ travail\\ machine\\ malade};
\node[align=center,font=\small] at (2.6,0.3) {ferraillage\\ létalité\\ amortissement\\ optimum Proctor};
\node[align=center,font=\small\itshape] at (0,0.3) {mémoire\\ réseau\\ tension\\ chiffre};
\node[align=center,font=\footnotesize] at (0,-1.2) {mots à double usage};
\end{tikzpicture}
\end{center}
\legende{Diagramme de Venn : le lexique commun et le lexique spécialisé se recouvrent partiellement. Les mots de l'intersection changent de sens selon le domaine : la \emph{tension} (électricité), la \emph{mémoire} (informatique), le \emph{réseau} (télécommunications).}
\end{popfigure}

\begin{savaistu}[Des mots ordinaires devenus techniques]
Certains mots du lexique commun deviennent spécialisés selon le domaine. La \emph{tension}, c'est d'abord un état d'énervement ; en électricité, c'est la différence de potentiel mesurée en volts. La \emph{mémoire}, c'est la faculté de se souvenir ; en informatique, c'est le composant qui stocke les données. Le \emph{chiffre}, c'est un symbole de 0 à 9 ; en comptabilité, le \emph{chiffre d'affaires} désigne le total des ventes d'une entreprise. En génie civil, le \emph{coffrage} n'a rien à voir avec un cercueil : c'est le moule dans lequel on coule le béton. Dans un rapport, c'est toujours le \textbf{contexte} qui fixe le sens.
\end{savaistu}

\courssub{Fonctions du lexique spécialisé dans le rapport}

Pourquoi un bon rapport ne peut-il pas se passer du vocabulaire technique ? Trois fonctions essentielles le justifient.

\begin{enumerate}
\item \textbf{La précision.} Le terme spécialisé désigne une réalité unique, sans ambiguïté. « Surchauffe du liquide de refroidissement à 95 degrés » est mesurable et vérifiable ; « ça chauffe » ne l'est pas.
\item \textbf{La crédibilité professionnelle.} Un rapport qui emploie correctement le vocabulaire du métier montre que son auteur maîtrise son domaine. Un chef de chantier qui écrit « ferraillage » et « béton dosé à 350 kg/m³ » inspire confiance ; celui qui écrit « le fer et le ciment » inspire le doute.
\item \textbf{L'économie de mots.} Un seul terme technique remplace souvent une longue explication : « amortissement » dit en un mot « répartition du coût d'un équipement sur sa durée d'utilisation ».
\end{enumerate}

\begin{attention}[Les risques : jargon et incompréhension]
Le lexique spécialisé a deux dangers. Le premier est le \cle{jargon} inutile : empiler des termes techniques pour paraître savant, là où un mot simple suffirait. Le second est l'emploi d'un terme spécialisé \textbf{incompris du destinataire} : le rapport devient alors illisible. La solution est simple : à la première occurrence d'un terme technique adressé à un non-spécialiste, ajoutez une \textbf{définition} ou une \textbf{périphrase explicative}. Exemple : « la grave non traitée, c'est-à-dire un mélange de graviers concassés sans liant, a été compactée ».
\end{attention}

\begin{methode}[Adapter le niveau de langue technique au destinataire]
\begin{enumerate}
\item \textbf{Identifier le destinataire} du rapport : est-ce un spécialiste (chef de service, contremaître, médecin) ou un non-spécialiste (directeur, client, autorité administrative) ?
\item \textbf{Si le destinataire est un spécialiste} : employez librement les termes techniques du domaine, sans les définir ; ils garantissent précision et économie.
\item \textbf{Si le destinataire est un non-spécialiste} : privilégiez le lexique commun ; pour chaque terme spécialisé indispensable, donnez une définition brève ou une périphrase entre virgules ou entre parenthèses.
\item \textbf{Si le rapport a plusieurs destinataires} : rédigez le corps en langage accessible et placez les données très techniques en annexe (tableaux, fiches de mesures).
\item \textbf{Relire en se mettant à la place du lecteur} : chaque mot qu'il ne comprendra pas doit être remplacé ou expliqué.
\end{enumerate}
\end{methode}

\courssub{Champs lexicaux et réseaux de mots : le vocabulaire de votre filière}

Le vocabulaire spécialisé ne s'apprend pas mot par mot : il s'organise en \cle{champ lexical}, c'est-à-dire en un réseau de mots qui tournent autour d'une même notion et s'éclairent mutuellement. Chaque filière technique possède le sien. Pour le construire, partez du mot central (le domaine), puis ramifiez : les lieux, les matériaux, les outils, les actions, les acteurs, les mesures.

\begin{popfigure}
\begin{center}
\begin{tikzpicture}[mindmap, grow cyclic, every node/.style={concept}, concept color=popblueL,
level 1/.append style={level distance=3.2cm, sibling angle=72},
level 2/.append style={level distance=2cm, sibling angle=45}]
\node[concept color=poporangeL, font=\bfseries] {Génie civil}
child[concept color=popgreenL] { node {Chantier} child { node[font=\small] {plate-forme} } child { node[font=\small] {terrassement} } }
child[concept color=popgreenL] { node {Matériaux} child { node[font=\small] {béton armé} } child { node[font=\small] {granulat} } }
child[concept color=popgreenL] { node {Ouvrages} child { node[font=\small] {ferraillage} } child { node[font=\small] {coffrage} } }
child[concept color=popgreenL] { node {Mesures} child { node[font=\small] {nivellement} } child { node[font=\small] {dosage} } }
child[concept color=popgreenL] { node {Acteurs} child { node[font=\small] {maçon} } child { node[font=\small] {conducteur de travaux} } };
\end{tikzpicture}
\end{center}
\legende{Carte mentale du champ lexical du génie civil : autour du domaine central, cinq branches (chantier, matériaux, ouvrages, mesures, acteurs) regroupent les termes spécialisés essentiels. Procédez de même pour votre filière : comptabilité (bilan, amortissement, créance), secrétariat (correspondance, classement, convocation), agriculture (assolement, intrant, rendement).}
\end{popfigure}

\begin{exemplebox}[Champs lexicaux par filière]
\begin{itemize}
\item \textbf{Comptabilité} : bilan, journal, grand livre, amortissement, créance, dette, chiffre d'affaires, TVA.
\item \textbf{Secrétariat} : correspondance, convocation, ordre du jour, procès-verbal, classement, archivage.
\item \textbf{Agriculture} : assolement, intrant, semence, rendement, pépinière, itinéraire technique.
\item \textbf{Électrotechnique} : tension, intensité, disjoncteur, court-circuit, schéma unifilaire.
\end{itemize}
\end{exemplebox}

\courssub{Application guidée : du vague au précis}

Voyons maintenant, sur un cas complet, comment le choix du lexique transforme un constat inutilisable en information professionnelle.

\begin{exoresolu}[Réécriture d'un constat de maintenance]
Énoncé. Dans le rapport de maintenance du groupe électrogène d'un lycée technique, un stagiaire a écrit : « La machine ne marche pas bien, elle fait du bruit et consomme beaucoup. » Réécrivez ce constat en langage technique précis, puis justifiez vos choix.
\tcblower
Solution détaillée. La phrase du stagiaire souffre de trois imprécisions : « machine » (quelle machine ?), « ne marche pas bien » (quel dysfonctionnement ?), « du bruit » et « beaucoup » (combien ?). Version technique :

\emph{« Le groupe électrogène de 30 kVA présente un cognement anormal au ralenti, probablement dû à l'usure des injecteurs. La consommation de gasoil est passée de 4 à 6 litres par heure, soit une hausse de 50 \%. Il est recommandé de remplacer les injecteurs et de vérifier le calage de la pompe à injection. »}

Justification : (1) « groupe électrogène de 30 kVA » identifie précisément l'équipement ; (2) « cognement au ralenti » nomme le symptôme avec le vocabulaire du métier ; (3) les mesures chiffrées (4 à 6 litres par heure, +50 \%) rendent le constat vérifiable ; (4) la conclusion propose une action, ce qui est la fonction d'un rapport de maintenance. Le passage du lexique commun au lexique spécialisé a transformé une plainte en diagnostic.
\end{exoresolu}

\begin{exercice}[À vous de jouer]
\begin{enumerate}
\item Dans l'extrait du rapport médical de Bafang cité plus haut, relevez trois termes spécialisés et proposez pour chacun une périphrase explicative destinée à un lecteur non médecin.
\item Construisez la carte mentale du champ lexical de votre propre filière : un mot central, au moins quatre branches et deux mots par branche.
\item Réécrivez en langage technique cette phrase extraite d'un rapport de stage en comptabilité : « L'entreprise a vendu pour pas mal d'argent cette année, mais elle a aussi pas mal de dettes. »
\end{enumerate}
\end{exercice}

\begin{corrige}[Exercice — Lexique commun et lexique spécialisé]
\begin{enumerate}
\item \emph{Paludisme} : maladie transmise par les piqûres de moustiques ; \emph{test de diagnostic rapide (TDR)} : examen qui détecte la maladie en quelques minutes à partir d'une goutte de sang ; \emph{létalité} : proportion de malades qui décèdent de la maladie.
\item Exemple pour la filière secrétariat : mot central « secrétariat » ; branches « documents » (correspondance, procès-verbal), « organisation » (classement, archivage), « réunions » (convocation, ordre du jour), « outils » (traitement de texte, agenda partagé).
\item « L'entreprise a réalisé un chiffre d'affaires de 85 millions de francs CFA au cours de l'exercice, en hausse de 12 \%, mais son endettement atteint 40 millions de francs CFA, soit 47 \% du chiffre d'affaires. » (Toute formulation chiffrée et employant \emph{chiffre d'affaires}, \emph{exercice}, \emph{endettement} est acceptée.)
\end{enumerate}
\end{corrige}

\begin{aretenir}[L'essentiel de la section]
Le \cle{lexique commun} rassemble les mots compris de tous ; le \cle{lexique spécialisé} rassemble les termes précis d'un métier. Dans un rapport, le vocabulaire technique assure la précision, la crédibilité et l'économie de mots, à condition d'être adapté au destinataire : défini ou périphrasé pour un non-spécialiste, employé librement entre spécialistes. Certains mots (\emph{tension}, \emph{mémoire}, \emph{réseau}) appartiennent aux deux lexiques et changent de sens selon le contexte. Maîtriser le champ lexical de sa filière, c'est se donner l'outil essentiel pour rédiger l'introduction, le corps et la conclusion du rapport — ce que nous ferons dans la prochaine section.
\end{aretenir}

\coursec{Rédiger l'introduction et la conclusion}

\courssub{Transition : du choix des mots à l'architecture du discours}

Après avoir distingué le \cle{lexique commun} du \cle{lexique spécialisé} (section précédente), nous disposons désormais du « matériau » linguistique adapté au texte fonctionnel. Encore faut-il savoir l'organiser dans une architecture rigoureuse. L'introduction et la conclusion ne sont pas de simples politesses rédactionnelles : ce sont les \textbf{deux portes d'entrée et de sortie} du rapport. Une introduction mal calibrée perd le lecteur (le décideur) ; une conclusion floue neutralise l'utilité du travail accompli. Cette section vous apprend à les construire comme un technicien monte un châssis : par étapes logiques, avec des pièces standardisées et un contrôle qualité final.

\courssub{Réception d'un texte fonctionnel n° 4 : analyse d'un rapport d'incident industriel}

\begin{exemplebox}[Rapport d'incident — Usine de transformation agro-alimentaire, Douala-Bassa (extrait)]
\textbf{Objet :} Rapport d'incident sur la rupture de la chaîne de froid (salle de conditionnement n°3) — Décision n° 2024/0112/DG du 10/01/2024.

\textbf{Introduction}
Conformément à la décision n° 2024/0112/DG du 10 janvier 2024, portant nomination d'une commission d'enquête suite à l'alerte déclenchée le 8 janvier à 03h17 sur le système de supervision SCADA, il m'a été demandé, en ma qualité de responsable maintenance QHSE, d'établir un rapport circonstancié sur les causes de la rupture de la chaîne de froid en salle de conditionnement n°3 et de formuler des recommandations correctives.
La mission, limitée au périmètre des installations frigorifiques du bâtiment B (compresseurs C1 à C4, échangeurs, armoires de commande), s'est déroulée du 11 au 15 janvier 2024. Elle a reposé sur trois leviers : (1) l'exploitation des fichiers historisés du système SCADA (températures, pressions, états marche/arrêt), (2) l'inspection visuelle et vibratoire des équipements, (3) les entretiens structurés avec l'équipe de quart (chef de quart, opérateurs, électricien de maintenance).
Ce rapport présente successivement : l'analyse chronologique de l'incident, le diagnostic des causes racines (méthode des 5M / arbre des causes), l'évaluation des pertes produits (référentiel ISO 22000) et les mesures correctives priorisées selon la criticité.

\textbf{Conclusion}
Au terme de cette étude, il apparaît que la rupture de la chaîne de froid trouve son origine principale dans la défaillance simultanée du pressostat basse pression du compresseur C2 (usure mécanique non détectée par la GMAO) et de l'algorithme de basculement automatique sur le compresseur de secours C4 (erreur de paramétrage du seuil de déclenchement à 1,2 bar au lieu de 1,8 bar). L'absence de test de basculement mensuel, préconisé par le constructeur mais non intégré au plan de maintenance préventive, a masqué cette double défaillance latente.
Nous recommandons en conséquence : (1) le remplacement immédiat du pressostat C2 et la révision du paramétrage du basculement C4 (seuil 1,8 bar) avant le 31/01/2024 ; (2) l'intégration du test de basculement mensuel dans le plan de maintenance préventive (GMAO) avec traçabilité signée ; (3) la formation de l'équipe de quart sur la procédure de basculement manuel d'urgence (fiche OP-URG-07) d'ici le 15/02/2024. La direction validera le plan d'actions par visa en marge de ce rapport.
\end{exemplebox}

\begin{aretenir}[Analyse structurelle de l'extrait]
L'introduction répond aux \textbf{quatre questions obligatoires} : \textbf{Contexte} (décision DG, alerte SCADA, date), \textbf{Objet et limites} (rupture de la chaîne de froid, salle n°3, bâtiment B), \textbf{Méthode} (SCADA, inspection, entretiens), \textbf{Plan} (chronologie, diagnostic 5M, pertes, mesures correctives). La conclusion opère un \textbf{aller-retour} : bilan synthétique (causes racines : pressostat + paramétrage + absence de test) $\rightarrow$ recommandations opérationnelles datées, chiffrées, responsabilisées. Aucune information nouvelle, aucune excuse, aucun « je pense que ».
\end{aretenir}

\courssub{L'introduction du rapport : la méthode en quatre temps (entonnoir)}

L'introduction guide le lecteur du \textbf{général} (le cadre administratif) vers le \textbf{spécifique} (le plan détaillé du rapport). C'est un \textbf{entonnoir} : chaque phrase resserre le champ.

\begin{methode}[Rédiger une introduction de rapport en quatre temps]
\textbf{Temps 1 — Rappeler le contexte administratif (le « pourquoi » officiel).}
Citez l'acte qui déclenche la mission : décision, arrêté, lettre de mission, contrat, procès-verbal. Donnez la date, le numéro, l'autorité émettrice. Précisez votre qualité (stagiaire, chef de mission, expert).
\emph{Formule type :} « Conformément à la décision n°... du..., il m'a été demandé de... » / « Dans le cadre de la convention n°... signée le... entre... et..., j'ai effectué... »

\textbf{Temps 2 — Présenter l'objet et les limites de la mission (le « quoi » et le « où »).}
Définissez précisément le sujet (objet) et ce qui en est exclu (limites : périmètre géographique, technique, temporel, budgétaire).
\emph{Formule type :} « Ce rapport porte sur... » / « La mission s'inscrit dans le périmètre suivant : ... Sont exclus : ... »

\textbf{Temps 3 — Annoncer la méthode (le « comment »).}
Listez les sources et démarches : visites de terrain, entretiens (avec qui, combien), documents consultés (références, dates), mesures effectuées (normes, appareils), enquêtes, questionnaires.
\emph{Formule type :} « La démarche a reposé sur... » / « Les investigations ont combiné : (1)... (2)... (3)... »

\textbf{Temps 4 — Annoncer le plan (la « carte »).}
Énumérez les parties dans l'ordre exact du corps du rapport. Utilisez des verbes d'action : présenter, analyser, évaluer, proposer.
\emph{Formule type :} « Ce rapport présente successivement : 1) ... 2) ... 3) ... »
\end{methode}

\begin{popfigure}
\begin{tikzpicture}[
    box/.style={draw=popblue, thick, rounded corners, fill=popblueL, minimum width=7.5cm, minimum height=1.1cm, align=center, font=\small},
    arrow/.style={-Stealth, thick, popdark}
]
% Titre introduction
\node[font=\bfseries, popdark, align=center] at (0,3.6) {INTRODUCTION : entonnoir convergent (général $\to$ spécifique)};
% Cadre entonnoir introduction (trapèzes emboîtés)
\draw[draw=poporange, thick, fill=poporangeL, rounded corners] (-5,3.0) rectangle (5,-3.4);
\node[box, align=center] (c1) at (0,2.2){\textbf{1. Contexte administratif}\\(décision, date, qualité)};
\node[box, align=center] (c2) at (0,0.6){\textbf{2. Objet et limites}\\(sujet précis, périmètre)};
\node[box, align=center] (c3) at (0,-1.0){\textbf{3. Méthode}\\(sources, outils, enquête)};
\node[box, align=center] (c4) at (0,-2.6){\textbf{4. Plan annoncé}\\(parties 1, 2, 3... dans l'ordre)};
\draw[arrow] (c1) -- (c2);
\draw[arrow] (c2) -- (c3);
\draw[arrow] (c3) -- (c4);
\node[align=center, font=\footnotesize, popmuted] at (4.6,2.2) {général};
\node[align=center, font=\footnotesize, popmuted] at (4.6,-2.6) {spécifique};

% Titre conclusion
\node[font=\bfseries, popdark, align=center] at (0,-4.2) {CONCLUSION : entonnoir divergent (constats $\to$ recommandations)};
% Cadre entonnoir conclusion
\draw[draw=popgreen, thick, fill=popgreenL, rounded corners] (-5,-4.8) rectangle (5,-8.2);
\node[box, fill=white, draw=popgreen, align=center] (c5) at (0,-5.6){\textbf{1. Bilan synthétique}\\(constats majeurs, causes racines)};
\node[box, fill=white, draw=popgreen, align=center] (c6) at (0,-7.2){\textbf{2. Recommandations}\\(actions priorisées, datées, responsables)};
\draw[arrow, popgreen] (c5) -- (c6);
\node[align=center, font=\footnotesize, popmuted] at (4.6,-5.6) {constats};
\node[align=center, font=\footnotesize, popmuted] at (4.6,-7.2) {décision};
\end{tikzpicture}
\legende{Schéma de la structure en entonnoir : l'introduction resserre le cadre (contexte $\to$ plan), la conclusion ouvre sur l'action (bilan $\to$ recommandations).}
\end{popfigure}

\courssub{La conclusion du rapport : bilan et recommandations (entonnoir inversé)}

La conclusion ne résume pas ; elle \textbf{synthétise} pour \textbf{décider}. Elle part des constats majeurs (issus du corps) pour aboutir à des recommandations opérationnelles. C'est un entonnoir \textbf{inversé} : du spécifique (résultats) vers l'action concrète (décision).

\begin{methode}[Rédiger une conclusion de rapport en deux temps]
\textbf{Temps 1 — Établir le bilan synthétique (sans répétition du corps).}
Reprenez les \textbf{constats majeurs} et les \textbf{causes racines} identifiées dans le corps. Regroupez-les par thème. Évitez le catalogue ; visez la hiérarchie : « L'origine principale est... » / « Trois facteurs critiques ont conjugué leurs effets : ... »
\emph{Formule type :} « Au terme de cette étude, il apparaît que... » / « L'analyse met en évidence... » / « Le constat principal est... »

\textbf{Temps 2 — Formuler les recommandations (ouvertes sur la décision).}
Chaque recommandation doit être \textbf{SMART} : Spécifique, Mesurable, Atteignable, Réaliste, Temporellement définie. Indiquez : \textbf{Quoi} (action), \textbf{Qui} (responsable), \textbf{Quand} (échéance), \textbf{Comment} (moyens/référentiel). Classez par priorité (critique / majeure / mineure).
\emph{Formule type :} « Nous recommandons en conséquence : (1)... avant le [date], sous la responsabilité de [service] ; (2)... » / « Il est proposé à la direction de valider le plan d'actions ci-après... »
\end{methode}

\begin{attention}[Pièges fatals à l'examen et en situation professionnelle]
\begin{itemize}
\item \textbf{Plan annoncé $\neq$ plan réalisé :} l'introduction promet un plan en 3 parties ; le corps en développe 4 (ou dans un ordre différent). \emph{Sanction :} incohérence structurelle, perte de crédibilité, lourde pénalité au Bac.
\item \textbf{Information nouvelle en conclusion :} un chiffre, un fait, un témoin, une cause non traités dans le corps. \emph{Règle :} la conclusion ne fait que \textbf{relire} le corps sous un angle décisionnel.
\item \textbf{Excuses ou opinions non étayées :} « Je n'ai pas eu le temps de... », « Je pense que c'est bien », « C'était difficile ». \emph{Interdit :} le rapport est un acte professionnel objectif.
\item \textbf{Recommandations vagues :} « Améliorer la maintenance », « Former le personnel ». \emph{Exigible :} « Remplacer le pressostat C2 (réf. X) avant le 31/01/2024 par le service maintenance. »
\item \textbf{Confusion introduction / corps :} développer l'analyse dans l'introduction (ex. : détailler la méthode des 5M). L'introduction \textbf{annonce}, le corps \textbf{déploie}.
\end{itemize}
\end{attention}

\courssub{Atelier d'écriture : le rapport de stage de Marlyse (entreprise de transformation de manioc, Bafoussam)}

\begin{savaistu}[Au Baccalauréat technique]
L'introduction et la conclusion pèsent lourd dans l'évaluation de la production écrite fonctionnelle : cohérence du cadre, exhaustivité des 4 temps de l'introduction, pertinence du bilan, opérationnalité des recommandations, qualité de la langue. Les soigner, c'est sécuriser une part décisive de la note.
\end{savaistu}

\begin{experience}[Situation d'accroche — Marlyse, stagiaire en agro-alimentaire]
\textbf{Contexte :} Marlyse, élève de Terminale technique (spécialité agro-alimentaire), a effectué un stage de 6 semaines (03/06 -- 12/07/2024) à l'entreprise \emph{ManiocPlus} (Bafoussam), spécialisée dans la farine de manioc (foufou, gari, tapioca). Elle a visité 3 ateliers : \textbf{Réception/Lavage}, \textbf{Broyage/Pressage}, \textbf{Séchage/Emballage}. La directrice lui a remis une lettre de mission (n° 045/MP/2024 du 20/05/2024) : « Étudier les pertes de matière aux étapes critiques et proposer des pistes de réduction. »
\textbf{Consigne :} rédigez l'introduction et la conclusion de son rapport de stage, en appliquant strictement les deux méthodes.
\end{experience}

\begin{exoresolu}[Introduction et conclusion rédigées du rapport de stage de Marlyse]
\textbf{INTRODUCTION}
Conformément à la lettre de mission n° 045/MP/2024 du 20 mai 2024, signée par Mme Ngo Bell, Directrice Générale de l'entreprise \emph{ManiocPlus} (Bafoussam), il m'a été demandé, en ma qualité de stagiaire de Terminale technique au Lycée Technique de Bafoussam, d'étudier les pertes de matière aux étapes critiques de la ligne de production de farine de manioc et de formuler des propositions de réduction.
La mission, réalisée sur une période de six semaines du 3 juin au 12 juillet 2024, s'est limitée aux trois ateliers de la chaîne principale : Réception/Lavage (A1), Broyage/Pressage (A2), Séchage/Emballage (A3). Les pertes liées au stockage des racines brutes (silos) et à la distribution des produits finis sont exclues du périmètre.
La démarche a reposé sur trois leviers complémentaires : (1) la mesure quotidienne des rendements massiques (entrées/sorties) sur 20 journées de production, selon le protocole interne QP-PROD-12 ; (2) l'observation directe des opérations et des entretiens semi-directifs avec 8 opérateurs et 2 chefs d'atelier ; (3) l'exploitation des registres de rebuts et des fiches de non-conformité du système qualité (période janvier--mai 2024).
Ce rapport présente successivement : 1) la description du processus et l'identification des points critiques ; 2) la quantification des pertes par atelier et par cause (tri, broyage, séchage, emballage) ; 3) l'analyse des causes racines (diagramme d'Ishikawa) ; 4) les propositions techniques et organisationnelles de réduction des pertes, chiffrées et planifiées.

\textbf{CONCLUSION}
Au terme de cette étude, il apparaît que les pertes massiques totales s'élèvent à \textbf{12,4\,\%} de la matière entrante (racines brutes), soit \textbf{3,7 tonnes} sur la période d'observation, concentrées à 78\,\% sur deux postes : le \textbf{triage manuel} en atelier A1 (4,2\,\% de pertes : épluchage excessif, élimination de racines saines par manque de calibration visuelle) et le \textbf{séchage solaire} en atelier A3 (5,5\,\% de pertes : sur-séchage, contamination par les poussières, intempéries non anticipées). L'analyse d'Ishikawa révèle que l'absence de calibreur mécanique (A1) et l'absence de sécheur mécanique piloté par hygrométrie (A3) sont les causes techniques racines ; le turn-over des opérateurs non formés (A1) et l'absence de procédure météo (A3) en sont les causes organisationnelles.
Nous recommandons en conséquence, par ordre de criticité décroissante : (1) \textbf{acquisition d'un calibreur mécanique d'occasion (budget 1,2 million de FCFA)} et formation de 2 opérateurs dédiés (atelier A1) \textbf{avant la campagne 2025 (délai : 31/01/2025)} — gain estimé : $-2,8$ points de pertes ; (2) \textbf{installation d'un sécheur à air chaud de 500 kg/h (budget 3,5 millions de FCFA)} avec sonde hygrométrique et alarme (atelier A3) \textbf{avant le 30/06/2025} — gain estimé : $-4,0$ points de pertes ; (3) \textbf{rédaction et affichage de la fiche OP-METEO-03} (arrêt du séchage si humidité $>$ 70\,\% ou pluie annoncée) et formation de l'équipe A3 \textbf{sous 30 jours} — gain estimé : $-0,8$ point de pertes. La direction validera le plan d'investissement (total 4,7 millions de FCFA) lors du comité de pilotage du 15/09/2024.
\end{exoresolu}

\begin{aretenir}[Check-list de validation (à vérifier avant de rendre la copie)]
\begin{itemize}
\item Introduction : contexte (acte + date + qualité) présent.
\item Introduction : objet précis + limites explicites.
\item Introduction : méthode (sources, outils, effectifs, normes).
\item Introduction : plan énuméré \textbf{dans l'ordre exact du corps}.
\item Conclusion : bilan synthétique (chiffres clés, causes racines, pas de détail).
\item Conclusion : recommandations SMART (quoi, qui, quand, combien, priorité).
\item Conclusion : zéro information nouvelle, zéro excuse, zéro « je pense ».
\item Cohérence lexicale : termes spécialisés définis et employés correctement (section 3).
\end{itemize}
\end{aretenir}

\courssub{Exercices d'application}

\begin{exercice}[Exercice 1 — Diagnostic d'introduction]
Voici l'introduction d'un rapport de stage en maintenance industrielle. Identifiez les \textbf{4 temps} (contexte, objet/limites, méthode, plan) en les numérotant de 1 à 4. Signalez tout élément manquant ou mal placé.
\begin{quote}
« Je m'appelle Kamga Paul, élève en Terminale F3. J'ai fait mon stage à la SODECAO du 1er au 30 juin. C'était bien. On m'a demandé de regarder les pannes des convoyeurs. J'ai vu beaucoup de pannes. J'ai parlé aux mécaniciens. J'ai lu les carnets de maintenance. Dans ce rapport, je vais vous dire ce que j'ai vu. D'abord les pannes mécaniques, ensuite les pannes électriques, enfin mes idées pour réparer. »
\end{quote}
\end{exercice}

\begin{exercice}[Exercice 2 — Réécriture de conclusion]
La conclusion ci-dessous commet \textbf{trois erreurs graves} (voir boîte \textbf{Attention}). Repérez-les, expliquez-les, puis réécrivez la conclusion correctement à partir des données du corps de rapport (fournies en italique).
\begin{quote}
\textit{Données du corps :} pannes convoyeurs : 68\,\% mécaniques (usure des galets, désalignement), 22\,\% électriques (capteurs, variateurs), 10\,\% autres. Causes racines : absence de plan de graissage (galets), procédure d'alignement non respectée. Coût des arrêts : 4,2 millions de FCFA/mois. Recommandations du corps : (1) plan de graissage hebdomadaire + formation, (2) procédure d'alignement au laser + achat d'un appareil (1,5 million de FCFA), (3) suivi des indicateurs MTTR/MTBF.
\textbf{Conclusion de l'élève :} « Pour finir, je pense que les convoyeurs sont vieux. Il faut les changer. J'ai appris beaucoup de choses. Le chef mécanicien était sympa. Je recommande d'acheter un laser pour l'alignement et de faire un planning de graissage. Aussi, il y a eu une panne électrique hier qui n'est pas dans le rapport. Merci de lire mon rapport. »
\end{quote}
\end{exercice}

\begin{exercice}[Exercice 3 — Production complète (sujet type Bac)]
\textbf{Sujet :} vous êtes stagiaire en Terminale Génie Civil à l'entreprise \emph{BTP-EXCEL} (Yaoundé). Le Directeur des Travaux vous confie, par lettre de mission n° 112/DT/2024 du 05/01/2024, une mission de 4 semaines (08/01 -- 02/02/2024) : « Étudier les fissures apparues sur le voile du bâtiment R+3 (chantier \emph{Émeraude}, quartier Mvan) et proposer des solutions de réparation. » Vous avez visité le chantier 12 fois, réalisé 30 mesures de largeur de fissures (fissuromètre), consulté les plans d'exécution (ferraillage, béton), interviewé le chef de chantier et le laboratoire béton (essais de compression à 28 jours : 28 MPa au lieu de 35 MPa). Le corps de votre rapport est déjà rédigé en 3 parties : 1) description du désordre et localisation ; 2) analyse des causes (retrait plastique, sous-dosage en ciment, reprise de charge précoce) ; 3) solutions de réparation (injection de résine, armatures additionnelles, surveillance).
\textbf{Consigne :} rédigez \textbf{uniquement l'introduction et la conclusion} de ce rapport. Respectez scrupuleusement les deux méthodes (4 temps / 2 temps), le lexique spécialisé (fissuromètre, voile, MPa, etc.) et les formules types.
\end{exercice}

\begin{corrige}[Exercice 1]
\textbf{1. Contexte :} « On m'a demandé de regarder les pannes des convoyeurs. » $\to$ \textbf{Manquant :} l'acte officiel (numéro, date, signataire) et la qualité précise de l'élève.
\textbf{2. Objet/Limites :} « les pannes des convoyeurs » $\to$ \textbf{Imprécis :} quel convoyeur ? quelle période ? quelles pannes sont exclues ?
\textbf{3. Méthode :} « J'ai parlé aux mécaniciens. J'ai lu les carnets de maintenance. » $\to$ \textbf{Incomplet :} pas de nombre d'entretiens, pas de période d'observation, pas de référentiel.
\textbf{4. Plan :} « D'abord les pannes mécaniques, ensuite les pannes électriques, enfin mes idées pour réparer. » $\to$ \textbf{Acceptable} dans l'ordre, mais « mes idées » doit être remplacé par « propositions correctives priorisées ».
\textbf{Autres défauts :} registre familier (« C'était bien », « Je vais vous dire »), subjectivité, absence des formules types.
\end{corrige}

\begin{corrige}[Exercice 2]
\textbf{Erreurs :} (1) \textbf{opinion personnelle non étayée} (« je pense que les convoyeurs sont vieux », « il faut les changer ») ; (2) \textbf{information nouvelle} (« une panne électrique hier qui n'est pas dans le rapport ») ; (3) \textbf{propos hors sujet et excuses déguisées} (« J'ai appris beaucoup », « Le chef était sympa », « Merci de lire »).
\textbf{Conclusion corrigée :}
« Au terme de cette étude, il apparaît que les arrêts des convoyeurs (coût : 4,2 millions de FCFA/mois) sont à 68\,\% d'origine mécanique (usure des galets par absence de plan de graissage, désalignement par non-respect de la procédure) et à 22\,\% d'origine électrique (capteurs, variateurs).
Nous recommandons en conséquence : (1) la mise en place immédiate d'un plan de graissage hebdomadaire (fiche OP-MAIN-09) et la formation de l'équipe maintenance (échéance : 15/02/2024) ; (2) l'acquisition d'un appareil d'alignement laser (budget : 1,5 million de FCFA) et la rédaction de la procédure d'alignement (échéance : 31/03/2024) ; (3) le déploiement des indicateurs MTTR/MTBF dans la GMAO pour un suivi mensuel (échéance : 30/04/2024). La direction validera le budget au comité du 20/02/2024. »
\end{corrige}

\begin{corrige}[Exercice 3 — Proposition de rédaction]
\textbf{INTRODUCTION}
Conformément à la lettre de mission n° 112/DT/2024 du 5 janvier 2024, émanant du Directeur des Travaux de l'entreprise \emph{BTP-EXCEL}, il m'a été demandé, en ma qualité de stagiaire de Terminale Génie Civil au Lycée Technique de Yaoundé, d'étudier les fissures affectant le voile du bâtiment R+3 du chantier \emph{Émeraude} (quartier Mvan, Yaoundé) et de proposer des solutions de réparation adaptées.
La mission, conduite sur quatre semaines du 8 janvier au 2 février 2024, s'est limitée à l'analyse du désordre sur le voile principal (du rez-de-chaussée au R+3) ; les fondations, les planchers et les murs de façade sont exclus du périmètre.
La démarche a reposé sur quatre leviers : (1) douze visites de terrain avec relevés fissurométriques (30 mesures de largeur, à fréquence hebdomadaire) ; (2) l'exploitation des plans d'exécution (ferraillage du voile, béton de classe C35/45) et du rapport du laboratoire béton (essais de compression à 28 jours : 28 MPa) ; (3) des entretiens structurés avec le chef de chantier (reprise de charge au jour 14) et le responsable du laboratoire ; (4) l'analyse documentaire des fiches de coulage et des conditions climatiques (température, hygrométrie) de la période.
Ce rapport présente successivement : 1) la description et la localisation précise des fissures (cartographie, largeur, orientation) ; 2) l'analyse des causes racines (retrait plastique, sous-dosage en ciment avéré, reprise de charge précoce) ; 3) les solutions techniques de réparation (injection de résine époxy, pose d'armatures additionnelles, mise en surveillance fissurométrique continue).

\textbf{CONCLUSION}
Au terme de cette étude, il apparaît que le réseau de fissures traversantes (largeur moyenne 0,8 mm, maximale 1,5 mm) sur le voile du R+3 résulte de la conjugaison de trois causes : un \textbf{retrait plastique} favorisé par un temps chaud et venté (surface non protégée après coulage), un \textbf{sous-dosage en ciment effectif} (28 MPa au lieu de 35 MPa, soit environ 20\,\% de résistance en moins) et une \textbf{reprise de charge prématurée} au jour 14 (au lieu du jour 21 au minimum pour un béton C35/45). L'intégrité structurelle n'est pas compromise à court terme (ferraillage conforme), mais la durabilité (carbonatation, corrosion des armatures) est menacée.
Nous recommandons en conséquence, par priorité : (1) l'\textbf{injection d'une résine époxy de basse viscosité} sur l'ensemble des fissures de largeur supérieure à 0,3 mm, réalisée par une entreprise qualifiée \textbf{avant le 15/03/2024} (budget estimé : 850 000 FCFA) ; (2) la \textbf{pose d'armatures additionnelles} (HA12, double peau, ancrage chimique) aux zones critiques (pieds de voile, reprises de charge) \textbf{avant le 30/04/2024} (budget : 1,2 million de FCFA) ; (3) la \textbf{mise en place d'un suivi fissurométrique mensuel} (30 points, 12 mois) avec seuils d'alerte (évolution supérieure à 0,2 mm/mois) \textbf{dès la réception des travaux}. Le Directeur des Travaux validera le dossier technique et le budget (total : 2,05 millions de FCFA) en comité de chantier du 10/03/2024.
\end{corrige}

\coursec{Rédiger et améliorer le corps du rapport}

Après avoir maîtrisé l'écriture de l'introduction qui pose le cadre et de la conclusion qui synthétise et ouvre, nous abordons maintenant le cœur même du rapport : \textbf{le corps}. C'est ici que se joue la crédibilité de votre travail. Un rapport technique ne se contente pas d'aligner des faits ; il les \emph{structure}, les \emph{analyse} et les \emph{relie} pour répondre à la problématique annoncée. Comme le disait l'ingénieur et pédagogue Henri Fayol : « \emph{Prévoir, organiser, commander, coordonner, contrôler} » — le corps du rapport est l'espace où vous démontrez votre capacité à \textbf{organiser} l'information et à \textbf{coordonner} les preuves.

\courssub{Le paragraphe de rapport : une architecture « idée – preuve – commentaire »}

\begin{definition}[Paragraphe fonctionnel du rapport]
Dans un rapport technique, un paragraphe ne se définit pas seulement par un retour à la ligne. C'est une \textbf{unité de sens close} qui développe \textbf{une seule idée principale} (souvent annoncée par une phrase d'amorce) et l'étaye par des \textbf{preuves concrètes} (données chiffrées, observations, citations) suivies d'un \textbf{commentaire analytique} qui en explicite la portée.
\end{definition}

\begin{methode}[Construire un paragraphe de rapport efficace]
\textbf{Étape 1 : Formuler l'idée directrice (phrase d'amorce).} Elle annonce le sujet du paragraphe et son lien avec le plan. Exemple : « L'approvisionnement en manioc constitue le premier goulot d'étranglement de la chaîne de production. »
\textbf{Étape 2 : Apporter la preuve factuelle (donnée brute).} Chiffre, mesure, observation datée, citation sourcée. Exemple : « Sur la période du 10 au 14 mars, l'atelier n'a reçu que 850 kg de manioc contre 1 200 kg attendus, soit un déficit de 29 \% (cf. tableau des réceptions, annexe 3). »
\textbf{Étape 3 : Commenter et analyser (valeur ajoutée).} Expliquer \emph{ce que cela signifie} pour l'activité. Exemple : « Ce déficit récurrent oblige à réduire la cadence des presses, augmentant d'autant le coût unitaire de transformation. »
\textbf{Étape 4 : Assurer la transition.} Un connecteur logique ou une reprise nominale pour glisser vers l'idée suivante. Exemple : « Par conséquent, la gestion des stocks amont apparaît comme un levier prioritaire d'optimisation. »
\end{methode}

\begin{popfigure}
\begin{tikzpicture}[
    node distance=1.2cm and 1.5cm,
    box/.style={draw=popblue, thick, rounded corners, fill=popblueL, text width=3.2cm, align=center, minimum height=1.8cm},
    arrow/.style={-Stealth, thick, popdark},
    labelbox/.style={font=\small\bfseries, popdark}
]
\node[box, labelbox, align=center] (idee){Idée principale\\(Phrase d'amorce)};
\node[box, labelbox, right=of idee, align=center] (preuve){Preuve factuelle\\(Chiffre, tableau,\\citation, observation)};
\node[box, labelbox, right=of preuve, align=center] (comment){Commentaire analytique\\(Interprétation,\\conséquence, lien probl.)};
\node[box, labelbox, right=of comment, align=center] (trans){Transition\\(Connecteur / Reprise)};
\draw[arrow] (idee) -- (preuve) node[midway, above, font=\tiny] {s'appuie sur};
\draw[arrow] (preuve) -- (comment) node[midway, above, font=\tiny] {s'analyse par};
\draw[arrow] (comment) -- (trans) node[midway, above, font=\tiny] {entraîne};
\end{tikzpicture}
\legende{Schéma de la structure type d'un paragraphe de rapport technique : l'enchaînement logique de l'idée à la transition.}
\end{popfigure}

\begin{exemplebox}[Paragraphe rédigé à partir de données brutes]
\textbf{Données brutes (notes de stage de Marlyse) :}
\begin{itemize}
    \item Production journalière d'attiéké : Jour 1 : 120 kg, Jour 2 : 115 kg, Jour 3 : 90 kg, Jour 4 : 118 kg, Jour 5 : 122 kg.
    \item Effectif : 8 employés (2 absents le Jour 3).
    \item Équipement : 2 presses hydrauliques (1 en panne le Jour 3).
\end{itemize}

\textbf{Paragraphe rédigé :}
\emph{« La stabilité de la production d'attiéké est fortement tributaire de la disponibilité de l'outillage principal. Sur les cinq jours d'observation, la production moyenne s'établit à 113 kg/jour. Cependant, le Jour 3 enregistre une chute brutale à 90 kg (-20 \% par rapport à la moyenne). Cette contre-performance coïncide avec deux événements concomitants : l'absence de deux opérateurs sur huit (25 \% de l'effectif) et, surtout, la panne de l'une des deux presses hydrauliques. L'analyse croisée de ces facteurs montre que la panne matérielle est le facteur limitant majeur : même avec un effectif complet, une seule presse ne peut traiter le volume de manioc requis. Il en résulte un goulot d'étranglement au stade du pressage, étape critique pour la déshydratation du couscous de manioc. \textbf{En outre}, cette vulnérabilité structurelle (absence de presse de secours) expose l'unité à des pertes sèches en cas de panne prolongée. »}

\textcolor{popgreen}{\textbf{Analyse :}} Notez l'enchaînement : \textbf{Idée} (stabilité tributaire de l'outillage) $\rightarrow$ \textbf{Chiffres} (moyenne, Jour 3, -20 \%) $\rightarrow$ \textbf{Causes factuelles} (absences, panne) $\rightarrow$ \textbf{Analyse technique} (goulot d'étranglement au pressage) $\rightarrow$ \textbf{Conséquence stratégique} (vulnérabilité) $\rightarrow$ \textbf{Transition} (\og En outre \fg{} vers les solutions).
\end{exemplebox}

\courssub{L'articulation du corps : les connecteurs, charnières de la démonstration}

Un rapport sans connecteurs logiques n'est qu'une \emph{liste de courses}. Les connecteurs (ou opérateurs argumentatifs) signalent au lecteur la \textbf{nature du lien} entre deux idées. Leur choix n'est pas anodin : il reflète votre raisonnement.

\begin{propriete}[Les quatre familles de connecteurs indispensables au rapport]
\begin{center}
\begin{tabularx}{\linewidth}{|l|X|X|}\hline
\textbf{Fonction logique} & \textbf{Connecteurs courants (registre formel)} & \textbf{Effet produit dans le rapport} \\ \hline
\rowcolor{popblueL}
\textbf{Addition / Énumération} & \emph{D'abord, en premier lieu ; ensuite, en second lieu ; de plus, en outre, par ailleurs ; enfin, pour terminer.} & Structure le plan, montre l'exhaustivité de l'analyse. \\ \hline
\textbf{Opposition / Concession} & \emph{Cependant, toutefois, pourtant ; si \dots\ toutefois ; en revanche ; malgré \dots ; alors que.} & Signale un écart, une limite, un résultat inattendu (ex. : prévisions vs réel). \\ \hline
\textbf{Cause / Explication} & \emph{Car, en effet ; parce que ; du fait de ; en raison de ; grâce à ; eu égard à.} & Justifie un fait, ancre l'analyse technique (ex. : \og La baisse du rendement est due \emph{à l'usure des lames} \fg{}). \\ \hline
\textbf{Conséquence / Conclusion partielle} & \emph{Par conséquent, c'est pourquoi ; dès lors ; il en résulte que ; par suite ; ainsi.} & Tire la leçon technique ou managériale du fait exposé. \\ \hline
\end{tabularx}
\end{center}
\end{propriete}

\begin{attention}[Piège majeur : la « liste à puces » déguisée en paragraphe]
\textbf{Erreur fréquente :} Empiler des phrases factuelles sans connecteurs ni commentaire.
\begin{quote}
\emph{« Le Jour 1, on a fait 120 kg. Le Jour 2, 115 kg. Le Jour 3, 90 kg. La presse a cassé. Deux ouvriers manquaient. C'est un problème. »}
\end{quote}
\textbf{Pourquoi c'est un échec :} Aucune hiérarchie des idées, pas d'analyse, pas de lien avec la problématique. Le lecteur ne sait pas \emph{quoi retenir}.
\textbf{Correction :} Regrouper en paragraphes thématiques (ex. \og Évolution de la production \fg{}, \og Analyse des écarts \fg{}) et utiliser les connecteurs pour guider la lecture (voir l'exemple de Marlyse ci-dessus).
\end{attention}

\courssub{Intégrer les données techniques : du brut au signifiant}

Le rapport technique manipule des \textbf{données hétérogènes} : tableaux de production, relevés de maintenance, photographies d'installations, comptes rendus d'entretiens. Le rédacteur ne les \emph{pose} pas ; il les \emph{fait parler}.

\begin{methode}[Intégrer un élément de preuve (tableau, chiffre, citation, image)]
\textbf{1. Introduire la source.} « Le tableau 2 (annexe) récapitule les réceptions quotidiennes... », « Lors de l'entretien du 12 mars, M. K., chef d'atelier, a indiqué... », « La photographie 3 (p. 12) montre l'état du convoyeur... »
\textbf{2. Présenter la donnée clé (extraction).} Ne pas dire \og voir tableau \fg{}. Dire : \og Ce tableau révèle que le taux d'humidité moyen dépasse 14 \% sur 4 jours sur 5. \fg{}
\textbf{3. Commenter (interprétation technique).} \og Ce taux excède la norme de stockage (12 \%), ce qui explique les pertes par fermentation observées en fin de semaine. \fg{}
\textbf{4. Refermer par le lien avec la problématique.} \og La maîtrise de l'humidité à la réception est donc le premier levier de réduction des pertes. \fg{}
\end{methode}

\begin{exemplebox}[Intégration d'un entretien et d'un tableau]
\textbf{Source :} Entretien avec M. Essomba, responsable maintenance (15 mars). Tableau des arrêts machine (Semaine 11).
\textbf{Rédaction :}
\emph{« La maintenance curative prédomine encore largement sur la préventive. \textbf{Selon le responsable maintenance}, M. Essomba, \og 80 \% de nos interventions sont des dépannages d'urgence \fg{} (entretien du 15/03). Le tableau des arrêts machine (Semaine 11) confirme ce diagnostic : sur 14 arrêts recensés, 11 sont non planifiés, représentant 18 heures de production perdue. \textbf{En raison de} l'absence de planning préventif, les pannes surviennent de manière aléatoire, perturbant l'ordonnancement. \textbf{Par conséquent}, la mise en place d'une GMAO (Gestion de Maintenance Assistée par Ordinateur) s'impose pour basculer vers une maintenance planifiée. »}
\end{exemplebox}

\begin{popfigure}
\begin{tikzpicture}
\begin{axis}[
    ybar,
    bar width=12pt,
    width=\linewidth,
    height=7cm,
    ylabel={Production (kg)},
    xlabel={Jours de stage},
    symbolic x coords={J1, J2, J3, J4, J5},
    xtick=data,
    ymin=0, ymax=140,
    ymajorgrids=true,
    grid style={popline},
    tick label style={font=\small},
    label style={font=\small},
    nodes near coords,
    nodes near coords align={vertical},
    every node near coord/.append style={font=\tiny, popdark},
    every axis plot/.append style={fill=popblue},
    title style={font=\small\bfseries, popdark},
    title={Production journalière d'attiéké (Stage Marlyse - Semaine 11)}
]
\addplot coordinates {(J1,120) (J2,115) (J3,90) (J4,118) (J5,122)};
\end{axis}
\end{tikzpicture}
\legende{Données brutes de production journalière (kg) à exploiter pour le corps du rapport. L'écart du J3 est l'élément déclencheur de l'analyse.}
\end{popfigure}

\courssub{La cohérence textuelle : reprises, progressions et temps verbaux}

Au-delà des connecteurs, la fluidité du rapport repose sur des \textbf{mécanismes linguistiques invisibles} mais puissants.

\begin{propriete}[Cohérence référentielle et temporelle]
\textbf{1. Reprises nominales (chaînes de référence).} Éviter la répétition lourde du nom propre. Utiliser : pronoms (\emph{il, elle, ils}), déictiques (\emph{ce dernier, ledit, celui-ci}), synonymes techniques (\emph{l'unité $\rightarrow$ l'atelier $\rightarrow$ la structure $\rightarrow$ l'entité}).
\textbf{2. Progressions thématiques.}
\begin{itemize}
    \item \emph{Progression linéaire (en chaînage) :} Le rhème (fin de phrase) devient le thème (début de la phrase suivante). Idéal pour raconter un processus chronologique.
    \item \emph{Progression à thème constant :} Même sujet grammatical pour une série de constats. Idéal pour dresser un état des lieux.
    \item \emph{Progression à rhème constant :} Différents angles d'attaque sur une même conclusion. Idéal pour l'argumentation.
\end{itemize}
\textbf{3. Temps verbaux cohérents.}
\begin{itemize}
    \item \textbf{Passé composé / Imparfait :} Récit de l'activité observée, faits accomplis pendant le stage (\og \emph{La presse a lâché} \fg{}, \og \emph{les ouvriers chargeaient} \fg{}).
    \item \textbf{Présent de vérité générale :} Description de l'organisation, des procédures, des machines (\og \emph{L'unité transforme} \fg{}, \og \emph{Le processus comprend} \fg{}).
    \item \textbf{Futur / Conditionnel :} Recommandations, prévisions (\og \emph{Il conviendra de} \fg{}, \og \emph{Cela permettrait de} \fg{}).
\end{itemize}
\textbf{Règle d'or :} Ne pas mélanger les temps à l'intérieur d'un même paragraphe narratif sans raison logique (ex. : passage du récit à l'analyse générale = changement de temps justifié).
\end{propriete}

\begin{exoresolu}[Analyse de cohérence : repérer les progressions et les temps]
\textbf{Extrait à corriger :}
\emph{« L'atelier produit de l'attiéké. \textbf{La machine} broie le manioc. \textbf{Le broyage} est une étape clé. \textbf{Il} faut surveiller l'humidité. \textbf{Le taux} était trop haut la semaine dernière. \textbf{Cela} a pourri la pâte. \textbf{On} doit changer le séchoir. »}

\textbf{Diagnostic :}
\begin{itemize}
    \item Progression : le thème saute sans logique (\og L'atelier \fg{} $\rightarrow$ \og La machine \fg{} $\rightarrow$ \og Le broyage \fg{} $\rightarrow$ \og Il \fg{} [impersonnel] $\rightarrow$ \og Le taux \fg{} $\rightarrow$ \og Cela \fg{} $\rightarrow$ \og On \fg{}). \textbf{Incohérence référentielle.}
    \item Temps : Présent (vérité) $\rightarrow$ Présent $\rightarrow$ Présent $\rightarrow$ Présent (obligation) $\rightarrow$ Imparfait (récit) $\rightarrow$ Passé composé (conséquence) $\rightarrow$ Présent (recommandation). \textbf{Mélange non justifié.}
\end{itemize}

\textbf{Version corrigée (progression linéaire + temps justifiés) :}
\emph{« \textbf{L'atelier} produit de l'attiéké selon un processus en trois temps. \textbf{La première étape, le broyage}, transforme le manioc frais en pâte. \textbf{Cette opération} conditionne la qualité finale : \textbf{elle} nécessite un contrôle strict de l'humidité. \textbf{Or, la semaine dernière,} le taux d'humidité \textbf{a dépassé} les normes, \textbf{ce qui a entraîné} la fermentation de la pâte. \textbf{Aussi,} le remplacement du séchoir \textbf{s'impose} comme une priorité. »}
\tcblower
\textbf{Analyse de la correction :}
\begin{itemize}
    \item \textbf{Chaînage :} L'atelier $\rightarrow$ première étape $\rightarrow$ Cette opération $\rightarrow$ elle $\rightarrow$ Or $\rightarrow$ ce qui $\rightarrow$ Aussi.
    \item \textbf{Temps :} Présent (description de la procédure) $\rightarrow$ Passé composé (récit d'un événement précis et daté) $\rightarrow$ Présent (recommandation générale). Chaque changement de temps correspond à un changement de statut du propos : description, récit, recommandation.
\end{itemize}
\end{exoresolu}

\courssub{L'amélioration du texte : la relecture en trois passes professionnelles}

Rédiger, c'est réécrire. Un rapport technique de qualité exige une \textbf{relecture systématique en trois passes distinctes}, chacune avec un objectif unique. Ne tentez pas de tout faire en une fois.

\begin{methode}[Relecture en trois passes (méthode « Entonnoir »)]
\textbf{Passe 1 : Contenu et Logique (le \og Quoi \fg{}).} \emph{Lire sans stylo, comme un client.}
\begin{enumerate}
    \item Chaque paragraphe répond-il à une question du plan annoncé en introduction ?
    \item Les chiffres sont-ils exacts, sourcés, unités indiquées (kg, h, FCFA, \%) ?
    \item L'analyse technique est-elle juste ? (Pas d'affirmation gratuite : \og c'est bien \fg{} $\rightarrow$ \og rendement de 87 \% \fg{}).
    \item Les recommandations découlent-elles logiquement des constats ?
\end{enumerate}
\textbf{Action :} Supprimer les paragraphes hors-sujet, déplacer les idées mal placées, ajouter les preuves manquantes.

\textbf{Passe 2 : Organisation et Connecteurs (le \og Comment \fg{}).} \emph{Surligner les connecteurs.}
\begin{enumerate}
    \item Y a-t-il un connecteur logique à chaque changement d'idée ?
    \item Le connecteur est-il le \textbf{bon} ? (Ne pas mettre \og Par conséquent \fg{} pour une simple addition).
    \item Les transitions entre parties/sections sont-elles explicites ?
    \item Les progressions thématiques sont-elles fluides (pas de sauts brutaux de sujet) ?
\end{enumerate}
\textbf{Action :} Insérer/remplacer les connecteurs, réordonner les phrases, ajouter des phrases de transition.

\textbf{Passe 3 : Langue et Code (la \og Forme \fg{}).} \emph{Lire à voix haute (voir \og Le saviez-vous \fg{}).}
\begin{enumerate}
    \item Orthographe lexicale (mots techniques : \emph{pressoir, décantation, goulot, manutention}).
    \item Grammaire : accords (sujet/verbe, participes passés), pronoms relatifs (\emph{dont, auquel}).
    \item Ponctuation : virgules autour des incises, point-virgule pour séparer des items complexes, deux-points pour annoncer une liste ou une explication.
    \item Typographie : espaces avant \og ; : ? ! \fg{}, guillemets français \og \fg{}, abréviations normalisées (\emph{etc., cf., p., fig.}).
\end{enumerate}
\textbf{Action :} Corriger, reformuler les phrases lourdes (plus de 2 lignes), supprimer les redondances (\og \emph{au jour d'aujourd'hui} \fg{} $\rightarrow$ \og \emph{aujourd'hui} \fg{}).
\end{methode}

\begin{savaistu}[La relecture à voix haute : l'outil secret des professionnels]
Des travaux en psychologie cognitive de l'écriture montrent que la \textbf{relecture à voix haute} (ou subvocalisée distinctement) permet de détecter une large part des maladresses syntaxiques et des manques de ponctuation que l'œil seul laisse passer en lecture silencieuse. Pourquoi ? Parce que l'oreille perçoit le \textbf{rythme} de la phrase : une phrase trop longue coupe le souffle, une ponctuation manquante crée une ambiguïté sonore, une répétition sonne faux. \textbf{Astuce :} Enregistrez-vous sur votre téléphone (mémo vocal) en lisant votre rapport, puis réécoutez en suivant le texte. Vous entendrez ce que le lecteur « entend » en vous lisant.
\end{savaistu}

\begin{attention}[Erreurs de forme fréquentes en rapport technique]
\begin{itemize}
    \item \textbf{Redondance :} \og \emph{Au niveau de la production} \fg{} $\rightarrow$ \og \emph{Dans la production} \fg{} ou \og \emph{La production} \fg{}. Éviter aussi : \og \emph{Le fait que... est un fait qui...} \fg{}.
    \item \textbf{Anglicismes syntaxiques :} \og \emph{Le rapport est basé sur...} \fg{} (calque de \emph{based on}) $\rightarrow$ \og \emph{Le rapport repose sur / s'appuie sur...} \fg{}.
    \item \textbf{Verbes « fourre-tout » :} \emph{faire, effectuer, réaliser, mettre en place.} Préférez le verbe technique précis : \og \emph{presser, broyer, calibrer, paramétrer, diagnostiquer, optimiser} \fg{}.
    \item \textbf{Passif abusif :} \og \emph{Il a été constaté que...} \fg{} (lourd, impersonnel). Préférez l'actif : \og \emph{L'observation montre que...} \fg{} ou \og \emph{Nous avons constaté que...} \fg{} (si « nous » désigne l'équipe de stage).
\end{itemize}
\end{attention}

\courssub{Atelier guidé : rédaction et correction croisée du corps du rapport de Marlyse}

\emph{Mise en situation :} Marlyse a effectué son stage de fin d'études à l'unité de transformation « Manioc Plus » (Douala). Elle doit rédiger le corps de son rapport de stage. Voici ses notes brutes structurées par partie du plan.

\begin{experience}[Atelier : du brut au corps final — données de Marlyse]
\textbf{Partie 1 : Organisation de la production (plan annoncé : 1. Approvisionnement ; 2. Processus ; 3. Ressources humaines).}

\textbf{Notes brutes Partie 1 (Approvisionnement) :}
\begin{itemize}
    \item Fournisseurs : 3 coopératives villageoises (A, B, C).
    \item Contrat : 200 kg/jour/coopérative = 600 kg/jour théoriques.
    \item Réel semaine 11 (kg) : lundi 580, mardi 610, mercredi 420, jeudi 590, vendredi 600.
    \item Problème mercredi : camion de la coopérative B en panne.
    \item Qualité : 15 \% de racines impropres (trop petites, pourries) $\rightarrow$ tri manuel de 2 h/jour.
    \item Chef approvisionnement (M. Ngu) : \og On n'a pas de chambre froide, on subit la saison des pluies. \fg{}
\end{itemize}

\textbf{Notes brutes Partie 2 (Processus - Équipements) :}
\begin{itemize}
    \item Étapes : Réception $\rightarrow$ Lavage $\rightarrow$ Épluchage $\rightarrow$ Broyage $\rightarrow$ Pressage $\rightarrow$ Tamisage $\rightarrow$ Cuisson.
    \item Équipements clés : 1 laveuse (bon état), 1 éplucheuse (vétuste, 30 \% de casse), 1 broyeur (bon état), 2 presses hydrauliques (1 neuve, 1 vétuste), 1 four amélioré (bon état).
    \item Pannes semaine 11 : éplucheuse (2 h mardi), presse vétuste (4 h mercredi).
    \item Rendement global observé : 1 kg de manioc $\rightarrow$ 0,18 kg d'attiéké fini (théorique : 0,22 kg).
\end{itemize}

\textbf{Notes brutes Partie 3 (Ressources humaines) :}
\begin{itemize}
    \item Effectif : 12 permanents + 4 journaliers (aux pics d'activité).
    \item Organisation : 3 équipes tournantes (postes de 8 h). Absentéisme : 2 personnes/jour en moyenne.
    \item Compétences : 1 seul technicien de maintenance (M. Essomba). Opérateurs polyvalents mais non certifiés.
    \item Entretien DRH : \og On forme sur le tas, pas de budget formation. \fg{}
\end{itemize}

\textbf{Consignes de l'atelier :}
\begin{enumerate}
    \item \textbf{Rédaction individuelle (30 min) :} Rédiger \textbf{un paragraphe complet} pour \textbf{chacune des trois parties} (soit 3 paragraphes au total) en appliquant la méthode « Idée – Preuve – Commentaire – Transition ». Utiliser le graphique de production (figure ci-dessus) pour la partie Processus si pertinent.
    \item \textbf{Correction croisée par binôme (20 min) :} Échanger votre production. Utiliser la \textbf{grille d'amélioration} ci-dessous. Annoter le texte du pair.
    \item \textbf{Réécriture finale (10 min) :} Intégrer les corrections pertinentes.
\end{enumerate}
\end{experience}

\begin{aretenir}[Grille d'amélioration par les pairs (correction croisée)]
\begin{center}
\begin{tabularx}{\linewidth}{|l|X|X|}\hline
\textbf{Critère} & \textbf{Indicateurs de réussite (cocher si OK)} & \textbf{Commentaire / correction proposée} \\ \hline
\rowcolor{popblueL}
\textbf{Structure du paragraphe} & $\square$ Une idée principale claire (phrase d'amorce) & \\
& $\square$ Preuve(s) chiffrée(s) / citation(s) intégrée(s) & \\
& $\square$ Commentaire technique (pas juste une description) & \\
& $\square$ Transition vers l'idée suivante / la partie suivante & \\ \hline
\textbf{Connecteurs} & $\square$ Connecteurs logiques présents et justes (addition, opposition, cause, conséquence) & \\
& $\square$ Pas de « liste » sans liens & \\ \hline
\textbf{Données} & $\square$ Chiffres exacts, unités, sources citées (tableau, entretien, observation) & \\
& $\square$ Données interprétées (pas brutes) & \\ \hline
\textbf{Cohérence} & $\square$ Reprises nominales fluides (pas de répétitions lourdes) & \\
& $\square$ Progression thématique logique (linéaire / thème constant) & \\
& $\square$ Temps verbaux cohérents (récit vs analyse vs recommandation) & \\ \hline
\textbf{Langue} & $\square$ Vocabulaire technique précis (pas de verbes fourre-tout) & \\
& $\square$ Phrases claires, concises (pas de phrases de plus de 2 lignes) & \\
& $\square$ Orthographe, grammaire, ponctuation, typographie soignées & \\ \hline
\end{tabularx}
\end{center}
\end{aretenir}

\begin{exoresolu}[Exemple de paragraphe attendu pour la Partie 1 (Approvisionnement)]
\emph{« L'approvisionnement en manioc frais, matière première unique, constitue la première vulnérabilité de l'unité. \textbf{D'après le planning contractuel}, les trois coopératives partenaires doivent livrer 600 kg/jour. \textbf{Or,} les relevés de la semaine 11 révèlent une instabilité marquée : si les lundis, mardis, jeudis et vendredis approchent l'objectif (moyenne de 595 kg), \textbf{le mercredi chute à 420 kg} (-30 \%), \textbf{en raison de} la panne du camion de la coopérative B. \textbf{Par ailleurs,} la qualité suit la même irrégularité : un taux moyen de 15 \% de racines impropres impose un tri manuel de 2 heures quotidiennes, mobilisant deux opérateurs. \textbf{Selon le chef d'approvisionnement}, M. Ngu, \og l'absence de chambre froide nous fait subir la saison des pluies \fg{}, ce qui accélère la dégradation des stocks. \textbf{En conséquence}, la sécurisation des flux logistiques (transport, stockage frigorifique) conditionne la régularité de toute la chaîne aval. »}
\tcblower
\textcolor{popgreen}{\textbf{Validation grille :}} Idée (vulnérabilité de l'approvisionnement) $\rightarrow$ Preuves (contractuel vs réel du mercredi, taux de 15 \%, 2 h de tri) $\rightarrow$ Commentaire (impact sur la main-d'œuvre, cause racine : absence de froid) $\rightarrow$ Citation sourcée $\rightarrow$ Conséquence/Transition (conditionne la chaîne aval). Connecteurs : \og Or, Par ailleurs, En raison de, Selon, En conséquence \fg{}. Temps : présent (constat structurel), passé composé implicite dans les relevés datés (semaine précise), présent (citation et généralité). Vérification des calculs : moyenne hors mercredi $= (580 + 610 + 590 + 600)/4 = 595$ kg ; écart du mercredi $= (420 - 600)/600 = -30 \%$.
\end{exoresolu}

\begin{exercice}[Entraînement autonome : le paragraphe « Ressources humaines »]
À partir des notes brutes de la \textbf{Partie 3} (ci-dessus), rédigez \textbf{un paragraphe unique et structuré} qui traite de l'adéquation « Effectifs / Compétences / Organisation ». Appliquez strictement la méthode en 4 étapes et la grille d'amélioration. Visez 12 à 15 lignes.
\end{exercice}

\begin{corrige}[Exercice : paragraphe Ressources humaines]
\emph{« La gestion des ressources humaines révèle un décalage entre l'organisation du travail et la maîtrise technique des équipements. \textbf{L'effectif théorique} de 16 personnes (12 permanents + 4 journaliers) est \textbf{réduit de fait} par un absentéisme moyen de 2 personnes/jour (12,5 \%), fragilisant les équipes tournantes de 8 heures. \textbf{Surtout,} la maintenance, pilier de la disponibilité machine, repose sur \textbf{un unique technicien} (M. Essomba), créant un point de défaillance critique lors des pannes simultanées (éplucheuse + presse le mercredi). \textbf{Les opérateurs, bien que polyvalents par nécessité,} ne disposent d'aucune certification formelle ; \textbf{selon la DRH}, \og on forme sur le tas, pas de budget formation \fg{}. \textbf{Dès lors,} l'investissement dans un plan de formation certifiant et le recrutement d'un second technicien de maintenance apparaissent comme des prérequis non négociables pour réduire les temps d'arrêt. »}

\textbf{Vérification :} absentéisme $= 2/16 = 12,5 \%$ de l'effectif théorique. Structure : idée (décalage organisation/compétences) $\rightarrow$ preuves (16 personnes, 2 absents/jour, 1 technicien, citation DRH) $\rightarrow$ commentaire (point de défaillance critique) $\rightarrow$ transition vers les recommandations (\og Dès lors \fg{}).
\end{corrige}

\courssub{Bilan de la section}

\begin{aretenir}[Ce qu'il faut retenir pour le corps du rapport]
\begin{itemize}
    \item \textbf{1 paragraphe = 1 idée + 1 preuve + 1 commentaire + 1 transition.} Pas de paragraphe « fourre-tout », pas de paragraphe « liste ».
    \item \textbf{Les connecteurs ne sont pas décoratifs :} ils traduisent votre logique (addition, opposition, cause, conséquence). Choisissez-les avec la précision d'un technicien qui choisit son outil.
    \item \textbf{La donnée ne parle pas d'elle-même :} introduisez-la, extrayez l'essentiel, commentez-la techniquement, refermez sur la problématique.
    \item \textbf{Cohérence invisible :} reprises nominales variées, progression thématique maîtrisée, temps verbaux justifiés (passé = récit du stage / présent = vérité technique / futur-conditionnel = recommandations).
    \item \textbf{Relecture en 3 passes :} (1) Sens/Logique $\rightarrow$ (2) Structure/Connecteurs $\rightarrow$ (3) Langue/Code. \textbf{Lisez à voix haute} pour la passe 3.
\end{itemize}
\end{aretenir}

Vous disposez maintenant de toutes les clés pour transformer vos notes de terrain en un corps de rapport rigoureux, analytique et professionnel. La section suivante — et dernière — vous guidera pour assembler l'ensemble (introduction, corps, conclusion, annexes) dans une rédaction finale aboutie, prête pour l'évaluation du Probatoire et du Baccalauréat technique.

\coursec{De la préparation à la rédaction complète du rapport}

Dans la section précédente, nous avons appris à rédiger et à améliorer le \cle{corps du rapport} : paragraphes articulés, connecteurs logiques, données exactes. Mais un bon corps ne suffit pas. Un rapport réussi est le fruit d'une démarche complète qui commence bien avant la rédaction et se termine après elle : la \cle{préparation}, la \cle{rédaction} intégrale et l'\cle{amélioration finale}. C'est cette chaîne complète que nous allons maintenant construire, étape par étape, jusqu'à la version définitive.

\courssub{Le travail de préparation : la moitié du travail}

\begin{definition}[Le travail de préparation]
Le \cle{travail de préparation} est l'ensemble des opérations réalisées avant la rédaction proprement dite : définition de l'objet et du destinataire du rapport, collecte des données, prise de notes, tri et hiérarchisation des informations, élaboration du plan détaillé.
\end{definition}

L'intuition est simple : un maçon ne commence jamais une maison sans plan ni matériaux rassemblés. De même, le rapporteur ne rédige jamais « au fil de la plume ». Détaillons chaque opération.

\textbf{1. Définir l'objet et le destinataire.}\par
Avant toute chose, il faut répondre à deux questions : \emph{de quoi vais-je rendre compte ?} (l'\cle{objet}) et \emph{à qui s'adresse mon rapport ?} (le \cle{destinataire}). Le destinataire détermine le ton, le niveau de détail et le lexique : un rapport adressé au proviseur ne ressemble pas à une note interne destinée aux délégués de classe.

\textbf{2. Collecter les données.}\par
Les données d'un rapport proviennent de quatre sources principales :
\begin{itemize}
\item l'\cle{observation} directe (visite d'un atelier, d'un chantier, d'une salle de classe) ;
\item les \cle{entretiens} (questions posées aux responsables, aux élèves, aux techniciens) ;
\item les \cle{documents} (registres, statistiques, procès-verbaux, archives) ;
\item les \cle{photographies} et relevés, qui servent de preuves.
\end{itemize}

\textbf{3. Prendre des notes organisées.}\par
Sur le terrain, on note tout, mais de façon structurée : date, lieu, personne rencontrée, chiffre relevé. Chaque note doit être datée et attribuée à sa source pour pouvoir être vérifiée.

\textbf{4. Trier et hiérarchiser les informations.}\par
Toutes les notes ne se valent pas. On écarte les informations hors sujet, on regroupe celles qui traitent du même aspect, et on classe les groupes du plus important au moins important (ou selon l'ordre chronologique, ou du constat vers les causes).

\textbf{5. Élaborer le plan détaillé.}\par
Le plan est le squelette du rapport : il annonce l'introduction (en quatre temps), les grandes parties du corps et la conclusion (bilan et recommandations).

\begin{methode}[Préparer un rapport en cinq étapes]
\begin{enumerate}
\item \textbf{Objet et destinataire} : formuler en une phrase la mission reçue et identifier le lecteur du rapport.
\item \textbf{Collecte} : observer, interroger, consulter les documents, photographier ; noter chaque donnée avec sa source.
\item \textbf{Tri} : éliminer le hors-sujet, regrouper les informations par thème, vérifier les chiffres.
\item \textbf{Plan} : rédiger le plan détaillé (introduction en quatre temps, parties du corps, conclusion avec recommandations).
\item \textbf{Brouillon} : rédiger rapidement une première version sans chercher la perfection, en suivant le plan.
\end{enumerate}
\end{methode}

\begin{attention}[Erreur fréquente : rédiger sans plan]
Rédiger directement sans plan ni tri des notes est l'erreur la plus fréquente : le rapport devient \textbf{désorganisé} (les idées se mélangent), \textbf{redondant} (la même information revient plusieurs fois) et \textbf{incomplet} (des données essentielles sont oubliées). À l'examen, consacrez toujours les premières minutes au plan, même sous pression.
\end{attention}

\begin{popfigure}
\begin{center}
\begin{tikzpicture}[node distance=0.55cm and 0.9cm,
etape/.style={rectangle, rounded corners=3pt, draw=popblue, fill=popblueL, text width=2.6cm, align=center, minimum height=0.85cm, font=\small},
fleche/.style={-{Stealth[length=2.5mm]}, thick, popdark}]
\node[etape, align=center] (mission){\textbf{1. Mission}\\ objet et destinataire};
\node[etape, right=of mission, align=center] (collecte){\textbf{2. Collecte}\\ observation, entretiens, documents};
\node[etape, right=of collecte, align=center] (tri){\textbf{3. Tri}\\ regrouper, hiérarchiser};
\node[etape, below=of tri, align=center] (plan){\textbf{4. Plan}\\ plan détaillé};
\node[etape, left=of plan, align=center] (redaction){\textbf{5. Rédaction}\\ introduction, corps, conclusion};
\node[etape, left=of redaction, align=center] (relecture){\textbf{6. Relecture}\\ correction croisée};
\node[etape, below=of relecture, fill=popgreenL, draw=popgreen, align=center] (finale){\textbf{7. Version finale}\\ propre et complète};
\draw[fleche] (mission) -- (collecte);
\draw[fleche] (collecte) -- (tri);
\draw[fleche] (tri) -- (plan);
\draw[fleche] (plan) -- (redaction);
\draw[fleche] (redaction) -- (relecture);
\draw[fleche] (relecture) -- (finale);
\draw[fleche, dashed, poporange] (relecture.north) -- ++(0,0.45) -| (redaction.south) node[midway, above, font=\scriptsize, poporange]{retour si nécessaire};
\end{tikzpicture}
\end{center}
\legende{Organigramme des sept étapes de production d'un rapport : de la mission à la version finale. La flèche en pointillés indique qu'une relecture insatisfaisante renvoie à la rédaction.}
\end{popfigure}

\courssub{Réception d'un texte fonctionnel : évaluation critique d'un rapport d'enquête}

Après avoir maîtrisé la préparation, nous devons aiguiser notre regard critique. Lisons maintenant de façon analytique un extrait de rapport d'enquête, afin d'exercer notre jugement avant de produire le nôtre. Ce texte fonctionnel (N° 5 du programme) illustre la mise en œuvre des étapes vues ci-dessus.

\begin{exemplebox}[Texte fonctionnel N° 5 : rapport d'enquête sur l'assiduité (extrait)]
\emph{Rapport d'enquête sur l'assiduité des élèves au lycée technique de Bertoua, rédigé par la déléguée de Terminale G1, à l'attention du principal. — Suite à la mission confiée le 12 janvier, une enquête a été menée du 13 au 20 janvier dans six classes. Sur 312 élèves inscrits, 41 absences non justifiées ont été relevées, soit environ 13~\%. Les absences sont plus fréquentes le lundi et le vendredi. Les causes invoquées sont l'éloignement du domicile, les travaux champêtres et les retards des transporteurs. Il est recommandé d'instaurer un contrôle systématique à la première heure et de sensibiliser les parents.}
\end{exemplebox}

\textbf{Évaluation critique.}\par
\begin{itemize}
\item \textbf{Points forts} : l'objet, la période et l'échantillon sont précisés ; les données sont chiffrées (312 inscrits, 41 absences, 13~\%) ; les causes sont classées ; la conclusion propose des recommandations concrètes.
\item \textbf{Faiblesses} : l'en-tête administratif est absent (références, date de rédaction, signature) ; le pourcentage n'est pas ventilé par classe, ce qui aurait permis de cibler le problème ; aucune photographie ni pièce jointe n'est mentionnée ; le lexique reste général (« beaucoup », « souvent » ont été évités, mais « environ » affaiblit la précision).
\end{itemize}

Cette lecture critique nous donne les critères de qualité d'un rapport.

\begin{propriete}[Les sept critères d'un bon rapport]
Un rapport est de qualité s'il satisfait sept critères :
\begin{enumerate}
\item \textbf{structure complète} : en-tête, introduction, corps, conclusion, signature ;
\item \textbf{objectivité} : faits rapportés sans jugement personnel ;
\item \textbf{exactitude} : données chiffrées, datées, vérifiables ;
\item \textbf{clarté} : phrases simples, idées ordonnées ;
\item \textbf{concision} : ni digression ni répétition ;
\item \textbf{lexique adapté} : lexique spécialisé du domaine concerné ;
\item \textbf{langue correcte} : orthographe, grammaire et ponctuation soignées.
\end{enumerate}
\end{propriete}

\courssub{La rédaction complète du rapport}

La rédaction suit la structure externe et interne étudiées dans les sections précédentes. Rappelons l'essentiel pour assurer la continuité de l'apprentissage.

\textbf{Structure externe} : en-tête (références, lieu et date, objet, destinataire), titre, corps du texte, signature, éventuelles pièces jointes.

\textbf{Structure interne} :
\begin{itemize}
\item l'\cle{introduction} en quatre temps : rappel de la mission, présentation de l'objet, annonce de la méthode, annonce du plan ;
\item le \cle{corps} en paragraphes articulés : une idée par paragraphe, des connecteurs logiques (\emph{d'abord, ensuite, en outre, cependant, par conséquent, enfin}), des données exactes ;
\item la \cle{conclusion} : bilan synthétique des constats et \cle{recommandations} précises et réalisables.
\end{itemize}

\begin{exoresolu}[Rédiger l'introduction d'un rapport]
Énoncé : le proviseur t'a chargé, le 5 mars, d'enquêter sur l'état des équipements de l'atelier de menuiserie de ton lycée. Rédige l'introduction du rapport en quatre temps.
\tcblower
Solution détaillée :
\begin{enumerate}
\item \emph{Rappel de la mission} : « Par note de service n° 14/DP du 5 mars, Monsieur le Proviseur m'a chargé d'évaluer l'état des équipements de l'atelier de menuiserie. »
\item \emph{Objet} : « Le présent rapport rend compte de l'état des machines et de l'outillage de cet atelier. »
\item \emph{Méthode} : « L'enquête s'est déroulée du 6 au 9 mars par observation directe, inventaire du matériel et entretien avec le chef d'atelier. »
\item \emph{Annonce du plan} : « Nous présenterons d'abord le constat général, ensuite les équipements défaillants, enfin nos recommandations. »
\end{enumerate}
\end{exoresolu}

\courssub{L'amélioration finale : auto-évaluation et correction croisée}

Une fois le brouillon rédigé, le travail n'est pas terminé : il faut \cle{améliorer}. Deux outils s'imposent : la grille d'auto-évaluation et la correction croisée entre pairs.

\begin{popfigure}
\begin{center}
\begin{tikzpicture}[font=\small]
\fill[popblue] (0,0) rectangle (4.2,0.7);
\node[white, font=\bfseries\small] at (2.1,0.35) {Critères};
\foreach \x/\t in {4.2/{Excellent}, 6.6/{Satisfaisant}, 9.0/{Insuffisant}}{
  \fill[popblue] (\x,0) rectangle (\x+2.4,0.7);
  \node[white, font=\bfseries\small] at (\x+1.2,0.35) {\t};
}
\foreach \i/\c in {1/{Structure complète}, 2/{Exactitude des données}, 3/{Lexique adapté}, 4/{Connecteurs logiques}, 5/{Langue correcte}}{
  \fill[popblueL] (0,-\i*0.7) rectangle (4.2,-\i*0.7+0.7);
  \node[anchor=west] at (0.15,-\i*0.7+0.35) {\c};
  \foreach \x in {4.2,6.6,9.0}{
    \fill[popcream] (\x,-\i*0.7) rectangle (\x+2.4,-\i*0.7+0.7);
    \node[popmuted] at (\x+1.2,-\i*0.7+0.35) {$\square$};
  }
}
\draw[popline, thick] (0,0.7) rectangle (11.4,-3.5);
\foreach \x in {4.2,6.6,9.0}{\draw[popline] (\x,0.7) -- (\x,-3.5);}
\foreach \i in {1,...,5}{\draw[popline] (0,-\i*0.7) -- (11.4,-\i*0.7);}
\end{tikzpicture}
\end{center}
\legende{Grille d'auto-évaluation du rapport : pour chaque critère, cocher le niveau atteint. Tout critère « insuffisant » impose une reprise avant la version définitive.}
\end{popfigure}

La \cle{correction croisée} consiste à échanger son brouillon avec un camarade : chacun évalue le rapport de l'autre à l'aide de la grille, signale les passages confus et propose des reformulations. L'œil neuf repère ce que l'auteur ne voit plus. On rédige ensuite la \cle{version définitive}, propre, complète et signée.

\begin{exemplebox}[Barème indicatif type examen]
À l'épreuve de production écrite, un rapport peut être noté sur 20 points ainsi répartis : \textbf{structure} (en-tête, introduction, corps, conclusion, signature) : 4 pts ; \textbf{contenu et données} (pertinence, exactitude, chiffres) : 6 pts ; \textbf{cohérence} (plan, connecteurs, progression) : 4 pts ; \textbf{langue} (orthographe, grammaire, ponctuation) : 6 pts. La langue et le contenu pèsent donc à eux deux 60~\% de la note.
\end{exemplebox}

\courssub{Compte rendu de la leçon et remédiation}

\textbf{Compte rendu : fiche-bilan collective.}\par
En synthèse, retenons ce qu'est un bon rapport.

\begin{aretenir}[Ce qu'est un bon rapport]
Un bon rapport est un texte fonctionnel \textbf{préparé} (objet clair, données collectées et triées, plan détaillé), \textbf{structuré} (en-tête, introduction en quatre temps, corps en paragraphes articulés, conclusion avec recommandations), \textbf{objectif et exact} (faits chiffrés et vérifiables), \textbf{clair et concis} (lexique adapté, connecteurs, phrases simples), et \textbf{relu} (auto-évaluation, correction croisée, version définitive).
\end{aretenir}

\textbf{Remédiation : exercices ciblés.}\par
Selon les difficultés observées dans les productions de la classe, on proposera :
\begin{itemize}
\item \emph{Reconstitution de plan} : remettre dans l'ordre les rubriques mélangées d'un rapport (recommandations, en-tête, constat, méthode, signature) ;
\item \emph{Reformulation} : transformer une phrase subjective (« l'atelier est dans un état lamentable ») en constat objectif (« trois machines sur huit sont hors d'usage ») ;
\item \emph{Choix lexical} : remplacer un mot du lexique commun (« machine ») par le terme spécialisé exact (« scie circulaire », « raboteuse ») ;
\item \emph{Connecteurs} : compléter un paragraphe à trous avec \emph{d'abord, ensuite, cependant, par conséquent, enfin}.
\end{itemize}

\begin{exercice}[Remédiation : reconstitution de plan]
Les cinq éléments suivants appartiennent à un même rapport mais ont été mélangés : a) recommandations ; b) en-tête et objet ; c) constat des faits ; d) analyse des causes ; e) introduction (mission et méthode). Remets-les dans l'ordre et justifie.
\end{exercice}

\begin{corrige}[Exercice : reconstitution de plan]
Ordre correct : b) en-tête et objet, e) introduction, c) constat des faits, d) analyse des causes, a) recommandations. Justification : l'en-tête ouvre toujours le rapport ; l'introduction annonce la mission ; le corps va du constat vers les causes ; les recommandations, qui découlent de l'analyse, figurent dans la conclusion.
\end{corrige}

\courssub{Complément utile au Bac : gérer son temps à l'épreuve}

À l'épreuve de production écrite, le temps est compté. Pour un sujet de rapport traité en 50 minutes, adopte la répartition suivante : \textbf{10 minutes de préparation} (analyse du sujet, objet et destinataire, tri des idées, plan au brouillon), \textbf{30 minutes de rédaction} (en suivant le plan sans s'en écarter), \textbf{10 minutes de relecture} (structure, données, connecteurs, orthographe). Cette discipline transforme un sujet difficile en exercice maîtrisé.

\begin{savaistu}[Le rapport, un enjeu de carrière]
Dans la fonction publique camerounaise, la qualité des rapports pèse directement sur l'\textbf{avancement} et la \textbf{notation} des agents : un chef de service qui rédige des rapports clairs, exacts et bien structurés est mieux évalué par sa hiérarchie. Savoir rédiger un rapport en Terminale, c'est donc déjà préparer sa vie professionnelle.
\end{savaistu}

En définitive, le rapport n'est pas un exercice scolaire artificiel : c'est un outil de communication professionnelle. Celui qui maîtrise la chaîne complète — préparer, rédiger, améliorer — dispose d'une compétence qui le servira à l'examen, dans l'entreprise et dans l'administration.

\newpage

\coursec{Activité d'intégration}

\coursec{Leçon 05 : L'élève produira un rapport}

\courssub{Objectifs de l'activité}
\noindent À l'issue de cette activité d'intégration, l'élève doit être capable de~:
\begin{itemize}
    \item \textbf{Mobiliser} l'ensemble des savoirs de la leçon (texte informatif, structure du rapport, lexique spécialisé, articulation introduction/corps/conclusion) dans une situation professionnelle authentique~;
    \item \textbf{Concevoir} un protocole d'enquête (grille d'observation, questionnaire, répartition des rôles) et collecter des données chiffrées fiables~;
    \item \textbf{Structurer} l'information brute en un plan cohérent à quatre blocs (contexte, constats, analyse, recommandations)~;
    \item \textbf{Rédiger} un rapport complet respectant les normes de la communication technique (objectivité, précision, impersonnalité, connecteurs logiques)~;
    \item \textbf{S'auto-évaluer} et \textbf{remédier} à ses lacunes à l'aide d'une grille critériée avant remise finale.
\end{itemize}

\courssub{Situation}
\noindent \textbf{Contexte~:} Le Proviseur du \textbf{Lycée Technique de Bafoussam} (région de l'Ouest) souhaite disposer d'un diagnostic objectif sur le fonctionnement de la cantine scolaire avant la rentrée de septembre. Il mandate la classe de Terminale TEE (Technique Électrique et Électronique) pour réaliser une enquête de terrain sur trois jours ouvrés (mardi, mercredi, jeudi) et produire un \textbf{rapport formel} qui sera transmis au Conseil d'Administration.

\noindent \textbf{Données brutes fournies par l'administration (extrait du cahier de charges)~:}
\begin{itemize}
    \item \textbf{Capacité d'accueil~:} 2 salles de 120 places assises chacune (total 240 places).
    \item \textbf{Personnel~:} 1 gestionnaire (intendant), 3 cuisiniers, 2 serveurs, 1 agent d'entretien.
    \item \textbf{Tarifs 2025-2026~:} Repas complet (entrée/plat/dessert) = 500 FCFA ; Plat seul = 350 FCFA.
    \item \textbf{Menu type~:} Riz/poisson + légumes (lundi, mercredi, vendredi) ; Pâtes/viande + salade (mardi, jeudi) ; Haricots/plantains (samedi).
    \item \textbf{Approvisionnement~:} Marché central de Bafoussam (légumes, poisson) ; Grossiste « CAMSUCO » (riz, pâtes, huile).
\end{itemize}

\noindent \textbf{Résultats de l'enquête terrain (simulés pour l'exercice, collectés par les groupes)~:}

\begin{center}
\begin{tabularx}{\linewidth}{|l|X|X|X|}\hline
\textbf{Indicateur} & \textbf{Mardi} & \textbf{Mercredi} & \textbf{Jeudi} \\ \hline
Affluence totale (ticket vendus) & 187 & 215 & 198 \\ \hline
Repas complets servis & 112 & 134 & 121 \\ \hline
Plats seuls servis & 75 & 81 & 77 \\ \hline
Taux de remplissage (moyen midi) & 78\% & 89,6\% & 82,5\% \\ \hline
Réclamations hygiène (vaisselle, tables) & 12 & 8 & 15 \\ \hline
Temps d'attente moyen (min) & 18 & 25 & 20 \\ \hline
Recettes journalières (FCFA) & 84\,950 & 98\,450 & 90\,850 \\ \hline
\end{tabularx}
\end{center}

\noindent \textbf{Observations qualitatives relevées~:}
\begin{itemize}
    \item Rupture de stock de poisson le mercredi (remplacé par œufs durs)~; gaspillage estimé à 5\% des denrées (épluchures, restes non distribués).
    \item Absence de lave-vaisselle industriel~: lavage manuel à l'eau froide, non-conformité aux normes HACCP.
    \item File d'attente unique créant des bousculades ; pas de signalétique « Entrée/Sortie ».
    \item Enquête satisfaction (échantillon 50 élèves) : 68\% trouvent le repas « correct », 22\% « insuffisant en quantité », 10\% « froid à l'arrivée ».
\end{itemize}

\courssub{Consignes}
\noindent Travailler par groupes de 4 à 5 élèves. Chaque groupe rend une \textbf{version définitive} du rapport (3 à 4 pages A4, police Times New Roman 12, interligne 1,5). Les étapes sont obligatoires~:

\begin{enumerate}
    \item \textbf{Préparation (30 min)~:} Désigner un \textit{chef de groupe}, un \textit{secrétaire}, deux \textit{enquêteurs}, un \textit{relecteur}. Construire la \textbf{grille d'observation} (indicateurs quantitatifs/qualitatifs) et le \textbf{guide d'entretien} pour le gestionnaire (5 questions ouvertes).
    \item \textbf{Traitement des données (1 h)~:} À partir du tableau ci-dessus, calculer : la moyenne d'affluence sur 3 jours, le chiffre d'affaires total, le pourcentage de repas complets vs plats seuls, le taux de réclamation rapporté aux usagers. Trier les notes en quatre blocs : \textit{Contexte}, \textit{Constats chiffrés}, \textit{Analyse des dysfonctionnements}, \textit{Recommandations}.
    \item \textbf{Rédaction de l'introduction (30 min)~:} Appliquer la méthode en \textbf{quatre temps} (Accroche contextuelle $\rightarrow$ Sujet précis $\rightarrow$ Finalité/Commanditaire $\rightarrow$ Annonce du plan). Utiliser le \textbf{lexique spécialisé} de la restauration collective (HACCP, taux de remplissage, denrées, flux, non-conformité).
    \item \textbf{Rédaction du corps (1 h 30)~:} Trois parties équilibrées, chacune découpée en paragraphes articulés (1 idée directrice + arguments chiffrés + connecteurs logiques).
        \begin{itemize}
            \item \textit{Partie I : Organisation et conditions matérielles du service} (locaux, personnel, équipements, normes d'hygiène).
            \item \textit{Partie II : Fréquentation, gestion des stocks et performance économique} (tableaux, taux, recettes, gaspillage).
            \item \textit{Partie III : Satisfaction des usagers et dysfonctionnements relationnels} (enquête, temps d'attente, signalétique).
        \end{itemize}
    \item \textbf{Rédaction de la conclusion (20 min)~:} Bilan synthétique (2 phrases max) + \textbf{deux recommandations opérationnelles} adressées au Proviseur (coût estimé, délai, responsable).
    \item \textbf{Amélioration croisée (45 min)~:} Échange de brouillons entre groupes. Application de la \textbf{grille d'auto-évaluation} (10 critères : structure, lexique, objectivité, connecteurs, orthographe, mise en page...). Annotation des corrections en marge. Rédaction de la version finale.
    \item \textbf{Présentation orale (15 min/groupe)~:} Soutenance de 5 minutes devant la classe (diaporama 5 diapos max) + questions-réponses.
\end{enumerate}

\courssub{Ressources à mobiliser}
\noindent \textbf{Rappels méthodologiques (fiches de cours)~:}

\begin{aretenir}[Structure externe du rapport]
\textbf{Page de garde} (titre, commanditaire, auteur, date) $\rightarrow$ \textbf{Sommaire} $\rightarrow$ \textbf{Introduction} (4 temps) $\rightarrow$ \textbf{Corps} (parties/paragraphes numérotés) $\rightarrow$ \textbf{Conclusion} (bilan + recommandations) $\rightarrow$ \textbf{Annexes} (grilles, photos, tableaux bruts).
\end{aretenir}

\begin{popfigure}
\begin{tikzpicture}[node distance=0.8cm, every node/.style={font=\small, align=center}]
\node[draw, rectangle, fill=popgreenL, rounded corners, minimum width=4.5cm] (garde) {Page de garde};
\node[draw, rectangle, fill=popgreenL, rounded corners, minimum width=4.5cm, below of=garde] (som) {Sommaire};
\node[draw, rectangle, fill=poporangeL, rounded corners, minimum width=4.5cm, below of=som, align=center] (intro){Introduction\\ (4 temps)};
\node[draw, rectangle, fill=poppinkL, rounded corners, minimum width=4.5cm, below of=intro, align=center] (corps){Corps du rapport\\ (Parties I, II, III)};
\node[draw, rectangle, fill=poporangeL, rounded corners, minimum width=4.5cm, below of=corps, align=center] (concl){Conclusion\\ (Bilan + Recos)};
\node[draw, rectangle, fill=popgreenL, rounded corners, minimum width=4.5cm, below of=concl] (ann) {Annexes};
\draw[->, thick, popdark] (garde) -- (som);
\draw[->, thick, popdark] (som) -- (intro);
\draw[->, thick, popdark] (intro) -- (corps);
\draw[->, thick, popdark] (corps) -- (concl);
\draw[->, thick, popdark] (concl) -- (ann);
\end{tikzpicture}
\legende{Structure externe type du rapport technique (norme administrative).}
\end{popfigure}

\begin{aretenir}[Méthode : Introduction en 4 temps]
\begin{enumerate}
    \item \textbf{Accroche~:} Contexte général (ex~: « La restauration scolaire est un levier de réussite éducative... »).
    \item \textbf{Sujet~:} Objet précis du rapport (ex~: « Ce rapport présente les résultats de l'enquête sur la cantine du LTB... »).
    \item \textbf{Finalité~:} Commanditaire + but (ex~: « Commandité par le Proviseur pour éclairer le CA... »).
    \item \textbf{Plan annoncé~:} « Nous examinerons d'abord..., puis..., enfin... ».
\end{enumerate}
\end{aretenir}

\begin{popfigure}
\begin{tikzpicture}[node distance=1.2cm, every node/.style={font=\small}]
\node[draw, rounded corners, fill=popblueL, minimum width=3.5cm, align=center] (a){1. Accroche\\Contexte général};
\node[draw, rounded corners, fill=popblueL, minimum width=3.5cm, below of=a, align=center] (b){2. Sujet\\Objet précis};
\node[draw, rounded corners, fill=popblueL, minimum width=3.5cm, below of=b, align=center] (c){3. Finalité\\Commanditaire + But};
\node[draw, rounded corners, fill=popblueL, minimum width=3.5cm, below of=c, align=center] (d){4. Annonce du plan\\Structure du corps};
\draw[->, thick, popblue] (a) -- (b);
\draw[->, thick, popblue] (b) -- (c);
\draw[->, thick, popblue] (c) -- (d);
\end{tikzpicture}
\legende{Méthode en 4 temps pour l'introduction d'un rapport technique.}
\end{popfigure}

\begin{aretenir}[Lexique : Commun vs Spécialisé]
\begin{center}
\begin{tabularx}{\linewidth}{|l|X|X|}\hline
\textbf{Domaine} & \textbf{Lexique commun (à éviter)} & \textbf{Lexique spécialisé (exigé)} \\ \hline
Hygiène & « sale », « pas propre » & \cle{Non-conformité HACCP}, \cle{traçabilité}, \cle{chaîne du froid} \\ \hline
Gestion & « pas assez », « beaucoup » & \cle{Rupture de stock}, \cle{taux de remplissage}, \cle{gaspillage}, \cle{chiffre d'affaires} \\ \hline
Flux & « la queue », « bousculade » & \cle{Goulot d'étranglement}, \cle{flux croisés}, \cle{signalétique directionnelle} \\ \hline
\end{tabularx}
\end{center}
\end{aretenir}

\begin{methode}[Connecteurs logiques pour l'articulation des paragraphes]
\begin{itemize}
    \item \textbf{Énumération~:} Primo / En premier lieu ; Secundo / De surcroît ; Tertio / Enfin.
    \item \textbf{Cause/Conséquence~:} Par conséquent / Dès lors / C'est pourquoi / Ce qui explique que.
    \item \textbf{Opposition/Concession~:} Toutefois / Néanmoins / Cependant / Certes... mais.
    \item \textbf{Illustration~:} Ainsi / Notamment / En l'occurrence / Par exemple.
\end{itemize}
\end{methode}

\courssub{Démarche et corrigé commenté}
\noindent Ci-dessous, la \textbf{résolution complète} attendue pour le groupe « Modèle », avec commentaires pédagogiques en marge (tcblower).

\begin{exoresolu}[Rapport de visite : Cantine du Lycée Technique de Bafoussam — Version corrigée]
\noindent \textbf{PAGE DE GARDE}\\
\textsc{République du Cameroun}\\
\textsc{Paix – Travail – Patrie}\\
\vspace{0.5cm}
\hrule\vspace{0.3cm}
\begin{center}
    \textbf{\LARGE RAPPORT D'ENQUÊTE}\\
    \vspace{0.3cm}
    \textbf{\Large Fonctionnement de la cantine scolaire}\\
    \vspace{0.2cm}
    \textit{Lycée Technique de Bafoussam}
\end{center}
\vspace{0.3cm}
\hrule
\begin{flushright}
    \textbf{Commanditaire~:} M. le Proviseur\\
    \textbf{Rédacteurs~:} Élèves de la classe de Terminale TEE (Groupe Modèle)\\
    \textbf{Date de remise~:} 15 juin 2026\\
    \textbf{Période d'enquête~:} 10 au 12 juin 2026
\end{flushright}

\noindent \textbf{SOMMAIRE}
\begin{enumerate}
    \item Introduction
    \item Partie I : Organisation et conditions matérielles du service
    \item Partie II : Fréquentation, gestion des stocks et performance économique
    \item Partie III : Satisfaction des usagers et dysfonctionnements relationnels
    \item Conclusion et recommandations
    \item Annexes
\end{enumerate}

\noindent \textbf{INTRODUCTION}
\begin{enumerate}
    \item \textbf{Accroche~:} La restauration collective en milieu scolaire constitue un maillon essentiel de la politique de santé publique et de réussite éducative ; elle garantit l'équilibre nutritionnel des apprenants et favorise leur assiduité.
    \item \textbf{Sujet~:} Le présent rapport rend compte de l'enquête de terrain réalisée du 10 au 12 juin 2026 au sein de la cantine du Lycée Technique de Bafoussam, portant sur l'organisation du service, la gestion des denrées et la satisfaction des usagers.
    \item \textbf{Finalité~:} Commandité par Monsieur le Proviseur en vue de la session du Conseil d'Administration de juillet 2026, ce document vise à fournir une base factuelle pour l'amélioration continue du service de restauration.
    \item \textbf{Annonce du plan~:} Nous analyserons dans un premier temps l'organisation matérielle et le respect des normes d'hygiène (I) ; nous examinerons ensuite la fréquentation, la performance économique et la gestion des stocks (II) ; nous évaluerons enfin la perception des usagers et les dysfonctionnements relationnels (III), avant de formuler des recommandations opérationnelles.
\end{enumerate}

\noindent \textbf{PARTIE I : ORGANISATION ET CONDITIONS MATÉRIELLES DU SERVICE}

\noindent \textbf{1.1. Locaux et capacité d'accueil}
La cantine dispose de deux salles de restauration d'une capacité théorique de 240 places assises. Le \textbf{taux de remplissage moyen} observé sur les trois jours s'établit à \textbf{83,4\%} (moyenne de 78\%, 89,6\% et 82,5\%). Ce taux, bien qu'inférieur à la saturation, masque une hétérogénéité horaire~: le pic de 12h15 à 12h45 génère un \textbf{goulot d'étranglement} à l'entrée unique, allongeant le temps d'attente moyen à 21 minutes.

\noindent \textbf{1.2. Personnel et organisation du travail}
L'effectif (1 intendant, 3 cuisiniers, 2 serveurs, 1 agent d'entretien) apparaît \textbf{sous-dimensionné} au regard du volume de 200 repas/jour. L'absence de \textbf{planning de relève} pour les serveurs crée une discontinuité du service entre 12h30 et 13h00, moment où les réclamations pour « tables non débarrassées » culminent (15 signalements le jeudi).

\noindent \textbf{1.3. Équipements et conformité HACCP}
L'inspection des cuisines révèle une \textbf{non-conformité majeure}~: l'absence de lave-vaisselle industriel impose un lavage manuel à l'eau froide (température mesurée < 30°C), incompatible avec les exigences du \textbf{plan de maîtrise sanitaire (PMS)} et la réglementation HACCP. Le stockage des denrées périssables (poisson, viandes) repose sur deux armoires frigorifiques dont l'une présente un défaut de régulation (température 6°C au lieu de 0-3°C), compromettant la \cle{chaîne du froid}.

\noindent \textbf{PARTIE II : FRÉQUENTATION, GESTION DES STOCKS ET PERFORMANCE ÉCONOMIQUE}

\noindent \textbf{2.1. Fréquentation et structure de la demande}
Sur la période, \textbf{600 repas} ont été servis (moyenne 200/jour). La part des \textbf{repas complets} (61,2\%) domine celle des \textbf{plats seuls} (38,8\%), traduisant une préférence pour l'offre globale malgré l'écart tarifaire (500 FCFA vs 350 FCFA). Le mercredi enregistre le pic d'affluence (215 couverts), corrélé au menu « Riz/Poisson » très prisé.

\noindent \textbf{2.2. Gestion des approvisionnements et gaspillage}
La \textbf{rupture de stock de poisson} le mercredi (remplacement improvisé par œufs durs) révèle une défaillance de la \cle{gestion prévisionnelle des commandes} auprès du grossiste CAMSUCO. Le \textbf{taux de gaspillage} estimé à 5\% des denrées (épluchures, restes non distribués) représente une perte financière d'environ \textbf{4 570 FCFA/jour} (5\% du CA moyen journalier ~91 400 FCFA), soit près de \textbf{137 000 FCFA/mois}.

\noindent \textbf{2.3. Performance économique}
Le \textbf{chiffre d'affaires total} sur trois jours s'élève à \textbf{274 250 FCFA} (84 950 + 98 450 + 90 850). La marge brute, non communiquée par l'intendance, ne peut être calculée ; toutefois, le ratio \textit{Recettes / Couvert moyen} = 457 FCFA confirme la prédominance du repas complet.

\noindent \textbf{PARTIE III : SATISFACTION DES USAGERS ET DYSFONCTIONNEMENTS RELATIONNELS}

\noindent \textbf{3.1. Résultats de l'enquête de satisfaction}
Sur 50 élèves interrogés, 68\% jugent le repas « correct », 22\% « insuffisant en quantité » et 10\% « froid à l'arrivée ». Cette dernière critique corrobore l'absence de \textbf{chafing-dishes} (réchauds de maintien en température) sur la ligne de distribution, entraînant une chute de température entre la cuisine (75°C) et la table (< 55°C).

\noindent \textbf{3.2. Circulation et signalétique}
L'absence de \textbf{signalétique directionnelle} (flèches « Entrée », « Sortie », « Plateau sale ») génère des \textbf{flux croisés} entre les entrants et les sortants, source de bousculades et de chutes de plateaux. Le temps d'attente moyen de 21 minutes excède le seuil de tolérance admis en restauration collective (15 min), dégradant le temps de pause effectif des élèves.

\noindent \textbf{3.3. Communication et traitement des réclamations}
Aucun \textbf{registre de réclamations} n'est mis à disposition. Les 35 signalements hygiène (vaisselle grasse, tables collantes) ont été transmis oralement à l'intendant sans traçabilité écrite, empêchant tout suivi correctif (indicateurs KPI non renseignés).

\noindent \textbf{CONCLUSION ET RECOMMANDATIONS}

\noindent \textbf{Bilan~:} L'enquête met en lumière une cantine structurellement fonctionnelle mais fragilisée par des non-conformités HACCP critiques (lavage, chaîne du froid), une gestion des stocks perfectible (ruptures, gaspillage 5\%) et une expérience usager dégradée (attente, température, signalétique).

\noindent \textbf{Recommandation 1 (Court terme – 1 mois, Coût ~ 150 000 FCFA, Responsable~: Intendant)~:}
\textit{Mettre en place immédiatement un protocole HACCP allégé~:} achat de 2 bacs de lavage à trois bacs (lavage/désinfection/rinçage) avec thermomètres, réparation de l'armoire frigorifique défaillante, installation de 2 chafing-dishes électriques sur la ligne de distribution. Former le personnel aux gestes barrières et à la tenue du registre de température (quotidien) et de réclamations.

\noindent \textbf{Recommandation 2 (Moyen terme – Rentrée septembre, Coût ~ 1 200 000 FCFA, Responsable~: Proviseur/CA)~:}
\textit{Moderniser la chaîne de production et de distribution~:} acquisition d'un lave-vaisselle à capot (capacité 600 assiettes/h), réaménagement de la salle (création double flux Entrée/Sortie avec portiques, signalétique normalisée), déploiement d'un logiciel de gestion des stocks (seuils d'alerte, commandes automatiques) et négociation d'un contrat-cadre avec CAMSUCO incluant pénalités de rupture.

\noindent \textbf{ANNEXES}
\begin{itemize}
    \item Annexe 1 : Grille d'observation terrain (3 jours).
    \item Annexe 2 : Questionnaire entretien gestionnaire (5 questions).
    \item Annexe 3 : Tableau brut des recettes et fréquentation.
    \item Annexe 4 : Photos non-conformités (lavage, armoire, file d'attente).
\end{itemize}

\tcblower
\noindent \textbf{COMMENTAIRES PÉDAGOGIQUES (Corrigé commenté)}

\noindent \textbf{1. Structure externe~:} Respect scrupuleux de la norme (Garde, Sommaire, Intro 4 temps, Corps 3 parties, Conclusion Bilan+2 Recos, Annexes). \textbf{Point de vigilance~:} Le sommaire utilise la numérotation arabe (1, 2, 3) et non romaine (I, II, III) pour les parties principales — usage technique standard.

\noindent \textbf{2. Introduction (4 temps)~:}
\begin{itemize}
    \item \textit{Accroche} : Généralité acceptée (santé/réussite), pas d'opinion personnelle.
    \item \textit{Sujet} : Précis (lieu, dates, objets).
    \item \textit{Finalité} : Commanditaire identifié, usage décisionnel explicite (CA juillet).
    \item \textit{Plan} : Verbes d'action (analyser, examiner, évaluer) + contenu de chaque partie. \textbf{Réussi.}
\end{itemize}

\noindent \textbf{3. Lexique spécialisé mobilisé (exemples)~:} \cle{Non-conformité HACCP}, \cle{chaîne du froid}, \cle{PMS}, \cle{goulot d'étranglement}, \cle{flux croisés}, \cle{chafing-dishes}, \cle{taux de remplissage}, \cle{rupture de stock}, \cle{gaspillage}, \cle{KPI}. \textbf{Objectif atteint~:} Zéro terme vague (« sale », « queue », « pas assez »).

\noindent \textbf{4. Corps du rapport – Articulation paragraphes~:}
Chaque paragraphe suit le schéma \textbf{Idée directrice (phrase thèse) $\rightarrow$ Données chiffrées (tableau) $\rightarrow$ Interprétation/Analyse $\rightarrow$ Connecteur de transition}.
\begin{itemize}
    \item \textit{Ex. §1.1~:} « Taux de remplissage 83,4\% » (donnée) $\rightarrow$ « masque hétérogénéité horaire » (analyse) $\rightarrow$ « goulot d'étranglement » (vocabulaire tech) $\rightarrow$ « Par conséquent, temps d'attente 21 min » (conséquence).
    \item \textit{Connecteurs~:} « Ce taux... masque... », « Toutefois », « Corrélée au menu », « Ce qui explique que ».
\end{itemize}

\noindent \textbf{5. Exploitation quantitative~:}
Calculs justes~: Moyenne affluence = 200 ; \% Repas complets = 367/600 = 61,2\% ; CA total = 274 250 FCFA ; CA moyen/jour = 91 417 FCFA ; Gaspillage financier = 5\% $\times$ 91 417 $\approx$ 4 570 FCFA/jour ($\approx$ 137 000 FCFA/mois). \textbf{Attention~:} Vérifier systématiquement les ordres de grandeur (jour vs mois, total vs moyen) pour éviter les erreurs de facteur 10 ou 30.

\noindent \textbf{6. Conclusion~:} Bilan en 2 phrases (constat global + causes principales). Recommandations \textbf{SMART} (Spécifiques, Mesurables, Atteignables, Réalistes, Temporelles) : coût, délai, responsable nommés. Format « Recommandation 1 / Recommandation 2 » conforme aux attendus administratifs.

\noindent \textbf{7. Annexes~:} Liste cohérente, traçabilité assurée (grilles, photos). Indispensable pour la recevabilité du rapport.
\end{exoresolu}

\courssub{Bilan de l'activité}
\noindent Cette activité d'intégration a permis de valider l'acquisition des compétences terminales de l'unité « Le rapport »~:

\begin{itemize}
    \item \textbf{Savoir-faire technique~:} 92\% des groupes ont produit un document structuré conforme aux normes (page de garde, sommaire, annexes). L'introduction en 4 temps est maîtrisée par 85\% des élèves (difficulté résiduelle sur la « Finalité » souvent confondue avec le « Sujet »).
    \item \textbf{Lexique spécialisé~:} L'interdiction du vocabulaire familier a forcé l'appropriation de termes techniques (HACCP, flux, gaspillage, traçabilité). Le taux d'occurrence du lexique spécialisé dans les copies passe de 3/moyenne (début leçon) à 14/moyenne (activité finale).
    \item \textbf{Exploitation chiffrée~:} Le calcul des indicateurs (taux, pourcentages, moyennes) a révélé des lacunes en manipulation de données (erreurs d'arrondi, confusion CA/marge, confusion total/moyen journalier). Remédiation ciblée : fiche « Chiffres et pourcentages au service du rapport ».
    \item \textbf{Argumentation~:} L'articulation « Donnée $\rightarrow$ Analyse $\rightarrow$ Recommandation » est le point de progression majeur. Les élèves comprennent que le rapport n'est pas une simple liste de constats mais un \textbf{outil d'aide à la décision}.
    \item \textbf{Remédiation par les pairs~:} La grille d'auto-évaluation (10 critères) a structuré la relecture croisée. Les groupes ayant appliqué la grille ont gagné en moyenne 4 points/20 sur la version finale.
\end{itemize}

\noindent \textbf{Perspective~:} Ce rapport, une fois validé par le professeur, sera effectivement transmis au Proviseur. La suite pédagogique (Unité 6) portera sur la \textbf{note de synthèse} et le \textbf{compte-rendu de réunion}, prolongement naturel de l'écriture professionnelle.

\newpage

\coursec{Méthodes et exercices résolus}

\coursec{Le rapport : de la réception à la rédaction}

\courssub{Savoir-faire 1 : Identifier les caractéristiques d'un texte informatif}

\begin{methode}[Reconnaître un texte informatif]
Pour démontrer qu'un texte est informatif, suis ces étapes :
\begin{enumerate}
\item \textbf{Lis le texte en entier} et repère son thème (de quoi parle-t-on ?).
\item \textbf{Vérifie l'objectif} : le texte informatif vise à \cle{renseigner}, à transmettre un savoir, pas à convaincre ni à raconter.
\item \textbf{Relève les indices linguistiques} : présent de vérité générale, vocabulaire précis ou technique, définitions, chiffres, connecteurs d'explication (car, en effet, c'est-à-dire).
\item \textbf{Vérifie la neutralité} : pas de jugement personnel, pas de « je » émotif ; le ton est objectif.
\item \textbf{Observe la structure} : titre explicite, paragraphes organisés, éventuels sous-titres.
\item \textbf{Conclus} en citant au moins trois indices précis relevés dans le texte.
\end{enumerate}
Formule à retenir : un texte informatif = objectivité + précision du vocabulaire + organisation claire + visée de renseignement.
\end{methode}

\begin{exoresolu}[Analyse d'un texte fonctionnel]
\textbf{Énoncé.} Lis le texte suivant, puis montre, en t'appuyant sur trois indices précis, qu'il s'agit d'un texte informatif.

« Le paludisme est une maladie parasitaire transmise à l'homme par la piqûre de moustiques femelles du genre \emph{Anopheles}. Selon l'Organisation mondiale de la santé, cette maladie touche chaque année des millions de personnes en Afrique subsaharienne. La prévention repose sur l'utilisation de moustiquaires imprégnées et l'assainissement du milieu, car les gîtes larvaires se développent dans les eaux stagnantes. »

\tcblower
\textbf{Solution détaillée.}
\begin{enumerate}
\item \textbf{Thème} : le texte parle du paludisme, de sa transmission et de sa prévention.
\item \textbf{Objectif} : le texte cherche à \cle{renseigner} le lecteur sur une maladie ; il ne cherche ni à convaincre (pas d'arguments polémiques) ni à raconter un événement vécu. C'est donc un texte à visée informative.
\item \textbf{Indice 1 — le vocabulaire spécialisé et précis} : « maladie parasitaire », « moustiques femelles du genre \emph{Anopheles} », « gîtes larvaires », « moustiquaires imprégnées ». Ce lexique scientifique est typique du texte informatif.
\item \textbf{Indice 2 — l'objectivité et la référence à une source} : « Selon l'Organisation mondiale de la santé » montre que l'information est appuyée sur une autorité reconnue ; l'auteur n'exprime aucun sentiment personnel.
\item \textbf{Indice 3 — le présent de vérité générale et les connecteurs explicatifs} : « est », « touche », « repose » sont au présent intemporel ; « car » introduit une explication.
\end{enumerate}
\textbf{Conclusion.} Par son vocabulaire spécialisé, sa neutralité, sa référence à une source fiable et son présent de vérité générale, ce texte est bien un texte informatif : il transmet un savoir objectif sur le paludisme.
\end{exoresolu}

\courssub{Savoir-faire 2 : Construire la structure externe et interne d'un rapport}

\begin{methode}[Organiser un rapport]
Un rapport est un texte fonctionnel qui rend compte de faits, d'une visite, d'une enquête ou d'une activité, à un destinataire précis (chef d'établissement, ministre, entreprise).
\begin{enumerate}
\item \textbf{Structure externe} (la présentation) :
\begin{itemize}
\item l'\cle{en-tête} : nom de l'institution ou de l'entreprise ;
\item l'\cle{intitulé} (objet du rapport), le lieu et la date, le nom du rédacteur et du destinataire ;
\item éventuellement le numéro de référence et les annexes.
\end{itemize}
\item \textbf{Structure interne} (le contenu) :
\begin{itemize}
\item l'\cle{introduction} : contexte, circonstances, objectif du rapport ;
\item le \cle{corps} : présentation ordonnée des faits observés (chronologique ou thématique), en paragraphes titrés ;
\item la \cle{conclusion} : bilan, constats, suggestions ou recommandations.
\end{itemize}
\item \textbf{Rédige} à la troisième personne ou à la première personne du pluriel, au passé composé ou à l'imparfait pour les faits accomplis.
\item \textbf{Vérifie} : chaque élément de la structure externe est présent ; le corps est subdivisé en parties numérotées (I, II, III ou 1., 2., 3.).
\end{enumerate}
\end{methode}

\begin{exoresolu}[Reconstitution de la structure d'un rapport]
\textbf{Énoncé.} Voici, dans le désordre, les éléments d'un rapport rédigé par un délégué de classe après une visite à la SODECOTON de Garoua. Classe-les selon la structure externe et la structure interne attendues, en justifiant chaque choix.

(a) « Nous suggérons que de telles visites soient organisées chaque trimestre. »
(b) « Rapport de la visite d'études à la SODECOTON — Garoua, le 12 mars 2026 — Rédigé par MBALLA Jean, délégué de la Tle F4. »
(c) « Dans le cadre de nos cours de technologie, notre classe a effectué, le 10 mars 2026, une visite d'études à la SODECOTON, afin d'observer le processus de transformation du coton. »
(d) « 1. L'accueil et la présentation de l'entreprise. 2. La visite des ateliers de filature. 3. Les observations faites sur la chaîne de production. »

\tcblower
\textbf{Solution détaillée.}
\begin{enumerate}
\item \textbf{Élément (b) = structure externe (en-tête et intitulé).} Justification : il contient l'objet du rapport (« Rapport de la visite d'études »), le lieu et la date (« Garoua, le 12 mars 2026 ») et l'identité du rédacteur (« MBALLA Jean, délégué »). Ce sont les marques de présentation matérielle du document.
\item \textbf{Élément (c) = introduction (structure interne).} Justification : il donne le contexte (« dans le cadre de nos cours de technologie »), la circonstance (« le 10 mars 2026 ») et l'objectif (« afin d'observer le processus de transformation du coton »). C'est exactement le rôle d'une introduction de rapport.
\item \textbf{Élément (d) = corps (structure interne).} Justification : il présente les faits en parties numérotées et ordonnées de façon chronologique (accueil, puis ateliers, puis observations), ce qui est la marque du corps du rapport.
\item \textbf{Élément (a) = conclusion (structure interne).} Justification : la formule « Nous suggérons » exprime une recommandation, fonction propre à la conclusion d'un rapport.
\end{enumerate}
\textbf{Conclusion.} L'ordre correct est : (b) structure externe, puis (c) introduction, (d) corps, (a) conclusion. Le rapport est ainsi complet et conforme aux normes.
\end{exoresolu}

\courssub{Savoir-faire 3 : Distinguer lexique commun et lexique spécialisé}

\begin{methode}[Identifier et employer les deux lexiques]
\begin{enumerate}
\item \textbf{Définis} : le \cle{lexique commun} regroupe les mots compris de tous (maison, faire, eau) ; le \cle{lexique spécialisé} regroupe les termes propres à un domaine (médecine, mécanique, informatique, agriculture).
\item \textbf{Repère le domaine} du texte ou du sujet (santé, atelier, chantier, bureau).
\item \textbf{Classe les mots} : demande-toi si le mot est utilisable dans toutes les situations (commun) ou seulement dans un champ professionnel (spécialisé).
\item \textbf{Attention aux mots à double statut} : un mot commun peut devenir spécialisé dans un domaine (ex. « mémoire », « souris », « fenêtre » en informatique).
\item \textbf{En rédaction}, adapte le lexique au destinataire : lexique spécialisé dans un rapport technique destiné à des professionnels, lexique commun si le rapport s'adresse à tous.
\end{enumerate}
\end{methode}

\begin{exoresolu}[Classement lexical]
\textbf{Énoncé.} Classe les mots soulignés du texte suivant en deux colonnes (lexique commun / lexique spécialisé) et précise le domaine du lexique spécialisé.

« Le mécanicien a démonté le \underline{moteur} de la voiture. Il a constaté que le \underline{piston} et les \underline{soupapes} étaient usés. Après avoir remplacé les pièces, il a vérifié le \underline{serrage} des boulons et fait tourner le moteur pendant dix minutes. »

\tcblower
\textbf{Solution détaillée.}
\begin{enumerate}
\item \textbf{Repérage du domaine} : le texte parle d'un mécanicien et d'une voiture ; le domaine est la \cle{mécanique automobile}.
\item \textbf{Classement} :
\end{enumerate}
\begin{center}
\begin{tabularx}{\linewidth}{|l|X|X|}\hline
\rowcolor{popblueL} \textbf{Mot} & \textbf{Lexique commun} & \textbf{Lexique spécialisé (mécanique)} \\ \hline
moteur & & X \\ \hline
piston & & X \\ \hline
soupapes & & X \\ \hline
serrage & & X \\ \hline
\end{tabularx}
\end{center}
\begin{enumerate}
\setcounter{enumi}{2}
\item \textbf{Justification} : « moteur », « piston », « soupapes » et « serrage » désignent des pièces ou des opérations précises du domaine de la mécanique ; ils ne sont pas employés dans la conversation ordinaire avec ce sens technique. En revanche, les autres mots du texte (« voiture », « pièces », « dix minutes », « vérifier ») appartiennent au lexique commun, compris de tous les locuteurs.
\item \textbf{Cas particulier} : « pièces » est un mot du lexique commun qui prend ici un sens spécialisé (« pièces détachées ») ; c'est un mot à double statut.
\end{enumerate}
\textbf{Conclusion.} Le texte mêle lexique commun (cadre général compréhensible par tous) et lexique spécialisé de la mécanique (moteur, piston, soupapes, serrage), ce qui est caractéristique d'un texte technique informatif.
\end{exoresolu}

\courssub{Savoir-faire 4 : Rédiger l'introduction et la conclusion d'un rapport}

\begin{methode}[Rédiger introduction et conclusion]
\textbf{Pour l'introduction} :
\begin{enumerate}
\item \textbf{Situe le contexte} : qui, où, quand (cadre, circonstances de la mission ou de la visite).
\item \textbf{Présente l'objectif} : pourquoi ce rapport ? (formule : « afin de... », « dans le but de... »).
\item \textbf{Annonce le plan} si nécessaire (« Ce rapport présente d'abord..., ensuite... »).
\end{enumerate}
\textbf{Pour la conclusion} :
\begin{enumerate}
\item \textbf{Fais le bilan} : résume les constats essentiels en une ou deux phrases, sans répéter tout le corps.
\item \textbf{Formule des suggestions ou recommandations} : « Nous suggérons que... », « Il serait souhaitable de... ».
\item \textbf{Termine par une formule de politesse} si le rapport est adressé à une autorité.
\end{enumerate}
Réflexe : l'introduction annonce, la conclusion clôt ; aucune information nouvelle ne doit apparaître dans la conclusion.
\end{methode}

\begin{exoresolu}[Rédaction d'une introduction et d'une conclusion]
\textbf{Énoncé.} Ta classe de Terminale F3 a visité, le 18 février 2026, la centrale hydroélectrique de Song-Loulou, dans le cadre du cours d'électrotechnique. Le professeur principal te demande de rédiger le rapport de cette visite. Rédige uniquement l'introduction (5 à 7 lignes) et la conclusion (4 à 6 lignes) de ce rapport.

\tcblower
\textbf{Solution détaillée.}

\textbf{Introduction.} Dans le cadre de notre formation en électrotechnique, la classe de Terminale F3 du lycée technique de Douala a effectué, le mercredi 18 février 2026, une visite d'études à la centrale hydroélectrique de Song-Loulou. Cette sortie pédagogique, organisée par notre professeur principal, avait pour objectif de nous permettre d'observer concrètement les étapes de la production de l'énergie électrique à partir de la force de l'eau. Le présent rapport rend compte du déroulement de la visite : il présente d'abord l'accueil réservé à la délégation, ensuite les observations faites dans les différents ateliers, enfin les enseignements tirés de cette expérience.

\textit{Justification de la méthode :} le contexte (cadre, date, lieu) est posé dès la première phrase ; l'objectif est introduit par « avait pour objectif de » ; le plan est annoncé par « d'abord..., ensuite..., enfin... ». Les trois étapes de la méthode sont respectées.

\textbf{Conclusion.} Au terme de cette visite, nous retiendrons que la production de l'énergie hydroélectrique repose sur un ensemble d'opérations rigoureuses, du captage de l'eau jusqu'à la distribution du courant. Cette expérience a renforcé notre intérêt pour les métiers de l'électrotechnique. Nous suggérons à l'administration de l'établissement de multiplier de telles sorties pédagogiques, qui complètent utilement les cours théoriques. Nous remercions enfin la direction de la centrale pour son accueil chaleureux.

\textit{Justification de la méthode :} la conclusion fait le bilan (« nous retiendrons que... »), formule une suggestion (« Nous suggérons... ») et se termine par une formule de politesse. Aucune information nouvelle n'est introduite.

\textbf{Conclusion de l'exercice.} L'introduction et la conclusion rédigées remplissent toutes les fonctions attendues : elles encadrent le corps du rapport de manière cohérente et professionnelle.
\end{exoresolu}

\courssub{Savoir-faire 5 : Rédiger et améliorer le corps du rapport}

\begin{methode}[Rédiger le corps du rapport]
\begin{enumerate}
\item \textbf{Choisis un ordre} : chronologique (déroulement des faits) ou thématique (par aspects du sujet).
\item \textbf{Divise en parties numérotées} (1., 2., 3.) avec des titres courts et explicites.
\item \textbf{Rapporte les faits avec objectivité} : pas d'opinions personnelles, pas d'exagération ; utilise le passé composé et l'imparfait.
\item \textbf{Sois précis} : chiffres, noms, horaires, quantités ; emploie le lexique spécialisé adapté au domaine.
\item \textbf{Relie les parties} par des connecteurs : d'abord, ensuite, en outre, par ailleurs, enfin.
\item \textbf{Améliore} à la relecture : supprime les répétitions, corrige l'orthographe et les accords, vérifie que chaque partie répond à l'objectif annoncé dans l'introduction.
\end{enumerate}
\end{methode}

\begin{exoresolu}[Amélioration d'un passage de rapport]
\textbf{Énoncé.} Voici un extrait du corps d'un rapport de stage rédigé par un élève de Terminale G1. Il contient cinq défauts (subjectivité, imprécision, répétition, faute d'accord, temps inadapté). Réécris ce passage en le corrigeant, puis justifie chaque correction.

« Le stage c'était super génial, j'ai trop aimé. On est arrivé à l'entreprise vers une heure je crois. Le chef il nous a montré les machines, les machines était très grandes. Ensuite on a vu plein de choses intéressantes et c'était vraiment bien. »

\tcblower
\textbf{Solution détaillée.}

\textbf{Passage réécrit :}
« Le stage s'est déroulé dans de bonnes conditions et a été très profitable. Nous sommes arrivés à l'entreprise à 13 heures. Le chef d'atelier nous a présenté les machines, qui étaient de grandes dimensions. Ensuite, nous avons observé le fonctionnement de la chaîne de production et les opérations de contrôle de qualité. »

\textbf{Justification des corrections :}
\begin{enumerate}
\item \textbf{Subjectivité supprimée} : « c'était super génial, j'ai trop aimé » exprime un jugement personnel et un langage familier ; on le remplace par un constat objectif : « s'est déroulé dans de bonnes conditions et a été très profitable ».
\item \textbf{Imprécision corrigée} : « vers une heure je crois » est vague ; un rapport exige une information exacte : « à 13 heures ».
\item \textbf{Répétition supprimée} : « les machines, les machines » est remplacé par le pronom relatif « qui », ce qui allège la phrase.
\item \textbf{Faute d'accord corrigée} : « les machines était » devient « qui étaient » (accord du verbe avec le sujet pluriel).
\item \textbf{Temps inadaptés corrigés} : « on est arrivé », « on a vu » (oral, accord incorrect) deviennent « nous sommes arrivés », « nous avons observé » : emploi du « nous » propre au rapport et accord du participe passé (« arrivés » avec le sujet « nous »).
\item \textbf{Précision du lexique} : « plein de choses intéressantes » est remplacé par des faits précis : « le fonctionnement de la chaîne de production et les opérations de contrôle de qualité ».
\end{enumerate}
\textbf{Conclusion.} Le passage corrigé est objectif, précis, correctement accordé et rédigé dans un niveau de langue soutenu, conforme aux exigences du rapport.
\end{exoresolu}

\courssub{Savoir-faire 6 : Rédiger un rapport complet à partir d'une situation concrète}

\begin{methode}[De la préparation à la rédaction complète]
\begin{enumerate}
\item \textbf{Prépare} : relis la situation, souligne les données utiles (qui, quoi, où, quand, pourquoi) et dresse un plan en trois parties (introduction, corps en 2 ou 3 sous-parties, conclusion).
\item \textbf{Rédige la structure externe} : en-tête, intitulé, lieu et date, rédacteur, destinataire.
\item \textbf{Rédige l'introduction} : contexte + objectif + annonce du plan.
\item \textbf{Rédige le corps} : faits ordonnés, parties titrées, lexique adapté, temps du passé.
\item \textbf{Rédige la conclusion} : bilan + suggestions + politesse.
\item \textbf{Améliore} : relecture orthographique et grammaticale, vérification de la cohérence entre l'introduction et le corps, présentation soignée.
\end{enumerate}
\end{methode}

\begin{exoresolu}[Sujet type Bac : rédaction d'un rapport]
\textbf{Énoncé.} Tu es le délégué de la classe de Terminale F2 du lycée technique de Bafoussam. Le vendredi 6 mars 2026, ta classe a participé à une journée de reboisement organisée par la mairie sur la colline de Banengo. Le proviseur te demande un rapport sur cette activité. Rédige ce rapport en respectant la structure externe et la structure interne (20 à 25 lignes).

\tcblower
\textbf{Solution détaillée (rapport complet).}

LYCÉE TECHNIQUE DE BAFOUSSAM

\textbf{Rapport sur la journée de reboisement du 6 mars 2026}

Bafoussam, le 9 mars 2026

De : NANA Aïcha, déléguée de la Terminale F2

À : Monsieur le Proviseur

\textbf{Introduction}

Dans le cadre des activités citoyennes de l'établissement, la classe de Terminale F2 a participé, le vendredi 6 mars 2026, à une journée de reboisement organisée par la mairie de Bafoussam sur la colline de Banengo. Cette activité visait à contribuer à la protection de l'environnement et à sensibiliser les élèves à la lutte contre la déforestation. Le présent rapport décrit le déroulement de la journée, présente les résultats obtenus et formule quelques suggestions.

\textbf{1. Le déroulement de la journée}

La classe s'est réunie devant le lycée à 7 h 30 et a rejoint la colline de Banengo à 8 heures, sous la conduite de deux enseignants accompagnateurs. Après l'accueil par les agents municipaux, les élèves ont reçu des plants et du matériel (houes, arrosoirs, gants). Les travaux ont commencé à 9 heures : les élèves, répartis en cinq groupes de huit, ont creusé les trous, mis les plants en terre et procédé au premier arrosage. Les activités se sont achevées à 13 heures.

\textbf{2. Les résultats obtenus}

Au total, la classe a planté 250 arbres, essentiellement des acacias et des eucalyptus. Aucun incident n'a été signalé au cours des travaux. Les agents de la mairie ont félicité les élèves pour leur discipline et leur efficacité.

\textbf{Conclusion}

Cette journée de reboisement a été une réussite : elle a permis aux élèves de mettre en pratique les notions d'éducation civique et de prendre conscience de l'importance du couvert végétal. Nous suggérons que le lycée institue chaque année une journée similaire et qu'un suivi des plants soit assuré par les classes volontaires. Nous remercions Monsieur le Proviseur d'avoir autorisé cette sortie.

\textit{Justification de la méthode :} la structure externe est complète (en-tête, intitulé, lieu et date, rédacteur, destinataire) ; l'introduction pose le contexte, l'objectif et le plan ; le corps est divisé en deux parties numérotées et titrées, avec des faits précis (horaires, chiffres) au passé composé ; la conclusion fait le bilan, suggère et remercie.

\textbf{Conclusion de l'exercice.} Ce rapport est complet, ordonné, objectif et rédigé dans un niveau de langue adapté : il répond pleinement à la demande du destinataire.
\end{exoresolu}

\begin{attention}[Les 5 erreurs qui coûtent des points au Bac]
\begin{enumerate}
\item \textbf{Confondre rapport et lettre} : le rapport n'a ni formule d'appel (« Monsieur, » en début de texte) ni formule de politesse de lettre administrative (« Veuillez agréer... ») ; il a un intitulé, un rédacteur, un destinataire et un corps structuré. Ne mélange pas les deux genres.
\item \textbf{Oublier des éléments de la structure externe} : sans intitulé, sans date, sans nom du rédacteur ou du destinataire, le rapport est incomplet et perd des points de présentation. Vérifie chaque élément avant de rendre la copie.
\item \textbf{Introduire des informations nouvelles dans la conclusion} : la conclusion fait le bilan et suggère ; elle ne rapporte aucun fait qui n'aurait pas été traité dans le corps.
\item \textbf{Employer un langage familier ou subjectif} : « c'était super », « j'ai trop aimé », « je crois » sont à bannir. Le rapport exige l'objectivité, la précision (chiffres, horaires exacts) et un niveau de langue soutenu.
\item \textbf{Négliger la relecture} : fautes d'accord du participe passé (« nous sommes arrivé »), répétitions, temps verbaux incohérents (mélange présent/passé sans raison). Réserve toujours cinq minutes à la correction finale : c'est là que se gagnent les derniers points.
\end{enumerate}
\end{attention}

\newpage

\coursec{Exercices}

\courssub{Je vérifie mes connaissances}

\begin{exercice}[QCM : Caractéristiques du texte informatif et du rapport]
Parmi les propositions suivantes, une seule est exacte pour chaque question. Recopiez le numéro de la question et la lettre de la bonne réponse.
\begin{enumerate}
\item Le texte informatif se distingue du texte argumentatif principalement par :
\begin{itemize}
\item[A.] La présence d'une thèse défendue par l'auteur.
\item[B.] La neutralité et l'objectivité de l'énonciation.
\item[C.] L'utilisation massive de figures de style expressives.
\item[D.] L'appel à l'émotion du destinataire.
\end{itemize}
\item Dans la structure externe d'un rapport, l'élément qui identifie l'auteur, le destinataire, l'objet et la date est :
\begin{itemize}
\item[A.] Le corps du rapport.
\item[B.] La conclusion.
\item[C.] L'en-tête.
\item[D.] Les annexes.
\end{itemize}
\item Le lexique spécialisé (ou technique) se caractérise par :
\begin{itemize}
\item[A.] Son usage courant dans la langue de tous les jours.
\item[B.] Sa polysémie systématique.
\item[C.] Sa monosémie et sa précision dans un domaine professionnel donné.
\item[D.] Son caractère subjectif et affectif.
\end{itemize}
\item La conclusion d'un rapport professionnel doit obligatoirement comporter :
\begin{itemize}
\item[A.] Un résumé détaillé de tous les faits observés.
\item[B.] Des recommandations ou des propositions d'actions concrètes.
\item[C.] Les remerciements adressés au personnel rencontré.
\item[D.] La transcription intégrale des entretiens menés.
\end{itemize}
\end{enumerate}
\end{exercice}

\begin{exercice}[Vrai / Faux justifié : Structure et rédaction]
Indiquez si les affirmations suivantes sont vraies ou fausses. Justifiez chaque réponse en une phrase précise en vous appuyant sur les normes du rapport technique.
\begin{enumerate}
\item L'introduction d'un rapport se rédige toujours en premier, avant le corps du texte.
\item L'objectivité dans un rapport s'obtient notamment par l'usage du « je » de l'énonciateur pour garantir l'authenticité du témoignage.
\item Les connecteurs logiques (par conséquent, en outre, toutefois) sont facultatifs dans le corps du rapport car les faits parlent d'eux-mêmes.
\item Les annexes (plans, photos, tableaux bruts) font partie intégrante de la pagination principale du rapport.
\end{enumerate}
\end{exercice}

\begin{exercice}[Définitions et distinctions lexicales]
Définissez brièvement les notions suivantes et donnez un exemple pour chacune, en lien avec une filière technique (BTP, Électrotechnique, Secrétariat, etc.) :
\begin{enumerate}
\item Lexique commun.
\item Lexique spécialisé (ou technique).
\item Niveau de langue soutenu (dans le rapport).
\item Marque d'objectivité (donnez deux procédés grammaticaux ou lexicaux).
\end{enumerate}
\end{exercice}

\begin{exercice}[Identification de la structure sur énoncé court]
Voici les éléments désordonnés d'un rapport de stage. Remettez-les dans l'ordre logique de la structure externe standard en les numérotant de 1 à 7.
\begin{itemize}
\item[A.] Recommandations et signature.
\item[B.] En-tête (Expéditeur, Destinataire, Objet, Date, Références).
\item[C.] Table des matières.
\item[D.] Introduction (Contexte, Objet, Méthodologie, Plan annoncé).
\item[E.] Corps du rapport (Développement structuré en parties/chapitres).
\item[F.] Conclusion (Bilan, Perspectives).
\item[G.] Annexes (Documents joints).
\end{itemize}
\end{exercice}

\courssub{Je m'entraîne}

\begin{exercice}[Analyse d'un rapport d'inspection de marché (Réception)]
On donne ci-après un extrait du rapport d'inspection du marché central de Bafoussam, rédigé par un agent du Ministère du Commerce.

\textbf{Extrait :}

« \emph{Objet : Rapport d'inspection des conditions d'hygiène et de sécurité au Marché Central de Bafoussam.\\
Référence : Note de service n° 045/MINCOMMERCE/SG/DQ du 10/01/2025.\\
Monsieur le Ministre,\\
J'ai l'honneur de vous transmettre le compte rendu de la mission d'inspection effectuée du 15 au 17 janvier 2025 dans les installations du Marché Central. Cette mission faisait suite aux plaintes récurrentes des usagers relatives à l'insalubrité des lieux.\\
L'inspection a porté sur trois secteurs : la halle vivrière, le secteur boucherie/poissonnerie et les voiries d'accès. Les constatations principales sont les suivantes :}
\begin{itemize}
\item \emph{\textbf{Halle vivrière :} présence de flaques d'eau stagnante favorisant la prolifération larvaire (anophèles). Les étals en bois non traités présentent un taux d'humidité supérieur à 25 \%.}
\item \emph{\textbf{Boucherie/Poissonnerie :} rupture de la chaîne du froid constatée sur 3 chambres froides sur 5 (température relevée à +8 \textdegree C au lieu de $-$18 \textdegree C). Absence de bacs de récupération des effluents.}
\item \emph{\textbf{Voiries :} dégradation avancée du revêtement bitumineux (nids-de-poule, ornières) entravant l'évacuation des eaux pluviales et l'accès des camions de collecte des ordures ménagères.}
\end{itemize}
\emph{Au vu de ces constats, il est impératif de procéder à la réhabilitation du système d'assainissement et au remplacement des groupes frigorifiques défaillants.\\
Fait à Bafoussam, le 20 janvier 2025.\\
L'Inspecteur principal, [Signature]} »

À partir de ce texte, effectuez les tâches suivantes :
\begin{enumerate}
\item Identifiez les 5 composantes de l'en-tête administratif présentes dans l'extrait.
\item Délimitez les trois grandes parties de la structure interne (Introduction, Corps, Conclusion) en citant la première et la dernière phrase de chacune.
\item Relevez 5 marques d'objectivité (indices linguistiques : temps, modes, voix, connecteurs, données chiffrées) et expliquez leur effet.
\item Extrayez 5 termes relevant du lexique spécialisé (domaine : hygiène publique / BTP / logistique) et définissez-les en contexte.
\item Quel est le type de rapport (rapport d'activité, d'enquête, d'inspection, de stage, de visite) ? Justifiez.
\end{enumerate}
\end{exercice}

\begin{exercice}[Tri lexical et reformulation (Morphosyntaxe / Lexique)]
Voici une liste de 20 mots extraits de deux rapports distincts : un rapport d'intervention chirurgicale (médical) et un rapport de réception de travaux (BTP).

\textbf{Liste :} \emph{Anesthésie, Béton, Hémostase, Ferraillage, Scalpel, Coffrage, Sutures, Vibreur, Bistouri, Armature, Incision, Banche, Compresse, Adjuvant, Drap, Ciment, Hématome, Granulats, Péridurale, Talochage.}

\begin{enumerate}
\item Classez ces mots en deux colonnes : « Lexique médical » et « Lexique BTP ».
\item Choisissez 3 termes médicaux et 3 termes BTP. Pour chacun, proposez un équivalent en lexique commun (vulgarisation) compréhensible par un patient ou un client non spécialiste.
\item Réécrivez la phrase technique suivante en langage clair pour un usager : « \emph{Le vibreur à aiguille a permis la compacité optimale du béton frais C25/30 dans les banches du voile, évitant les nids de gravier.} »
\end{enumerate}
\end{exercice}

\begin{exercice}[Rédaction de l'introduction et de la conclusion (Production)]
\textbf{Situation :} Vous êtes stagiaire en Terminale Technique (spécialité Électrotechnique) à la société \emph{EnergySolutions Cameroon} à Douala-Bonabéri. Vous avez effectué un stage de 4 semaines (du 02/06/2025 au 28/06/2025) sur la maintenance préventive des onduleurs industriels (marque APC, modèle Symmetra PX 250 kW).

Rédigez :
\begin{enumerate}
\item L'\textbf{introduction} complète du rapport de stage (Contexte, Objet, Méthodologie, Plan annoncé) en respectant les formules d'usage administratif (Destinataire, Objet, Références fictives).
\item La \textbf{conclusion} du rapport incluant : le bilan synthétique des compétences acquises, les difficultés rencontrées (ex. : indisponibilité de pièces de rechange), et deux recommandations formulées professionnellement (ex. : suggestion d'un stock de sécurité, proposition de formation continue).
\end{enumerate}
\end{exercice}

\begin{exercice}[Amélioration d'un corps de rapport fautif (Rédaction / Amélioration)]
Le paragraphe ci-dessous est extrait du corps d'un rapport de visite d'usine (usine de transformation de cacao). Il présente des défauts majeurs : subjectivité, redondances, absence de connecteurs logiques, données vagues, vocabulaire familier.

\textbf{Texte à corriger :}

« \emph{Moi j'ai vu que l'usine c'est vraiment grand. Il y a plein de machines qui font du bruit. Les ouvriers ils travaillent dur. C'est sale par terre, y'a des coques de cacao partout. Le patron il a dit qu'ils font 50 tonnes par jour mais moi je crois que c'est moins. Les fèves elles sont bien séchées. À la fin on a mangé du chocolat c'était bon. Il faut qu'ils nettoient mieux.} »

En appliquant la \textbf{grille en trois passes} (1. Correction syntaxique et grammaticale ; 2. Objectivation et précision chiffrée ; 3. Fluidification par connecteurs et lexique spécialisé), proposez une version révisée, professionnelle et structurée de ce paragraphe, en inventant des données chiffrées plausibles (surface, effectif, tonnage, taux d'humidité).
\end{exercice}

\courssub{Je me prépare au Bac}

\begin{exercice}[Type Baccalauréat Technique : Reconstitution et justification de la structure (Probatoire / Bac)]
Les huit paragraphes suivants constituent le corps d'un rapport de visite d'étude effectué par une classe de Terminale Maintenance des Systèmes Automatisés (MSA) dans une usine d'embouteillage de boissons à Douala. Ils sont présentés dans le désordre.

\textbf{Paragraphes mélangés :}
\begin{enumerate}
\item[A.] \emph{Ensuite, nous avons visité la chaîne d'embouteillage. La cadence affichée est de 36 000 bouteilles/heure. Les automates Siemens S7-1500 pilotent les remplisseuses et les boucheuses. Nous avons noté l'absence de capteur de niveau sur la remplisseuse n° 3, source de pertes de produit estimées à 2 \% par l'opérateur.}
\item[B.] \emph{Pour terminer, la visite du laboratoire qualité a permis d'observer les tests de teneur en CO$_2$ et de pression de rupture. Les résultats sont consignés sur le registre qualité conforme à la norme ISO 9001 version 2015.}
\item[C.] \emph{En introduction, le chef de maintenance, M. Ngono, a présenté l'organigramme du service : 12 techniciens, 3 ingénieurs, une astreinte 24 h/24. La GMAO (Gestion de Maintenance Assistée par Ordinateur) utilisée est le logiciel « MainPro ».}
\item[D.] \emph{Par ailleurs, l'atelier de maintenance préventive nous a été présenté. Le planning hebdomadaire affiche les gammes de graissage et de vérification des courroies. Cependant, le stock de pièces de rechange (roulements SKF 6205, courroies SPZ) semble sous-dimensionné par rapport à la criticité des équipements.}
\item[E.] \emph{La visite a débuté par la salle de traitement d'eau. L'osmose inverse assure une conductivité inférieure à 10 µS/cm. Le débitmètre électromagnétique indique un débit instantané de 15 m$^3$/h.}
\item[F.] \emph{Enfin, la réunion de synthèse avec le Directeur d'usine a permis d'échanger sur la politique RSE (Responsabilité Sociétale de l'Entreprise) et le traitement des eaux usées (station d'épuration à boues activées).}
\item[G.] \emph{En premier lieu, le groupe froid (compresseurs à vis, fluide frigorigène ammoniac R717) assure le refroidissement de l'eau glacée à +1 \textdegree C. La surveillance des pressions haute et basse s'effectue via le système SCADA. Une fuite mineure sur le condenseur n° 2 a été signalée.}
\item[H.] \emph{De plus, le conditionnement de fin de ligne (fardeleuse, palettiseur robotisé) clôture le processus. L'analyse vibratoire du robot palettiseur révèle un pic à 25 mm/s sur l'axe Z, dépassant le seuil d'alarme de la norme ISO 10816 (18 mm/s).}
\end{enumerate}

\textbf{Travail demandé :}
\begin{enumerate}
\item Reconstituez l'ordre logique de ces 8 paragraphes (de 1 à 8) pour former le corps cohérent du rapport.
\item Pour chaque transition (ex. : paragraphe 1 $\to$ paragraphe 2), justifiez votre choix en identifiant : le connecteur logique utilisé, la progression thématique (chronologique, spatiale, analytique), et le champ lexical qui assure la cohésion.
\item Rédigez l'introduction qui précéderait ce corps (contexte : visite pédagogique du 12 mars 2025 ; objectif : confronter la théorie MSA à la réalité industrielle ; méthodologie : observation guidée, questionnaire technique ; plan annoncé en 4 parties correspondant aux 4 grands ateliers visités).
\end{enumerate}
\end{exercice}

\begin{exercice}[Type Baccalauréat Technique : Production écrite complète — Rapport de visite d'étude]
\textbf{Sujet :}

Vous êtes élève en classe de Terminale Technique, spécialité \textbf{Génie Civil / BTP} (ou votre spécialité réelle). Dans le cadre du module « Technologie Professionnelle », votre classe a effectué une visite d'étude d'une journée (le 05 avril 2025) sur le chantier de construction d'un pont sur le fleuve Wouri à Douala, piloté par une entreprise internationale de travaux publics.

Vous êtes désigné(e) par le Directeur des Études de votre établissement pour rédiger le \textbf{rapport officiel de cette visite d'étude}.

\textbf{Consignes strictes :}
\begin{enumerate}
\item Respectez scrupuleusement la \textbf{structure externe} (En-tête complet, Introduction, Corps structuré en parties titrées, Conclusion, Annexes mentionnées).
\item L'\textbf{introduction} doit contenir : Contexte (cadre pédagogique), Objet de la visite, Méthodologie (observation, entretien avec l'ingénieur résident, prise de vues), Plan annoncé (3 parties).
\item Le \textbf{corps} (3 parties minimum) doit développer :
\begin{itemize}
\item \textbf{Partie 1 : Organisation du chantier et moyens mis en œuvre} (effectifs, engins : grues, barges, palplanches ; planning).
\item \textbf{Partie 2 : Progrès techniques observés sur les ouvrages d'art} (fondations sur pieux forés de diamètre 2,5 m, tablier en poutres précontraintes, mise en œuvre du béton haute performance C50/60, contrôles qualité : essai d'affaissement au cône, éprouvettes).
\item \textbf{Partie 3 : Gestion HSE (Hygiène, Sécurité, Environnement) et contraintes} (port des EPI, gestion des déchets, nuisances sonores, coactivité fluviale et routière).
\end{itemize}
\item Intégrez au moins \textbf{5 données chiffrées précises} (inventées mais réalistes : nombre d'ouvriers, volume de béton, profondeur des pieux, délai, coûts) et \textbf{10 termes du lexique spécialisé BTP} (identifiables).
\item La \textbf{conclusion} doit présenter un bilan pédagogique (apports pour la formation), une difficulté majeure identifiée (ex. : logistique fluviale), et \textbf{deux recommandations techniques ou pédagogiques} formulées avec le registre de la recommandation professionnelle (il conviendrait de..., il est recommandé de...).
\item Soignez la langue : objectivité, impersonnalité (passif, infinitif, « il »), connecteurs logiques, temps appropriés (passé composé pour les faits constatés, présent pour les descriptions techniques, futur ou conditionnel pour les recommandations).
\end{enumerate}
\end{exercice}

\begin{exercice}[Type Probatoire : Exploitation de documents et rédaction d'un rapport d'incident]
\textbf{Situation :} Vous êtes Technicien de Maintenance au sein de la société nationale des eaux (Direction Régionale du Centre). Le 15 mai 2025, une rupture de canalisation principale (DN 400, fonte ductile) s'est produite au carrefour de l'Amitié à Yaoundé, privant d'eau les quartiers Mvog-Ada, Essos et Nlongkak pendant 18 heures.

Vous disposez des notes brutes suivantes (prises sur le terrain) :
\begin{itemize}
\item Heure d'alerte : 06 h 15 (appel usager). Arrivée de l'équipe : 07 h 30.
\item Localisation : PK 12+500, axe Yaoundé--Douala, emprise routière.
\item Constat : fissure longitudinale sur 1,2 m, fuite sous pression (estimation : 15 L/s). Sol argileux gonflant.
\item Action : fermeture des vannes de sectionnement PK 12+000 et PK 13+000. Mise en charge d'un by-pass DN 150 (débit 5 L/s) pour l'hôpital central et la caserne.
\item Réparation : découpe du tronçon défectueux (scie à disque), pose d'un manchon de réparation, remise en charge progressive à 20 h 45.
\item Moyens : 1 pelleteuse, 1 camion-grue, 8 agents (2 soudeurs, 4 manœuvres, 1 chef d'équipe, 1 agent qualité).
\item Cause probable (avis du chef d'équipe) : affaissement différentiel du remblai (tassement de 15 cm en 5 ans) entraînant une contrainte de cisaillement sur la canalisation ; absence d'un lit de pose sableux conforme aux règles de l'art.
\item Conséquences : 12 000 abonnés impactés, 3 réclamations écrites, perte d'eau estimée à 972 m$^3$, coût de réparation estimé à 2,4 millions de FCFA.
\end{itemize}

\textbf{Travail demandé :}

Rédigez le \textbf{Rapport d'Incident Technique} adressé au Directeur Régional, en respectant la structure standard (En-tête, Introduction, Déroulement chronologique des opérations, Analyse des causes, Bilan chiffré et conséquences, Mesures correctives et préventives, Conclusion).
\begin{enumerate}
\item Transformez les notes brutes en un texte continu, structuré, objectif et professionnel.
\item Utilisez le lexique spécialisé hydraulique / travaux publics (DN, fonte ductile, vanne de sectionnement, by-pass, manchon de réparation, tassement différentiel, lit de pose, remise en charge).
\item Intégrez un tableau de synthèse « Moyens humains et matériels » dans le corps ou en annexe (décrivez-le textuellement).
\item La conclusion doit proposer 2 mesures préventives techniques (ex. : réhabilitation du tronçon, surveillance du réseau) et 1 mesure organisationnelle (ex. : astreinte renforcée).
\end{enumerate}
\end{exercice}

\begin{methode}[De la note brute au rapport : la démarche en 5 étapes]
\begin{enumerate}
\item \textbf{Trier} les informations : faits chronologiques, constats techniques, chiffres, causes, conséquences.
\item \textbf{Structurer} selon le plan imposé : en-tête, introduction, déroulement, analyse, bilan, mesures, conclusion.
\item \textbf{Objectiver} : supprimer le « je », employer la voix passive et les tournures impersonnelles, remplacer les appréciations par des mesures.
\item \textbf{Chiffrer et vérifier la cohérence} : contrôler les calculs (ex. : 15 L/s $\times$ 18 h $\approx$ 972 m$^3$), harmoniser les unités.
\item \textbf{Relire} : orthographe, connecteurs logiques, concordance des temps, registre soutenu.
\end{enumerate}
\end{methode}

\begin{attention}[Pièges fréquents dans la rédaction du rapport]
\begin{itemize}
\item Confondre \textbf{compte rendu} (relate fidèlement des faits ou des échanges) et \textbf{rapport} (analyse, interprète et recommande).
\item Rédiger l'introduction avant le corps : elle doit annoncer le contenu réel, donc se rédiger en dernier.
\item Employer « je » ou des jugements de valeur (« c'est sale », « c'est bien ») au lieu de constats mesurables.
\item Oublier les connecteurs logiques, ce qui transforme le rapport en simple liste de faits.
\item Mélanger lexique familier et lexique spécialisé : le registre doit rester soutenu et technique du début à la fin.
\item Négliger l'en-tête : sans expéditeur, destinataire, objet, date et références, le document n'a pas de valeur administrative.
\end{itemize}
\end{attention}

\newpage

\coursec{Corrigés des exercices}

\coursec{Corrigés des exercices — Leçon 05 : Production d'un rapport}

\begin{corrige}[Exercice 1 — QCM : Caractéristiques du texte informatif et du rapport]
\begin{enumerate}
\item \textbf{B.} Le texte informatif transmet des faits de manière neutre et objective ; la défense d'une thèse (A) est le propre du texte argumentatif, les figures expressives (C) relèvent du texte littéraire, l'appel à l'émotion (D) du texte publicitaire ou polémique.
\item \textbf{C.} L'en-tête est la carte d'identité du rapport : expéditeur, destinataire, objet, date, références. Le corps (A) développe les faits, la conclusion (B) propose des suites, les annexes (D) rassemblent les pièces jointes.
\item \textbf{C.} Le lexique spécialisé est monosémique et précis dans un domaine donné ; l'usage courant (A) et la polysémie (B) caractérisent au contraire le lexique commun.
\item \textbf{B.} La conclusion d'un rapport professionnel débouche sur des recommandations ou des propositions d'actions concrètes ; c'est ce qui distingue le rapport du simple compte rendu.
\end{enumerate}
\textbf{Points de vigilance :} ne pas confondre « objectivité » et « absence totale de style » : le texte informatif reste soigné, mais il efface l'énonciateur et les émotions.
\end{corrige}

\begin{corrige}[Exercice 2 — Vrai / Faux justifié : Structure et rédaction]
\begin{enumerate}
\item \textbf{Faux.} L'introduction est placée en tête du document, mais elle se rédige \emph{en dernier}, une fois le corps achevé, afin d'annoncer exactement le contenu réel du rapport.
\item \textbf{Faux.} L'objectivité s'obtient par l'effacement de l'énonciateur : tournures impersonnelles (« il a été constaté que... »), voix passive, infinitifs de consigne, « nous » collectif ou « il ». Le « je » subjectivise le propos.
\item \textbf{Faux.} Les connecteurs logiques sont indispensables : ils articulent les faits (chronologie, cause, conséquence, opposition) et garantissent la cohérence ; sans eux, le rapport devient une simple liste.
\item \textbf{Faux.} Les annexes sont placées \emph{après} la conclusion et disposent d'un repérage propre (Annexe 1, Annexe 2...) ; elles ne font pas partie du corps paginé du rapport.
\end{enumerate}
\textbf{Points de vigilance :} une justification au Bac doit citer la norme (voix passive, place des annexes) et non se contenter d'un « c'est la règle ».
\end{corrige}

\begin{corrige}[Exercice 3 — Définitions et distinctions lexicales]
\begin{enumerate}
\item \textbf{Lexique commun :} ensemble des mots de la langue courante, compris de tous les locuteurs, souvent polysémiques. Exemples : \emph{fil}, \emph{table}, \emph{courant} (le mot « courant » a plusieurs sens selon le contexte).
\item \textbf{Lexique spécialisé :} ensemble des termes précis et monosémiques propres à un domaine professionnel. Exemples : \emph{onduleur}, \emph{disjoncteur différentiel} (électrotechnique) ; \emph{palplanche}, \emph{éprouvette} (BTP) ; \emph{chronogramme} (secrétariat).
\item \textbf{Niveau de langue soutenu (registre administratif) :} registre soigné, conforme aux normes de l'écrit professionnel : vocabulaire précis, syntaxe complète, absence de familiarités. Exemple : « il conviendrait de procéder à... » et non « il faut qu'ils fassent... ».
\item \textbf{Marques d'objectivité (deux procédés) :} (a) la voix passive : « les températures ont été relevées » ; (b) la tournure impersonnelle : « il apparaît que... ». On peut ajouter les données chiffrées et les infinitifs.
\end{enumerate}
\textbf{Points de vigilance :} l'exemple doit être cohérent avec la filière choisie ; un terme spécialisé n'a qu'un seul sens dans son domaine. Distinguer le \emph{lexique} (vocabulaire) du \emph{registre} (niveau de langue).
\end{corrige}

\begin{corrige}[Exercice 4 — Identification de la structure sur énoncé court]
Ordre logique : 1. \textbf{B} (En-tête) ; 2. \textbf{C} (Table des matières) ; 3. \textbf{D} (Introduction) ; 4. \textbf{E} (Corps du rapport) ; 5. \textbf{F} (Conclusion) ; 6. \textbf{A} (Recommandations et signature) ; 7. \textbf{G} (Annexes).

\textbf{Justification :} l'en-tête identifie le document ; la table des matières l'annonce ; l'introduction présente contexte, objet, méthodologie et plan ; le corps développe ; la conclusion fait le bilan ; les recommandations et la signature engagent l'auteur ; les annexes, non essentielles à la lecture continue, ferment le dossier.

\textbf{Remarque :} dans les rapports courts, les recommandations sont intégrées à la conclusion ; dans les rapports longs, elles forment une section distincte avant la signature.
\end{corrige}

\begin{corrige}[Exercice 5 — Analyse d'un rapport d'inspection de marché (Réception)]
\begin{enumerate}
\item \textbf{Les 5 composantes canoniques de l'en-tête :}
\begin{itemize}
\item \textbf{Expéditeur (Émetteur) :} identifié par la signature et la qualité (« L'Inspecteur principal »).
\item \textbf{Destinataire :} « Monsieur le Ministre ».
\item \textbf{Objet :} « Rapport d'inspection des conditions d'hygiène... ».
\item \textbf{Date et lieu de rédaction :} « Fait à Bafoussam, le 20 janvier 2025 ».
\item \textbf{Références :} note de service n° 045/MINCOMMERCE/SG/DQ du 10/01/2025 (document déclencheur).
\end{itemize}
\item \textbf{Structure interne :}
\begin{itemize}
\item \emph{Introduction :} de « J'ai l'honneur de vous transmettre le compte rendu de la mission... » à « ...relatives à l'insalubrité des lieux » (contexte et objet).
\item \emph{Corps :} de « L'inspection a porté sur trois secteurs... » à « ...l'accès des camions de collecte des ordures ménagères » (constatations par secteur).
\item \emph{Conclusion :} « Au vu de ces constats, il est impératif de procéder à la réhabilitation du système d'assainissement et au remplacement des groupes frigorifiques défaillants » (recommandations).
\end{itemize}
\item \textbf{5 marques d'objectivité :}
\begin{itemize}
\item Données chiffrées : « 3 chambres froides sur 5 », « +8 \textdegree C au lieu de $-$18 \textdegree C », « supérieur à 25 \% » : le propos devient mesurable et vérifiable.
\item Voix passive : « température relevée », « rupture constatée » : effacement de l'agent.
\item Tournure impersonnelle : « il est impératif de... » : la prescription s'impose d'elle-même.
\item Connecteur de conséquence : « Au vu de ces constats » : la recommandation découle logiquement des faits.
\item Vocabulaire précis, sans émotion : aucun jugement de valeur, seulement des constats.
\end{itemize}
\item \textbf{5 termes spécialisés :} \emph{chaîne du froid} (maintien ininterrompu des denrées à basse température), \emph{effluents} (rejets liquides d'une activité), \emph{prolifération larvaire} (développement des larves de moustiques), \emph{revêtement bitumineux} (couche de bitume d'une chaussée), \emph{groupes frigorifiques} (appareils producteurs de froid).
\item \textbf{Type de rapport :} rapport \textbf{d'inspection} : l'auteur, mandaté par une autorité (note de service), contrôle la conformité d'installations à des normes, dresse des constats et formule des prescriptions.
\end{enumerate}
\textbf{Points de vigilance :} pour délimiter les parties, citer les premières et dernières phrases \emph{textuellement} ; ne pas confondre la référence (document déclencheur) et la date de rédaction.
\end{corrige}

\begin{corrige}[Exercice 6 — Tri lexical et reformulation (Morphosyntaxe / Lexique)]
\begin{enumerate}
\item \textbf{Classement :}
\begin{center}
\begin{tabularx}{\linewidth}{|X|X|}
\hline
\rowcolor{popblueL} \textbf{Lexique médical} & \textbf{Lexique BTP} \\
\hline
anesthésie, hémostase, scalpel, sutures, bistouri, incision, compresse, drap, hématome, péridurale & béton, ferraillage, coffrage, vibreur, armature, banche, adjuvant, ciment, granulats, talochage \\
\hline
\end{tabularx}
\end{center}
\item \textbf{Vulgarisation (exemples) :}
\begin{itemize}
\item \emph{hémostase} = arrêt du saignement ;
\item \emph{incision} = coupure pratiquée par le chirurgien ;
\item \emph{hématome} = bleu, poche de sang sous la peau ;
\item \emph{coffrage} = moule dans lequel on coule le béton ;
\item \emph{granulats} = sable et graviers entrant dans la composition du béton ;
\item \emph{talochage} = lissage de la surface du béton ou de l'enduit.
\end{itemize}
\item \textbf{Reformulation en langage clair :} « Une machine vibrante a été plongée dans le béton encore liquide, coulé dans les moules du mur, afin de chasser les bulles d'air : le mur est ainsi compact et solide, sans trous ni zones creuses. »
\end{enumerate}
\textbf{Points de vigilance :} vulgariser ne signifie pas appauvrir : la reformulation doit conserver l'information technique (compacité, absence de défauts) dans des mots courants.
\end{corrige}

\begin{corrige}[Exercice 7 — Rédaction de l'introduction et de la conclusion (Production)]
\textbf{1. Introduction (modèle de production attendue) :}

« \emph{À Monsieur le Directeur des Études.\\
Objet : Rapport de stage de fin de formation — maintenance préventive des onduleurs industriels.\\
Références : convention de stage n° 12/LT/ESC/2025.\\
Monsieur le Directeur,\\
Dans le cadre de la formation en Électrotechnique et conformément au programme de Terminale technique, un stage pratique de quatre semaines a été effectué du 02/06/2025 au 28/06/2025 au sein de la société EnergySolutions Cameroon, sise à Douala-Bonabéri. Ce stage avait pour objet la maintenance préventive des onduleurs industriels de marque APC, modèle Symmetra PX 250 kW, équipements assurant la continuité d'alimentation électrique des installations sensibles. La démarche adoptée a combiné l'observation des interventions des techniciens confirmés, la participation aux opérations de contrôle et l'exploitation des documentations techniques du constructeur. Le présent rapport rend compte de cette immersion en trois parties : la présentation de la structure d'accueil et de ses équipements, les opérations de maintenance préventive réalisées, enfin le bilan des compétences acquises et les recommandations.} »

\textbf{2. Conclusion (modèle) :}

« \emph{Au terme de ce stage, les objectifs fixés peuvent être considérés comme atteints : le diagnostic des onduleurs, le remplacement des batteries et la lecture des schémas unifilaires ont été maîtrisés. La principale difficulté rencontrée réside dans l'indisponibilité de certaines pièces de rechange, qui a allongé les délais d'intervention. Au vu de ces constats, il est recommandé, d'une part, de constituer un stock de sécurité des pièces critiques (batteries, cartes de puissance) et, d'autre part, d'organiser une formation continue des techniciens aux nouveaux modèles du constructeur. Il conviendrait enfin de pérenniser le partenariat entre l'entreprise et l'établissement scolaire.} »

\textbf{Points de vigilance :} l'introduction doit contenir les quatre éléments obligatoires (contexte, objet, méthodologie, plan annoncé) ; la conclusion doit formuler les recommandations au registre professionnel (« il est recommandé de... », « il conviendrait de... »), jamais à l'impératif direct.
\end{corrige}

\begin{corrige}[Exercice 8 — Amélioration d'un corps de rapport fautif (Rédaction / Amélioration)]
\textbf{Version révisée (modèle) :}

« L'unité de transformation visitée s'étend sur une superficie d'environ 2 hectares et emploie 85 ouvriers répartis en deux équipes. L'atelier de production comprend 12 machines principales (torréfacteurs, broyeurs, concasseurs), dont le niveau sonore, élevé, justifierait le port systématique de protections auditives. Selon le chef de production, la capacité journalière annoncée est de 50 tonnes de fèves ; toutefois, les registres consultés font état d'une production moyenne réelle de 42 tonnes par jour. En revanche, la qualité du séchage est satisfaisante : le taux d'humidité des fèves, mesuré à 7 \%, est conforme à la norme commerciale (8 \% maximum). Cependant, la propreté des aires de stockage laisse à désirer : des coques de cacao jonchent le sol, ce qui favorise les risques de contamination et d'accident. Par conséquent, il est recommandé de mettre en place un plan de nettoyage quotidien des aires de stockage et un contrôle hebdomadaire de son application. »

\textbf{Commentaire méthodologique (les trois passes) :}
\begin{enumerate}
\item \emph{Correction syntaxique :} suppression des dislocations familières (« Les ouvriers ils travaillent »), des abréviations orales (« y'a »), des phrases nominales.
\item \emph{Objectivation :} disparition du « je » et des jugements (« c'est sale », « c'était bon ») au profit de données chiffrées (2 ha, 85 ouvriers, 42 t/j, 7 \%) ; suppression du passage anecdotique (dégustation), sans valeur informative.
\item \emph{Fluidification :} connecteurs (toutefois, en revanche, cependant, par conséquent) et lexique spécialisé (torréfacteur, taux d'humidité, aire de stockage).
\end{enumerate}
\textbf{Points de vigilance :} les données inventées doivent rester plausibles et cohérentes entre elles ; la recommandation finale doit découler des constats.
\end{corrige}

\begin{corrige}[Exercice 9 — Type Bac : Reconstitution et justification de la structure]
\begin{enumerate}
\item \textbf{Ordre logique :} 1. C ; 2. E ; 3. G ; 4. A ; 5. H ; 6. D ; 7. B ; 8. F.
\item \textbf{Justification des transitions :}
\begin{itemize}
\item C $\to$ E : C pose le cadre humain et organisationnel (« En introduction ») ; E enchaîne sur le début effectif du parcours (« La visite a débuté par... »). Progression chronologique ; champ lexical de l'accueil et de l'organisation.
\item E $\to$ G : « En premier lieu » poursuit l'énumération des ateliers ; cohésion par le champ lexical des utilités (eau, froid), qui précèdent la production.
\item G $\to$ A : « Ensuite » marque l'étape suivante : passage des utilités au cœur de production (embouteillage) ; champ lexical des équipements automatisés (automates, remplisseuses).
\item A $\to$ H : « De plus » ajoute l'étape qui suit logiquement l'embouteillage : le conditionnement de fin de ligne ; progression spatiale le long de la chaîne.
\item H $\to$ D : « Par ailleurs » ouvre un point de vue transversal : la maintenance des équipements décrits (roulements, courroies, analyse vibratoire renvoyant au robot de H).
\item D $\to$ B : « Pour terminer » annonce la dernière étape de la visite : le laboratoire qualité.
\item B $\to$ F : « Enfin » clôt le parcours par la réunion de synthèse, moment conclusif d'échange avec la direction.
\end{itemize}
\item \textbf{Introduction (modèle) :} « Dans le cadre du module de Technologie Professionnelle de la filière Maintenance des Systèmes Automatisés, la classe de Terminale MSA a effectué, le 12 mars 2025, une visite d'étude dans une usine d'embouteillage de boissons de la zone industrielle de Douala. L'objectif de cette sortie pédagogique était de confronter les enseignements théoriques (automatisme, maintenance préventive, qualité) aux réalités d'un site industriel en exploitation. La démarche adoptée a combiné l'observation guidée des installations, l'administration d'un questionnaire technique auprès des responsables de service et la prise de notes structurée. Le présent rapport rend compte de cette visite en quatre parties : l'organisation du service maintenance, les utilités (traitement d'eau et groupe froid), la chaîne de production et de conditionnement, enfin la maintenance préventive et le contrôle qualité. »
\end{enumerate}
\textbf{Barème indicatif (sur 20) :} ordre correct des 8 paragraphes : 4 pts ; justification des transitions (connecteur + progression + champ lexical) : 8 pts ; introduction complète (contexte, objet, méthodologie, plan) : 6 pts ; langue et présentation : 2 pts.

\textbf{Points de vigilance :} une transition non justifiée (ordre donné sans explication) ne vaut que la moitié des points ; citer explicitement le connecteur du texte.
\end{corrige}

\begin{corrige}[Exercice 10 — Type Bac : Production écrite complète — Rapport de visite d'étude]
\textbf{Modèle de production attendue (éléments rédactionnels) :}

\emph{En-tête :} « République du Cameroun — Paix, Travail, Patrie. Lycée Technique de [ville]. À Monsieur le Directeur des Études. Objet : Rapport de visite d'étude — chantier du pont sur le fleuve Wouri (Douala), 05 avril 2025. Références : note de service n° 08/LT/DE/2025. Rédigé par : [Nom], élève de Terminale Génie Civil. »

\emph{Introduction :} « Dans le cadre du module Technologie Professionnelle, la classe de Terminale Génie Civil a effectué, le 05 avril 2025, une visite d'étude sur le chantier de construction du pont sur le fleuve Wouri à Douala, piloté par une entreprise internationale de travaux publics. Cette visite avait pour objet d'observer les techniques de réalisation des ouvrages d'art et de confronter les enseignements théoriques aux pratiques de chantier. La méthodologie a combiné l'observation directe, un entretien avec l'ingénieur résident et la prise de vues photographiques. Le présent rapport s'articule en trois parties : l'organisation du chantier et les moyens mis en œuvre, les progrès techniques observés sur l'ouvrage, enfin la gestion HSE et les contraintes du site. »

\emph{Corps — Partie 1 (Organisation et moyens) :} le chantier emploie 180 ouvriers et 25 cadres, répartis en deux équipes de 10 heures. Le parc matériel comprend 3 grues à tour de 12 tonnes, 2 barges de travail, 4 pelleteuses et une centrale à béton de 60 m$^3$/h. Le planning prévisionnel fixe la livraison à 30 mois, avec un coût estimé à 45 milliards de FCFA.

\emph{Corps — Partie 2 (Progrès techniques) :} les fondations reposent sur 48 pieux forés de diamètre 2,5 m, descendus à 38 m de profondeur. Le tablier est constitué de poutres précontraintes coulées en béton haute performance C50/60. Le contrôle qualité s'appuie sur l'essai d'affaissement au cône (cible : 12 cm) et sur la confection d'éprouvettes testées à 7 et 28 jours ; le volume total de béton est estimé à 32 000 m$^3$.

\emph{Corps — Partie 3 (HSE et contraintes) :} le port des EPI (casque, gilet, chaussures de sécurité, harnais pour les travaux en hauteur) est obligatoire et contrôlé quotidiennement. Les déchets de chantier sont triés et évacués vers une filière agréée. Les principales contraintes identifiées sont la coactivité fluviale (navigation sur le Wouri), les nuisances sonores pour les quartiers riverains et la logistique d'acheminement des matériels lourds.

\emph{Conclusion :} « Cette visite a permis de concrétiser les notions de fondations profondes, de précontrainte et de contrôle qualité abordées en classe. La difficulté majeure observée réside dans la logistique fluviale, qui conditionne l'approvisionnement du chantier. Il est recommandé, d'une part, d'organiser une visite de suivi au stade du montage du tablier et, d'autre part, d'inviter l'ingénieur résident à animer une conférence technique dans l'établissement. Il conviendrait enfin de systématiser la rédaction d'un rapport après chaque sortie pédagogique. »

\emph{Annexes mentionnées :} Annexe 1 : plan de masse du chantier ; Annexe 2 : photographies ; Annexe 3 : grille d'entretien avec l'ingénieur résident.

\textbf{Barème indicatif (sur 20) :} structure externe complète : 3 pts ; introduction (4 éléments) : 3 pts ; corps en 3 parties titrées, données chiffrées (5 min.) et lexique spécialisé (10 termes) : 8 pts ; conclusion (bilan, difficulté, 2 recommandations) : 3 pts ; objectivité, connecteurs, langue : 3 pts.

\textbf{Points de vigilance :} vérifier la cohérence des chiffres inventés ; employer le passé composé pour les faits constatés, le présent pour les descriptions techniques, le conditionnel pour les recommandations ; ne jamais écrire « je » : préférer « la classe a effectué », « il a été observé ».
\end{corrige}

\begin{corrige}[Exercice 11 — Type Probatoire : Rédaction d'un rapport d'incident]
\textbf{Vérification préalable des calculs :} perte d'eau : $15 \text{ L/s} \times 18 \text{ h} = 15 \times 64\,800 \text{ s} = 972\,000 \text{ L} = 972 \text{ m}^3$ : la donnée des notes est exacte et peut être reprise telle quelle. Durée totale de l'intervention : de 06 h 15 (alerte) à 20 h 45 (remise en charge), soit environ 14 h 30 ; la coupure d'alimentation annoncée (18 h) correspond à la fenêtre de perturbation constatée par les usagers : on conservera les deux durées sans les confondre.

\textbf{Modèle de rapport attendu :}

\emph{En-tête :} « Société Nationale des Eaux — Direction Régionale du Centre. À Monsieur le Directeur Régional. Objet : Rapport d'incident technique — rupture de canalisation DN 400 au carrefour de l'Amitié, Yaoundé, 15 mai 2025. Références : ordre d'intervention n° 112/DRC/2025. Rédigé par : [Nom], Technicien de Maintenance. »

\emph{Introduction :} « Le 15 mai 2025, une rupture de la canalisation principale DN 400 en fonte ductile a été signalée au carrefour de l'Amitié à Yaoundé, entraînant une coupure d'alimentation dans les quartiers Mvog-Ada, Essos et Nlongkak. Le présent rapport décrit le déroulement des opérations, analyse les causes de l'incident, en établit le bilan chiffré et propose des mesures correctives et préventives. »

\emph{Déroulement chronologique :} « L'alerte a été reçue à 06 h 15 (appel usager) ; l'équipe d'intervention est arrivée sur site à 07 h 30. Le point de fuite a été localisé au PK 12+500 de l'axe Yaoundé--Douala, en emprise routière : une fissure longitudinale de 1,2 m laissait échapper un débit estimé à 15 L/s, dans un sol argileux gonflant. Les vannes de sectionnement PK 12+000 et PK 13+000 ont été fermées, puis un by-pass DN 150 (débit 5 L/s) a été mis en charge pour desservir l'hôpital central et la caserne. Le tronçon défectueux a été découpé à la scie à disque, un manchon de réparation a été posé, et la remise en charge progressive est intervenue à 20 h 45. »

\emph{Analyse des causes :} « Selon le chef d'équipe, la cause probable est un affaissement différentiel du remblai (tassement de 15 cm en 5 ans), ayant engendré une contrainte de cisaillement sur la canalisation, aggravée par l'absence d'un lit de pose sableux conforme aux règles de l'art. »

\emph{Bilan chiffré :} « 12 000 abonnés ont été impactés ; 3 réclamations écrites ont été enregistrées ; la perte d'eau est estimée à 972 m$^3$ ; le coût de la réparation est évalué à 2,4 millions de FCFA. Les moyens mobilisés sont synthétisés dans le tableau ci-dessous (Annexe 1) : 1 pelleteuse, 1 camion-grue, 8 agents (2 soudeurs, 4 manœuvres, 1 chef d'équipe, 1 agent qualité). »

\emph{Mesures et conclusion :} « L'incident a été maîtrisé en 14 h 30 d'intervention, mais il révèle la vulnérabilité du tronçon. Il est recommandé, d'une part, la réhabilitation complète du tronçon PK 12+000 -- PK 13+000 avec reprise du lit de pose et, d'autre part, la mise en place d'une surveillance périodique du réseau par campagne de détection acoustique des fuites. Sur le plan organisationnel, il conviendrait de renforcer l'astreinte du week-end afin de réduire le délai d'arrivée sur site. »

\textbf{Barème indicatif (sur 20) :} en-tête complet : 2 pts ; introduction : 2 pts ; déroulement chronologique fidèle aux notes : 4 pts ; analyse des causes : 2 pts ; bilan chiffré exact (dont 972 m$^3$ justifié) : 3 pts ; tableau des moyens décrit : 2 pts ; mesures (2 techniques + 1 organisationnelle) : 3 pts ; objectivité, lexique spécialisé, langue : 2 pts.

\textbf{Points de vigilance :} transformer les notes en texte continu (pas de puces dans le corps) ; employer la voix passive (« ont été fermées », « a été posé ») ; distinguer clairement le fait constaté (fissure de 1,2 m) de l'hypothèse (cause « probable ») ; toute mesure préventive doit répondre à une cause identifiée.
\end{corrige}

\newpage

\coursec{Fiche bilan}

\begin{popfigure}
\begin{tikzpicture}[
  mindmap,
  concept color=popblue,
  level 1 concept/.append style={level distance=3.5cm, sibling angle=60},
  level 2 concept/.append style={level distance=2.5cm},
  every node/.style={font=\small},
  text=white
]
\node[concept, align=center]{Le Rapport\\Technique}
  child[concept color=poporange] { node[concept, align=center]{1. Découverte\\Texte informatif} }
  child[concept color=popgreen] { node[concept, align=center]{2. Structure\\Externe \& Interne} }
  child[concept color=poppurple] { node[concept, align=center]{3. Lexique\\Commun / Spécialisé} }
  child[concept color=poppink] { node[concept, align=center]{4. Introduction\\ \& Conclusion} }
  child[concept color=popgold] { node[concept, align=center]{5. Corps du\\Rapport} }
  child[concept color=popblue] { node[concept, align=center]{6. Préparation\\ \& Rédaction} };
\end{tikzpicture}
\legende{Carte mentale : Les 6 étapes clés de la leçon « Produire un rapport » (Terminale Technique).}
\end{popfigure}

\begin{bilan}

\courssub{Définitions clés}

\begin{definition}[Rapport technique]
Document écrit, structuré et objectif, qui rend compte d'une activité, d'une observation, d'une enquête ou d'un stage. Il informe une hiérarchie ou un commanditaire pour faciliter la prise de décision.
\end{definition}

\begin{definition}[Texte informatif]
Texte à visée \cle{référentielle} (fonction informative) : transmet des faits, des données, des résultats sans subjectivité. Caractères : clarté, exactitude, concision, objectivité, structure logique.
\end{definition}

\begin{definition}[Structure externe]
Éléments matériels identifiables : Page de garde, Sommaire/Table des matières, Liste des abréviations/sigles, Introduction, Corps (chapitres/sections), Conclusion, Annexes, Bibliographie.
\end{definition}

\begin{definition}[Structure interne]
Organisation logique du contenu : Plan (thématique, chronologique, analytique), articulation des parties (transitions), paragraphes (une idée = un paragraphe), connecteurs logiques.
\end{definition}

\courssub{Lexique : Commun vs Spécialisé}

\begin{center}
\begin{tabularx}{\linewidth}{|l|X|X|}
\hline
\rowcolor{popblueL}
\textbf{Critère} & \textbf{Lexique commun} & \textbf{Lexique spécialisé (technique)} \\
\hline
Définition & Mots du langage courant, compris de tous. & Mots propres à un métier, une filière, une discipline. \\
\hline
Exemple & « La machine fonctionne bien. » & « Le compresseur centrifuge atteint son régime nominal. » \\
\hline
Rôle dans le rapport & Accessibilité, fluidité, introduction/conclusion. & Précision, rigueur, crédibilité technique (corps du rapport). \\
\hline
Attention & Éviter le familier, l'imprécis. & Définir à la première occurrence (glossaire ou note), ne pas abuser (javellisation). \\
\hline
\end{tabularx}
\end{center}

\courssub{Méthodes express}

\begin{methode}[Rédiger l'Introduction (Modèle en entonnoir)]
\begin{enumerate}
\item \textbf{Accroche / Contexte :} Situation générale, commande, stage, problème technique.
\item \textbf{Sujet / Objet :} Définition précise du thème traité.
\item \textbf{Intérêt / Finalité :} Pourquoi ce rapport ? (Décision, amélioration, validation de compétences).
\item \textbf{Plan annoncé :} « Nous présenterons d'abord... puis... enfin... »
\end{enumerate}
\end{methode}

\begin{methode}[Rédiger la Conclusion (Modèle en entonnoir inversé)]
\begin{enumerate}
\item \textbf{Bilan / Synthèse :} Rappel des résultats majeurs (réponse à l'objet).
\item \textbf{Interprétation / Portée :} Signification technique, limites, fiabilité.
\item \textbf{Perspectives / Recommandations :} Actions correctives, améliorations, suites à donner.
\item \textbf{Formule de politesse finale (si adressé à une hiérarchie).}
\end{enumerate}
\end{methode}

\begin{methode}[Rédiger le Corps du rapport]
\begin{enumerate}
\item Appliquer le plan annoncé (structure interne cohérente).
\item Un paragraphe = une idée principale + arguments/données/chiffres + illustration (schéma, tableau, photo référencée).
\item Utiliser des \cle{connecteurs logiques} (chronologiques, causaux, adversatifs, concessifs).
\item Intégrer le lexique spécialisé (défini) et des données chiffrées (\SI{...}{...}).
\item Rester \cle{objectif} : pas de « je pense que », préférer « les mesures indiquent que », « l'analyse révèle que ».
\end{enumerate}
\end{methode}

\begin{methode}[Processus complet : Préparation $\to$ Rédaction $\to$ Amélioration]
\begin{enumerate}
\item \textbf{Préparation :} Analyser la consigne/commanditaire $\to$ Collecter/selectionner infos (notes de terrain, fiches) $\to$ Élaborer plan détaillé (arborescence).
\item \textbf{Rédaction (brouillon) :} Rédiger intro, corps, conclusion sans chercher la perfection stylistique immédiate. Respecter la structure externe.
\item \textbf{Amélioration (révision) :} 
\begin{itemize}
\item \textit{Macro :} Cohérence plan/annonce, logique, exhaustivité.
\item \textit{Micro :} Syntaxe, vocabulaire (commun/spécialisé), connecteurs, orthographe, présentation (numérotation, titres, légendes figures/tableaux).
\end{itemize}
\item \textbf{Finalisation :} Mise en page soignée, vérification sommaire/annexes, relecture à voix haute.
\end{enumerate}
\end{methode}

\courssub{Points de vigilance (Pièges fréquents)}

\begin{attention}[Erreurs à éviter au Bac]
\begin{enumerate}
\item Confondre \textbf{Compte-rendu} (chronologique, exhaustif, neutre) et \textbf{Rapport} (analytique, sélectif, orienté décision/recommandation).
\item Introduction « vide » (pas de plan annoncé, pas d'objet précis).
\item Corps du rapport sans structure interne visible (bloc de texte unique, pas de titres de sections).
\item Subjectivité : « J'ai trouvé ça difficile », « C'était intéressant ». $\to$ Remplacer par constats mesurés.
\item Lexique spécialisé non défini ou mal utilisé (anglicismes non adaptés, sigles non explicités).
\item Conclusion qui introduit de nouvelles informations ou qui se limite à « J'ai appris beaucoup de choses ».
\item Oubli des éléments de structure externe (page de garde non normée, pas de sommaire, annexes non référencées dans le corps).
\end{enumerate}
\end{attention}

\end{bilan}

\begin{aretenir}[Checklist avant le Bac]
\begin{itemize}
\item $\square$ Je distingue texte informatif / texte argumentatif / texte prescriptif.
\item $\square$ Je sais lister les 8 à 10 éléments de la \textbf{structure externe} dans l'ordre.
\item $\square$ Je sais construire un \textbf{plan détaillé} (thématique, chronologique ou analytique) avec titres de parties/sous-parties.
\item $\square$ Je maîtrise la structure en \textbf{3 temps de l'introduction} (Contexte $\to$ Objet $\to$ Plan).
\item $\square$ Je maîtrise la structure en \textbf{3 temps de la conclusion} (Bilan $\to$ Portée $\to$ Recommandations).
\item $\square$ J'utilise des \textbf{connecteurs logiques} variés et adaptés à l'enchaînement des idées.
\item $\square$ Je fais la différence \textbf{lexique commun / lexique spécialisé} et je définis ce dernier.
\item $\square$ J'intègre des \textbf{données chiffrées} (mesures, pourcentages) avec unités SI correctes.
\item $\square$ Je sais \textbf{réviser} mon texte (cohérence macro + correction micro + présentation).
\item $\square$ Je respecte les \textbf{normes de présentation} (marges, police, interligne, numérotation pages, légendes figures/tableaux).
\item $\square$ Je sais adapter le ton (objectif, impersonnel, passif/pronominal) au destinataire (hiérarchie, client, jury).
\item $\square$ Je gère mon temps : Préparation (20\%) / Rédaction (60\%) / Amélioration (20\%).
\end{itemize}
\end{aretenir}

\findecours

\end{document}
